% Generated mechanically from the frozen more-etale.tex target.
% Do not edit this generated file; edit the bound locale source instead.

% OK, start here.
%

\setcounter{chapter}{62}
\chapter{进一步的 \'Etale 上同调}



\phantomsection
\label{more-etale-section-phantom}



\section{引言}
\label{more-etale-section-introduction}

\noindent
本章是概形 \'etale 上同调系列的第二章。第一章见“\'Etale 上同调”
第 \ref{etale-cohomology-section-introduction} 节。

\medskip\noindent
粗略地说，本章与前一章的分界如下：凡涉及“叹号函子”（紧支撑上同调及其
右伴随）以及使用这些材料的内容，都归入本章。






\section{延拓截面}
\label{more-etale-section-growing}

\noindent
本节讨论如下类型的结果。

\begin{lemma}
\label{more-etale-lemma-section-support-in-locally-closed-pre}
设 $X$ 为概形，$\mathcal{F}$ 为 $X_\etale$ 上的阿贝尔层。
设 $\varphi : U' \to U$ 为 $X_\etale$ 中的态射。设 $Z' \subset U'$ 为
闭子概形，使得 $Z' \to U' \to U$ 是闭浸入，其像为 $Z \subset U$。
则存在典范双射
$$
\{s \in \mathcal{F}(U) \mid \text{Supp}(s) \subset Z\} =
\{s' \in \mathcal{F}(U') \mid \text{Supp}(s') \subset Z'\}
$$
若 $\varphi^{-1}(Z) = Z'$，则该双射由限制给出。
\end{lemma}

\begin{proof}
考虑 $U'$ 的闭子概形 $Z'' = \varphi^{-1}(Z)$。由于 $Z'$ 在 $U'$ 中闭，
故 $Z' \subset Z''$ 闭。另一方面，$Z' \to Z''$ 是 \'etale 态射
（这是 $Z$ 上两个 \'etale 概形之间的态射），因而是开映射。因此，
对某个闭子集 $T$ 有 $Z'' = Z' \amalg T$。开覆盖
$U' = (U' \setminus T) \cup (U' \setminus Z')$ 表明
$$
\{s' \in \mathcal{F}(U') \mid \text{Supp}(s') \subset Z'\} =
\{s' \in \mathcal{F}(U' \setminus T) \mid \text{Supp}(s') \subset Z'\}
$$
而 \'etale 覆盖
$\{U' \setminus T \to U, U \setminus Z \to U\}$ 表明
$$
\{s \in \mathcal{F}(U) \mid \text{Supp}(s) \subset Z\} =
\{s' \in \mathcal{F}(U' \setminus T) \mid \text{Supp}(s') \subset Z'\}
$$
证明完成。
\end{proof}

\begin{lemma}
\label{more-etale-lemma-section-support-in-locally-closed}
设 $X$ 为概形，$Z \subset X$ 为局部闭子概形，$\mathcal{F}$ 为
$X_\etale$ 上的阿贝尔层。给定开子概形 $U, U' \subset X$，其中
$Z$ 均为闭子概形，则存在典范双射
$$
\{s \in \mathcal{F}(U) \mid \text{Supp}(s) \subset Z\} =
\{s \in \mathcal{F}(U') \mid \text{Supp}(s) \subset Z\}
$$
若 $U' \subset U$，则该双射由限制给出。
\end{lemma}

\begin{proof}
由于 $Z$ 是 $U \cap U'$ 的闭子概形，只须证明 $U' \subset U$ 的情形；
此时结论是引理 \ref{more-etale-lemma-section-support-in-locally-closed-pre} 的特例。
\end{proof}

\noindent
下面引入一个稍非标准、但后文很有用的记号。具体而言，在上述引理
\ref{more-etale-lemma-section-support-in-locally-closed} 的情形下，记
$$
H_Z(\mathcal{F}) = \{s \in \mathcal{F}(U) \mid \text{Supp}(s) \subset Z\}
$$
其中 $U \subset X$ 可任取为一个以 $Z$ 为闭子概形的开子概形。
若读者对这一写法的不精确性感到不安，可以取
$U = X \setminus \partial Z$，其中
$\partial Z = \overline{Z}\setminus Z$ 是 $Z$ 在 $X$ 中的“边界”。
不过，在下文许多论证中，能够灵活选择不同开子概形会发挥作用。
这一构造具有如下性质：
\begin{enumerate}
\item
\label{more-etale-item-inclusion}
若 $Z \subset Z'$ 是 $X$ 的局部闭子概形，且 $Z$ 在 $Z'$ 中闭，
则存在自然单射
$$
H_Z(\mathcal{F}) \to H_{Z'}(\mathcal{F}).
$$
\item
\label{more-etale-item-pullback}
若 $f : Y \to X$ 是概形态射，且 $Z \subset X$ 是局部闭子概形，
则存在自然拉回映射
$f^* : H_Z(\mathcal{F}) \to H_{f^{-1}Z}(f^{-1}\mathcal{F})$。
\end{enumerate}
把这一记号推广到如下情形会很方便：设 $W \in X_\etale$，且
$Z \subset W$ 为局部闭子概形。此时记
$$
H_Z(\mathcal{F}) =
\{s \in \mathcal{F}(U) \mid \text{Supp}(s) \subset Z\} =
H_Z(\mathcal{F}|_{W_\etale})
$$
其中 $U \subset W$ 与上面一样，可任取为一个以 $Z$ 为闭子概形的
开子概形\footnote{事实上，引理
\ref{more-etale-lemma-section-support-in-locally-closed-pre}
表明：若 $X$ 上的 $Z$ 同构于 $X_\etale$ 中某个对象 $W$ 的局部闭子概形，
则 $W$ 的选择无关紧要。}。







\section{紧支撑截面}
\label{more-etale-section-compact-support}

\noindent
本节可参考 \cite[报告 XVII，第 6 节]{SGA4}。
设 $f : X \to Y$ 为分离且局部有限型的概形态射。本节定义函子
$f_! : \textit{Ab}(X_\etale) \to \textit{Ab}(Y_\etale)$
如下：令 $f_!\mathcal{F} \subset f_*\mathcal{F}$ 为由相对于 $Y$ 具有
固有支撑（按适当方式定义）的截面所成的子层。

\medskip\noindent
警告：函子 $f_!$ 是导出范畴上某个函子 $Rf_!$ 的第零上同调层
（此处待补未来引用），但 $Rf_!$ 不是 $f_!$ 的导出函子。

\begin{lemma}
\label{more-etale-lemma-f-shriek-separated}
设 $f : X \to Y$ 为局部有限型概形态射，$\mathcal{F}$ 为
$X_\etale$ 上的阿贝尔层。规则
$$
Y_\etale \longrightarrow \textit{Ab},\quad
V \longmapsto \{s \in f_*\mathcal{F}(V) = \mathcal{F}(X_V) \mid
\text{Supp}(s) \subset X_V \text{ 固有地映到 }V\}
$$
定义了 $f_*\mathcal{F}$ 的阿贝尔子层。
\end{lemma}

\noindent
警告：若 $f$ 不分离，这个层并非“正确的”层。

\begin{proof}
回忆截面的支撑是闭的（“\'Etale 上同调”引理
\ref{etale-cohomology-lemma-support-section-closed}），故“概形上同调”
第 \ref{coherent-section-proper-over-base} 节的材料适用。由上面的引理及
“概形上同调”引理 \ref{coherent-lemma-union-closed-proper-over-base}，
可知我们在 $f_*\mathcal{F}(V)$ 中取出的子集是子群。由“概形上同调”
引理 \ref{coherent-lemma-base-change-closed-proper-over-base}，
可知该规则定义了子预层。
最后，设 $s \in f_*\mathcal{F}(V)$，并设有 \'etale 覆盖
$\{V_i \to V\}$，使得 $s|_{V_i}$ 的支撑在 $V_i$ 上固有。
注意，$s|_{V_i}$ 的支撑是 $s|_V$ 的支撑的逆像（使用以茎刻画支撑的方式
以及“\'Etale 上同调”引理
\ref{etale-cohomology-lemma-stalk-pullback}）。于是由“下降”引理
\ref{descent-lemma-descending-property-proper-over-base}，$s$ 的支撑在
$V$ 上固有。这证明该规则满足层条件。
\end{proof}

\begin{lemma}
\label{more-etale-lemma-separated-etale-shriek}
设 $j : U \to X$ 为分离 \'etale 态射，$\mathcal{F}$ 为
$U_\etale$ 上的阿贝尔层。“\'Etale 上同调”引理
\ref{etale-cohomology-lemma-shriek-into-star-separated-etale}
中的单射 $j_!\mathcal{F} \to j_*\mathcal{F}$ 之像，正是引理
\ref{more-etale-lemma-f-shriek-separated} 的子层。
\end{lemma}

\noindent
也可以把这个引理移到后面，再用两个层的茎的描述来证明。

\begin{proof}
在“\'Etale 上同调”引理
\ref{etale-cohomology-lemma-shriek-into-star-separated-etale} 的证明中，
$j_!\mathcal{F} \to j_*\mathcal{F}$ 是通过构造预层间的映射
$j_{p!}\mathcal{F} \to j_*\mathcal{F}$ 得到的；后者的像显然包含在引理
\ref{more-etale-lemma-f-shriek-separated} 的子层中。由于 $j_!\mathcal{F}$ 是
$j_{p!}\mathcal{F}$ 的层化，故 $j_!\mathcal{F} \to j_*\mathcal{F}$ 的像
也包含在该子层中。
反之，设 $s \in j_*\mathcal{F}(V)$ 的支撑 $Z$ 在 $V$ 上固有。则
$Z \to V$ 有限，且具有闭像 $Z' \subset V$；见“更多态射”引理
\ref{more-morphisms-lemma-characterize-finite}。$s$ 在
$V \setminus Z'$ 上的限制为零，而零截面属于
$j_!\mathcal{F} \to j_*\mathcal{F}$ 的像。
另一方面，若 $v \in Z'$，则可找到 \'etale 邻域
$(V', v') \to (V, v)$，使得有分解
$U_{V'} = W \amalg U'_1 \amalg \ldots \amalg U'_n$，
其中各项都是既开又闭的子概形，$U'_i \to V'$ 为同构，且
$T_{V'} \subset U'_1 \amalg \ldots \amalg U'_n$；见“\'Etale 态射”
引理 \ref{etale-lemma-etale-etale-local-technical}。把同构
$U'_i \to V'$ 取逆，并拉回 $s$，便得到 $n$ 个态射
$\varphi'_i : V' \to U$ 以及 $V'$ 上的截面 $s'_i$。
于是 $V'$ 上 $j_{p!}\mathcal{F}$ 的截面
$\sum (\varphi'_i, s'_i)$（参见“\'Etale 上同调”引理
\ref{etale-cohomology-lemma-shriek-into-star-separated-etale} 的证明中
$j_{p!}\mathcal{F}(V')$ 的公式）按构造映到 $s$ 在 $V'$ 上的限制。
因此，$s$ 在 \'etale 局部属于 $j_!\mathcal{F} \to j_*\mathcal{F}$ 的像，
证明完成。
\end{proof}

\begin{definition}
\label{more-etale-definition-f-shriek-separated}
设 $f : X \to Y$ 为分离（！）且局部有限型的概形态射，$\mathcal{F}$ 为
$X_\etale$ 上的阿贝尔层。引理 \ref{more-etale-lemma-f-shriek-separated} 所构造的
子层 $f_!\mathcal{F} \subset f_*\mathcal{F}$ 称为{\it 紧支撑直像}。
\end{definition}

\noindent
由引理 \ref{more-etale-lemma-separated-etale-shriek}，这不与“\'Etale 上同调”定义
\ref{etale-cohomology-definition-extension-zero} 冲突，因为两种定义均适用时
彼此一致。下面作一个基本检验。

\begin{lemma}
\label{more-etale-lemma-proper-f-shriek}
设 $f : X \to Y$ 为概形的固有态射。则 $f_! = f_*$。
\end{lemma}

\begin{proof}
由 $f_!$ 的构造立即得到。
\end{proof}

\noindent
下面的观察非常有用。

\begin{remark}[关于开嵌入的协变性]
\label{more-etale-remark-covariance-f-shriek-separated}
设 $f : X \to Y$ 为分离且局部有限型的概形态射，$\mathcal{F}$ 为
$X_\etale$ 上的阿贝尔层。设 $X' \subset X$ 为开子概形，并以
$f' : X' \to Y$ 表示 $f$ 的限制。存在典范单射
$$
f'_!(\mathcal{F}|_{X'}) \longrightarrow f_!\mathcal{F}
$$
具体地，设 $V \in Y_\etale$，并考虑截面
$s' \in f'_*(\mathcal{F}|_{X'})(V) = \mathcal{F}(X' \times_Y V)$，
其支撑 $Z'$ 在 $V$ 上固有。则 $Z'$ 在 $X \times_Y V$ 中也闭；见
“概形上同调”引理
\ref{coherent-lemma-functoriality-closed-proper-over-base}。因此存在唯一截面
$s \in \mathcal{F}(X \times_Y V) = f_*\mathcal{F}(V)$，使得它在
$X' \times_Y V$ 上的限制为 $s'$，而在
$X \times_Y V \setminus Z'$ 上的限制为零；见引理
\ref{more-etale-lemma-section-support-in-locally-closed}。这一构造与限制映射相容，
因而诱导所需的层映射 $f'_!(\mathcal{F}|_{X'}) \to f_!\mathcal{F}$，
它显然是单射。按构造得到交换图
$$
\xymatrix{
f'_!(\mathcal{F}|_{X'}) \ar[r] \ar[d] &
f_!\mathcal{F} \ar[d] \\
f'_*(\mathcal{F}|_{X'}) &
f_*\mathcal{F} \ar[l]
}
$$
它关于 $\mathcal{F}$ 具有函子性。显然，若 $X'' \subset X'$ 为开子概形，
且 $f'' = f|_{X''} : X'' \to Y$，则刚构造的典范映射之复合
$f''_!\mathcal{F}|_{X''} \to f'_!\mathcal{F}|_{X'} \to f_!\mathcal{F}$
就是典范映射 $f''_!\mathcal{F}|_{X''} \to f_!\mathcal{F}$。
\end{remark}

\begin{lemma}
\label{more-etale-lemma-compactify-f-shriek-separated}
设 $Y$ 为概形。设 $j : X \to \overline{X}$ 为 $Y$ 上概形的开浸入，
且 $\overline{X}$ 在 $Y$ 上固有。以 $f : X \to Y$ 和
$\overline{f} : \overline{X} \to Y$ 表示结构态射。对
$\mathcal{F} \in \textit{Ab}(X_\etale)$，存在典范同构（见证明）
$$
f_!\mathcal{F} \longrightarrow \overline{f}_!j_!\mathcal{F}
$$
由引理 \ref{more-etale-lemma-proper-f-shriek} 有
$\overline{f}_! = \overline{f}_*$，故作为
$\textit{Ab}(X_\etale) \to \textit{Ab}(Y_\etale)$ 的函子，有
$\overline{f}_* \circ j_! = f_!$。
\end{lemma}

\begin{proof}
由“\'Etale 上同调”引理
\ref{etale-cohomology-lemma-jshriek-open}，有
$(j_!\mathcal{F})|_X = \mathcal{F}$。因此，所示箭头就是评注
\ref{more-etale-remark-covariance-f-shriek-separated} 中对
$\mathcal{G} = j_!\mathcal{F}$ 所得的单射
$f_!(\mathcal{G}|_X) \to \overline{f}_!\mathcal{G}$
。由该映射的显式构造，现在只须证明：若 $V \in Y_\etale$，且
$s \in \overline{f}_!\mathcal{G}(V) = \overline{f}_*\mathcal{G}(V) =
\mathcal{G}(\overline{X}_V)$ 为截面，则 $s$ 的支撑包含在开子概形
$X_V \subset \overline{X}_V$ 中。这立即来自如下事实：
在 $\overline{X} \setminus X$ 的几何点处，$\mathcal{G}$ 的茎为零。
\end{proof}

\noindent
我们希望把 $f_!\mathcal{F}$ 的茎与纤维上的紧支撑截面联系起来。
为叙述这一点，需要先作一个定义。

\begin{definition}
\label{more-etale-definition-compact-support}
设 $X$ 为域 $k$ 上分离且局部有限型的概形，$\mathcal{F}$ 为
$X_\etale$ 上的阿贝尔层。令
$H^0_c(X, \mathcal{F}) \subset H^0(X, \mathcal{F})$
为支撑在 $k$ 上固有的截面之集。$H^0_c(X, \mathcal{F})$ 的元素称为
{\it 紧支撑截面}。
\end{definition}

\noindent
警告：若 $X$ 在 $k$ 上不分离，这个定义并非“正确的”定义。

\begin{lemma}
\label{more-etale-lemma-proper-compact-support}
设 $X$ 为域 $k$ 上的固有概形。则
$H^0_c(X, \mathcal{F}) = H^0(X, \mathcal{F})$。
\end{lemma}

\begin{proof}
由 $H^0_c$ 的构造立即得到。
\end{proof}

\begin{remark}[开嵌入与紧支撑截面]
\label{more-etale-remark-covariance-compact-support}
设 $X$ 为域 $k$ 上分离且局部有限型的概形，$\mathcal{F}$ 为
$X_\etale$ 上的阿贝尔层。完全类似于评注
\ref{more-etale-remark-covariance-f-shriek-separated}，对开子概形
$X' \subset X$ 存在单射
$$
H^0_c(X', \mathcal{F}|_{X'}) \longrightarrow H^0_c(X, \mathcal{F})
$$
这些映射使 $H^0_c$ 成为 $X$ 的 Zariski 位点上的一个“余层”。
\end{remark}

\begin{lemma}
\label{more-etale-lemma-compactify-compact-support}
设 $k$ 为域。设 $j : X \to \overline{X}$ 为 $k$ 上概形的开浸入，且
$\overline{X}$ 在 $k$ 上固有。对 $\mathcal{F} \in \textit{Ab}(X_\etale)$，
存在典范同构（见证明）
$$
H^0_c(X, \mathcal{F}) \longrightarrow
H^0_c(\overline{X}, j_!\mathcal{F}) =
H^0(\overline{X}, j_!\mathcal{F})
$$
其中右端等号由引理 \ref{more-etale-lemma-proper-compact-support} 得到。
\end{lemma}

\begin{proof}
由“\'Etale 上同调”引理
\ref{etale-cohomology-lemma-jshriek-open}，有
$(j_!\mathcal{F})|_X = \mathcal{F}$。因此，所示箭头就是评注
\ref{more-etale-remark-covariance-compact-support} 中对
$\mathcal{G} = j_!\mathcal{F}$ 所得的单射
$H^0_c(X, \mathcal{G}|_X) \to H^0_c(\overline{X}, \mathcal{G})$
。由该映射的显式构造，现在只须证明：若
$s \in H^0(\overline{X}, \mathcal{G})$ 为截面，则 $s$ 的支撑包含在开子概形
$X$ 中。这立即来自如下事实：在 $\overline{X} \setminus X$ 的几何点处，
$\mathcal{G}$ 的茎为零。
\end{proof}

\begin{lemma}
\label{more-etale-lemma-stalk-f-shriek-separated}
设 $f : X \to Y$ 为分离且局部有限型的概形态射，$\mathcal{F}$ 为
$X_\etale$ 上的阿贝尔层。则对任意几何点
$\overline{y} : \Spec(k) \to Y$，存在典范同构
$$
(f_!\mathcal{F})_{\overline{y}}
\longrightarrow
H^0_c(X_{\overline{y}}, \mathcal{F}|_{X_{\overline{y}}})
$$
\end{lemma}

\begin{proof}
回忆
$(f_*\mathcal{F})_{\overline{y}} = \colim f_*\mathcal{F}(V)$，其中余极限取遍
$\overline{y}$ 的 \'etale 邻域 $(V, \overline{v})$。若
$s \in f_*\mathcal{F}(V) = \mathcal{F}(X_V)$，则可将 $s$ 拉回为
$(X_V)_{\overline{v}} = X_{\overline{y}}$ 上的 $\mathcal{F}$ 的截面。
于是得到典范映射
$$
c_{\overline{y}} :
(f_*\mathcal{F})_{\overline{y}}
\longrightarrow 
H^0(X_{\overline{y}}, \mathcal{F}|_{X_{\overline{y}}})
$$
我们断言，此映射在子群 $(f_!\mathcal{F})_{\overline{y}}$ 与
$H^0_c(X_{\overline{y}}, \mathcal{F}|_{X_{\overline{y}}})$ 之间诱导双射。
该断言蕴含本引理，而且稍为精确一些：它说明引理中的等同是将
$\mathcal{F}$ 的截面拉回到 $f$ 的几何纤维所得。

\medskip\noindent
注意，任意元素
$s \in (f_!\mathcal{F})_{\overline{y}} \subset (f_*\mathcal{F})_{\overline{y}}$
都由 $c_{\overline{y}}$ 映到
$H^0_c(X_{\overline{y}}, \mathcal{F}|_{X_{\overline{y}}}) \subset
H^0(X_{\overline{y}}, \mathcal{F}|_{X_{\overline{y}}})$ 中的元素。
这是因为取截面的支撑与拉回可交换，且固有性在基变换下保持。
这至少构造出了引理所述的映射。为证明它是同构，可在 $Y$ 上作
Zariski 局部检验，因此以下设 $Y$ 为仿射概形。

\medskip\noindent
下面要用到如下观察：给定开子概形 $X' \subset X$，并令
$f' = f|_{X'}$，则得到交换图
$$
\xymatrix{
(f'_!(\mathcal{F}|_{X'}))_{\overline{y}} \ar[r] \ar[d] &
H^0_c(X'_{\overline{y}}, \mathcal{F}|_{X'_{\overline{y}}}) \ar[d] \\
(f_!\mathcal{F})_{\overline{y}} \ar[r] &
H^0_c(X_{\overline{y}}, \mathcal{F}|_{X_{\overline{y}}})
}
$$
其中横箭头是上述构造的映射，竖箭头由评注
\ref{more-etale-remark-covariance-f-shriek-separated} 和
\ref{more-etale-remark-covariance-compact-support} 给出。其理由如下：给定
$\overline{y}$ 的 \'etale 邻域 $(V, \overline{v})$ 以及截面
$s \in f_*\mathcal{F}(V) = \mathcal{F}(X_V)$，若其支撑 $Z$ 恰好包含在
$X'_V$ 中且在 $V$ 上固有，则 $s$ 同时给出
$(f'_!(\mathcal{F}|_{X'}))_{\overline{y}}$ 与
$(f_!\mathcal{F})_{\overline{y}}$ 中由图中竖箭头对应的两个元素；这两个元素
经横箭头都映到 $s$ 在 $X_{\overline{y}}$ 上的拉回
$s|_{X_{\overline{y}}}$。后者的支撑 $Z_{\overline{y}}$ 包含在
$X'_{\overline{y}}$ 中，故这一限制给出
$H^0_c(X'_{\overline{y}}, \mathcal{F}|_{X'_{\overline{y}}})$
与 $H^0_c(X_{\overline{y}}, \mathcal{F}|_{X_{\overline{y}}})$
中相容的一对元素。

\medskip\noindent
设 $s \in (f_!\mathcal{F})_{\overline{y}}$ 在
$H^0_c(X_{\overline{y}}, \mathcal{F}|_{X_{\overline{y}}})$ 中映为零。
设 $s$ 由 $s \in f_*\mathcal{F}(V) = \mathcal{F}(X_V)$ 表示，其支撑
$Z$ 在 $V$ 上固有。可设 $V$ 为仿射概形，因而 $Z$ 拟紧。于是可取包含
$Z$ 之像的拟紧开子概形 $X' \subset X$。此时 $Z$ 包含在 $X'_V$ 中，
从而 $s$ 是某个元素 $s' \in f'_!(\mathcal{F}|_{X'})(V)$ 的像，其中
$f' = f|_{X'}$ 如上一段所述。于是 $s'$ 在
$H^0_c(X'_{\overline{y}}, \mathcal{F}|_{X'_{\overline{y}}})$ 中映为零。
因此，为证明单射性，可用 $X'$ 替换 $X$，即可以假设 $X$ 拟紧。
这个情形将在下面证明。

\medskip\noindent
设 $t \in H^0_c(X_{\overline{y}}, \mathcal{F}|_{X_{\overline{y}}})$。
则 $t$ 的支撑包含在某个拟紧开子概形
$W \subset X_{\overline{y}}$ 中。因此，可找到拟紧开子概形
$X' \subset X$，使 $X'_{\overline{y}}$ 包含 $W$。显然，$t$ 属于单射
$H^0_c(X'_{\overline{y}}, \mathcal{F}|_{X'_{\overline{y}}}) \to
H^0_c(X_{\overline{y}}, \mathcal{F}|_{X_{\overline{y}}})$ 的像。
因此，为证明满射性，可用 $X'$ 替换 $X$，即可以假设 $X$ 拟紧。
这个情形将在下面证明。

\medskip\noindent
在证明的最后一段，我们证明 $X$ 拟紧且 $Y$ 仿射的情形。由
“关于平坦性的进一步结果”定理 \ref{flat-theorem-nagata}，存在 $Y$ 上的紧化
$j : X \to \overline{X}$。令 $\mathcal{G} = j_!\mathcal{F}$；由
“\'Etale 上同调”引理 \ref{etale-cohomology-lemma-jshriek-open}，有
$\mathcal{F} = \mathcal{G}|_X$。由上述讨论得到交换图
$$
\xymatrix{
(f_!\mathcal{F})_{\overline{y}} \ar[r] \ar[d] &
H^0_c(X_{\overline{y}}, \mathcal{F}|_{X_{\overline{y}}}) \ar[d] \\
(\overline{f}_!\mathcal{G})_{\overline{y}} \ar[r] &
H^0_c(\overline{X}_{\overline{y}}, \mathcal{G}|_{\overline{X}_{\overline{y}}})
}
$$
由引理 \ref{more-etale-lemma-compactify-f-shriek-separated} 和
\ref{more-etale-lemma-compactify-compact-support}，竖箭头均为同构。这把问题化归到固有态射
$\overline{X} \to Y$ 的情形。对固有态射，由引理
\ref{more-etale-lemma-proper-f-shriek}、\ref{more-etale-lemma-proper-compact-support} 以及直像的
固有基变换，我们的映射是同构；见“\'Etale 上同调”引理
\ref{etale-cohomology-lemma-proper-pushforward-stalk}。
\end{proof}

\begin{lemma}
\label{more-etale-lemma-base-change-f-shriek-separated}
考虑概形的笛卡尔方块
$$
\xymatrix{
X' \ar[r]_{g'} \ar[d]_{f'} & X \ar[d]^f \\
Y' \ar[r]^g & Y
}
$$
，其中 $f$ 分离且局部有限型。则对 $X_\etale$ 上任意阿贝尔层
$\mathcal{F}$，有
$f'_!(g')^{-1}\mathcal{F} = g^{-1}f_!\mathcal{F}$.
\end{lemma}

\begin{proof}
在非常一般的情形下，存在拉回映射
$g^{-1}f_*\mathcal{F} \to f'_*(g')^{-1}\mathcal{F}$；见“位点”第
\ref{sites-section-pullback} 节。我们断言，此映射把
$g^{-1}f_!\mathcal{F}$ 映入子层 $f'_!(g')^{-1}\mathcal{F}$，并诱导引理中的同构。

\medskip\noindent
取几何点 $\overline{y}': \Spec(k) \to Y'$，并以
$\overline{y} = g \circ \overline{y}'$ 表示其在 $Y$ 中的像。存在交换图
$$
\xymatrix{
(f_*\mathcal{F})_{\overline{y}} \ar[r] \ar[d] &
H^0(X_{\overline{y}}, \mathcal{F}|_{X_{\overline{y}}}) \ar[d] \\
(f'_*(g')^{-1}\mathcal{F})_{\overline{y}'} \ar[r] &
H^0(X'_{\overline{y}'}, (g')^{-1}\mathcal{F}|_{X'_{\overline{y}'}})
}
$$
其中横映射用于引理 \ref{more-etale-lemma-stalk-f-shriek-separated} 的证明，
竖映射是上述拉回映射。此图交换，因为所涉及的四个映射都是沿概形态射拉回
局部截面而得，且相应的概形态射图交换。由于引理所述方块为笛卡尔方块，
有 $X'_{\overline{y}'} = X_{\overline{y}}$。因此，由引理
\ref{more-etale-lemma-stalk-f-shriek-separated} 及其证明得到交换图
$$
\xymatrix{
(f_*\mathcal{F})_{\overline{y}} \ar[rrr] \ar[ddd] & & &
H^0(X_{\overline{y}}, \mathcal{F}|_{X_{\overline{y}}}) \ar[ddd] \\
& (f_!\mathcal{F})_{\overline{y}} \ar[r] \ar@{..>}[d] \ar[lu] &
H^0_c(X_{\overline{y}}, \mathcal{F}|_{X_{\overline{y}}}) \ar[d] \ar[ru] \\
& (f'_!(g')^{-1}\mathcal{F})_{\overline{y}'} \ar[r] \ar[ld] &
H^0_c(X'_{\overline{y}'}, (g')^{-1}\mathcal{F}|_{X'_{\overline{y}'}}) \ar[rd]\\
(f'_*(g')^{-1}\mathcal{F})_{\overline{y}'} \ar[rrr] & & &
H^0(X'_{\overline{y}'}, (g')^{-1}\mathcal{F}|_{X'_{\overline{y}'}})
}
$$
其中内方块的横箭头均为同构，右侧两条竖箭头均为恒等。另有，
东南、西南、东北、西北方向的箭头均为单射。因此，存在唯一的双射虚线箭头
使图交换。我们得出
$g^{-1}f_!\mathcal{F} \subset g^{-1}f_*\mathcal{F} \to f'_*(g')^{-1}\mathcal{F}$
的像包含在子层
$f'_!(g')^{-1}\mathcal{F} \subset f'_*(g')^{-1}\mathcal{F}$
中，因为这在茎上成立；见“\'Etale 上同调”定理
\ref{etale-cohomology-theorem-exactness-stalks}。同一定理继而说明所诱导的映射
为同构，证明完毕。
\end{proof}

\begin{lemma}
\label{more-etale-lemma-f-shriek-composition}
设 $f : X \to Y$ 与 $g : Y \to Z$ 为可复合的概形态射，且二者均分离并
局部有限型。设 $\mathcal{F}$ 为 $X_\etale$ 上的阿贝尔层。则作为
$(g \circ f)_*\mathcal{F}$ 的子层，有
$g_!f_!\mathcal{F} = (g \circ f)_!\mathcal{F}$。
\end{lemma}

\begin{proof}
我们强烈建议读者自行证明此结论。设 $W \in Z_\etale$，且
$s \in (g \circ f)_*\mathcal{F}(W) = \mathcal{F}(X_W)$。以
$T \subset X_W$ 表示 $s$ 的支撑；这是闭子集。注意，$s$ 是
$(g \circ f)_!\mathcal{F}$ 的截面，当且仅当 $T$ 在 $W$ 上固有。我们有
$f_!\mathcal{F} \subset f_*\mathcal{F}$，因而
$g_!f_!\mathcal{F} \subset g_!f_*\mathcal{F} \subset g_*f_*\mathcal{F}$.
另一方面，$s$ 是 $g_!f_!\mathcal{F}$ 的截面，当且仅当：(a) $T$ 在
$Y_W$ 上固有；(b) 将 $s$ 看作 $f_!\mathcal{F}$ 的截面时，其支撑
$T'$ 在 $W$ 上固有。若 (a) 成立，则 $T$ 在 $Y_W$ 中的像为闭集；又因
$f_!\mathcal{F} \subset f_*\mathcal{F}$，可知 $T' \subset Y_W$ 正是
$T$ 的像（略去细节；可考察茎）。

\medskip\noindent
归结起来，我们须证明：闭子集 $T \subset X_W$ 在 $W$ 上固有，当且仅当
$T$ 在 $Y_W$ 上固有且 $T$ 在 $Y_W$ 中的像在 $W$ 上固有。赋予 $T$
诱导的既约闭子概形结构。若 $T$ 在 $W$ 上固有，则由“态射”引理
\ref{morphisms-lemma-image-proper-scheme-closed}，$T \to Y_W$ 固有；
并由“概形上同调”引理
\ref{coherent-lemma-functoriality-closed-proper-over-base}，$T$ 在 $Y_W$
中的像在 $W$ 上固有。反之，若 $T$ 在 $Y_W$ 上固有，且 $T$ 在 $Y_W$
中的像在 $W$ 上固有，则态射 $T \to W$ 是固有态射的复合，因而固有
（这里给 $T$ 在 $Y_W$ 中的闭像赋予诱导的既约概形结构，从而把问题化为
概形态射的问题）；见“态射”引理 \ref{morphisms-lemma-composition-proper}。
\end{proof}

\begin{remark}
\label{more-etale-remark-f-shriek-base-change-composition}
上述构造的函子间同构满足以下两个性质：
\begin{enumerate}
\item 设 $f : X \to Y$、$g : Y \to Z$ 和 $h : Z \to T$ 为可复合的概形态射，
且均分离并局部有限型。则图
$$
\xymatrix{
(h \circ g \circ f)_! \ar[r] \ar[d] &
(h \circ g)_! \circ f_! \ar[d] \\
h_! \circ (g \circ f)_! \ar[r] &
h_! \circ g_! \circ f_!
}
$$
交换，其中箭头来自引理 \ref{more-etale-lemma-f-shriek-composition}。
\item 假设有概形图
$$
\xymatrix{
X' \ar[d]_{f'} \ar[r]_c & X \ar[d]^f \\
Y' \ar[d]_{g'} \ar[r]_b & Y \ar[d]^g \\
Z' \ar[r]^a & Z
}
$$
，其中两个方块均为笛卡尔方块，且 $f$ 与 $g$ 分离并局部有限型。则图
$$
\xymatrix{
a^{-1} \circ (g \circ f)_! \ar[d] \ar[rr] & &
(g' \circ f')_! \circ c^{-1} \ar[d] \\
a^{-1} \circ g_! \circ f_! \ar[r] &
g'_! \circ b^{-1} \circ f_! \ar[r] &
g'_! \circ f'_! \circ c^{-1}
}
$$
交换，其中横箭头来自引理 \ref{more-etale-lemma-base-change-f-shriek-separated}，
箭头来自引理 \ref{more-etale-lemma-f-shriek-composition}。
\end{enumerate}
(1) 成立是因为直像有类似的交换图。(2) 成立是因为直像的基变换映射具有
非常一般的相容性（“位点”评注 \ref{sites-remark-compose-base-change}），
且引理 \ref{more-etale-lemma-base-change-f-shriek-separated} 与
\ref{more-etale-lemma-f-shriek-composition} 中的同构均由相应的直像映射构造。
\end{remark}

\begin{lemma}
\label{more-etale-lemma-colim-f-shriek-separated}
设 $f : X \to Y$ 为分离且局部有限型的概形态射。设
$X = \bigcup_{i \in I} X_i$ 为开覆盖，且对任意 $i, j \in I$，都存在
$k$ 使 $X_i \cup X_j \subset X_k$。以 $f_i : X_i \to Y$ 表示 $f$
的限制。则
$$
f_!\mathcal{F} = \colim_{i \in I} f_{i, !}(\mathcal{F}|_{X_i})
$$
关于 $\mathcal{F} \in \textit{Ab}(X_\etale)$ 函子性地成立，其中转移映射
由评注 \ref{more-etale-remark-covariance-f-shriek-separated} 构造。
\end{lemma}

\begin{proof}
只须证明，从右到左的典范映射在 $Y_\etale$ 的拟紧对象 $V$ 上取值时为双射。
注意，右端余极限是有向的，且转移映射均为单射。因此可用“位点”引理
\ref{sites-lemma-directed-colimits-sections} 计算该余极限。于是结论归结为如下观察：
在 $V$ 上固有的闭子集 $Z \subset X_V$ 是拟紧的，故对某个 $i$，
它包含在 $X_{i, V}$ 中。
\end{proof}

\begin{lemma}
\label{more-etale-lemma-f-shriek-separated-direct-sums}
设 $f : X \to Y$ 为分离且局部有限型的概形态射。则函子 $f_!$ 与直和可交换。
\end{lemma}

\begin{proof}
令 $\mathcal{F} = \bigoplus \mathcal{F}_i$。为证明映射
$\bigoplus f_!\mathcal{F}_i \to f_!\mathcal{F}$ 为同构，只须证明这些层在
$Y_\etale$ 的任意拟紧对象 $V$ 上具有相同的截面。用 $V$ 替换 $Y$ 后，
只须证明
$H^0(Y, f_!\mathcal{F}) \subset H^0(X, \mathcal{F})$
等于
$\bigoplus H^0(Y, f_!\mathcal{F}_i)
\subset \bigoplus H^0(X, \mathcal{F}_i)
\subset H^0(X, \bigoplus \mathcal{F}_i)$.
把 $X$ 写成其拟紧开子概形的并，并应用引理
\ref{more-etale-lemma-colim-f-shriek-separated}，便可化归到 $X$ 也拟紧的情形。
此时，由“\'Etale 上同调”定理 \ref{etale-cohomology-theorem-colimit}，有
$H^0(X, \mathcal{F}) = \bigoplus H^0(X, \mathcal{F}_i)$。
考察截面的支撑即可容易得出结论。
\end{proof}

\begin{lemma}
\label{more-etale-lemma-lqf-f-shriek-separated-colimits}
设 $f : X \to Y$ 为分离且局部拟有限的概形态射。则
\begin{enumerate}
\item 对 $\textit{Ab}(X_\etale)$ 中的 $\mathcal{F}$ 及几何点
$\overline{y} : \Spec(k) \to Y$，有
$$
(f_!\mathcal{F})_{\overline{y}} =
\bigoplus\nolimits_{f(\overline{x}) = \overline{y}} \mathcal{F}_{\overline{x}}
$$
关于 $\mathcal{F}$ 函子性地成立；
\item 函子 $f_!$ 是正合的。
\end{enumerate}
\end{lemma}

\begin{proof}
由构造，函子 $f_!$ 左正合。右正合性可在茎上检验（“\'Etale 上同调”定理
\ref{etale-cohomology-theorem-exactness-stalks}）。因此只须证明 (1)。

\medskip\noindent
设 $\overline{y} : \Spec(k) \to Y$ 为几何点。概形 $X_{\overline{y}}$
的底拓扑空间是离散的（“态射”引理
\ref{morphisms-lemma-locally-quasi-finite-fibres}），且各点处的剩余域都等于
$k$（它们是 $k$ 的有限扩张）。因此，集合
$\{\overline{x} : \Spec(k) \to X : f(\overline{x}) = \overline{y}\}$
等于 $X_{\overline{y}}$ 的点集。于是茎的计算由更一般的引理
\ref{more-etale-lemma-stalk-f-shriek-separated} 得出。
\end{proof}









\section{有限支撑截面}
\label{more-etale-section-finite-support}

\noindent
本节把第 \ref{more-etale-section-compact-support} 节的构造推广到未必分离的局部拟有限态射。

\medskip\noindent
设 $f : X \to Y$ 为局部拟有限的概形态射，$\mathcal{F}$ 为
$X_\etale$ 上的阿贝尔层。给定 $Y_\etale$ 中的 $V$，以
$X_V = X \times_Y V$ 表示基变换。我们考虑有限形式和所成的群
\begin{equation}
\label{more-etale-equation-formal-sum}
s = \sum\nolimits_{i = 1, \ldots, n} (Z_i, s_i)
\end{equation}
，其中 $Z_i \subset X_V$ 为局部闭子概形，且态射 $Z_i \to V$ 有限
\footnote{由于 $f$ 局部拟有限，态射 $Z_i \to V$ 有限当且仅当它固有。}，
并且 $s_i \in H_{Z_i}(\mathcal{F})$。这里与第 \ref{more-etale-section-growing} 节一样，令
$$
H_{Z_i}(\mathcal{F}) =
\{s_i \in \mathcal{F}(U_i) \mid \text{Supp}(s_i) \subset Z_i\}
$$
，其中 $U_i \subset X_V$ 为把 $Z_i$ 包含为闭子概形的开子概形。我们对这些
形式和模去以下关系：
\begin{enumerate}
\item
\label{more-etale-item-sum}
$(Z, s) + (Z, s') = (Z, s + s')$,
\item
\label{more-etale-item-sub}
$(Z, s) = (Z', s)$（当 $Z \subset Z'$ 时）。
\end{enumerate}
注意，第二个关系确有意义：由于 $Z \to V$ 有限而 $Z' \to V$ 分离，
包含映射 $Z \to Z'$ 为闭浸入，故可使用 (\ref{more-etale-item-inclusion}) 中讨论的映射。

\medskip\noindent
以 $f_{p!}\mathcal{F}(V)$ 表示形式和 (\ref{more-etale-equation-formal-sum}) 所成阿贝尔群
模去这些关系的商。第一个关系说明，$f_{p!}\mathcal{F}(V)$ 是所有在 $V$
上有限的局部闭子概形 $Z \subset X_V$ 所对应的阿贝尔群
$H_Z(\mathcal{F})$ 之直和的商。第二个关系说明，我们实际上取的是余极限
\begin{equation}
\label{more-etale-equation-colimit-definition}
f_{p!}\mathcal{F}(V) = \colim_Z H_Z(\mathcal{F})
\end{equation}
这个公式为理解该构造提供了方便的抽象方式。

\medskip\noindent
接着注意，该构造可自然地成为 $Y_\etale$ 上的阿贝尔群预层
$f_{p!}\mathcal{F}$。具体地，给定 $Y_\etale$ 中的 $V' \to V$，得到基变换态射
$X_{V'} \to X_V$。若 $Z \subset X_V$ 是在 $V$ 上有限的局部闭子概形，则其
概形论逆像 $Z' \subset X_{V'}$ 在 $V'$ 上有限。此外，若 $U \subset X_V$
为开子概形且 $Z$ 在 $U$ 中闭，则逆像 $U' \subset X_{V'}$ 为开子概形，
且 $Z'$ 在 $U'$ 中闭。因此，$\mathcal{F}$ 的限制映射
$\mathcal{F}(U) \to \mathcal{F}(U')$ 把 $H_Z(\mathcal{F})$ 映入
$H_{Z'}(\mathcal{F})$；这是上文 (\ref{more-etale-item-pullback}) 中所述函子性的特例。
显然，这些映射与 $X_V$ 中此类局部闭子概形的包含
$Z_1 \subset Z_2$ 相容，于是得到映射
$$
f_{p!}\mathcal{F}(V) = \colim_Z H_Z(\mathcal{F})
\longrightarrow
\colim_{Z'} H_{Z'}(\mathcal{F}) =
f_{p!}\mathcal{F}(V')
$$
这些映射确实使 $f_{p!}\mathcal{F}$ 成为 $Y_\etale$ 上的阿贝尔群预层。
略去细节。

\medskip\noindent
最后注意，$f_{p!}\mathcal{F}$ 的构造关于
$\textit{Ab}(X_\etale)$ 中的 $\mathcal{F}$ 具有函子性。因此，给定局部拟有限态射
$f : X \to Y$，我们构造了函子
$$
f_{p!} :
\textit{Ab}(X_\etale)
\longrightarrow
\textit{PAb}(Y_\etale)
$$
，从 $X_\etale$ 上阿贝尔层的范畴到 $Y_\etale$ 上阿贝尔预层的范畴。
在把 $f_!$ 定义为此函子的层化之前，先检验在两种构造均适用时，它与
第 \ref{more-etale-section-compact-support} 节的构造以及“\'Etale 上同调”第
\ref{etale-cohomology-section-extension-by-zero} 节的构造一致。

\begin{lemma}
\label{more-etale-lemma-finite-support-f-shriek-separated}
设 $f : X \to Y$ 为分离且局部拟有限的概形态射。关于
$\mathcal{F} \in \textit{Ab}(X_\etale)$ 函子性地存在阿贝尔预层的典范同构（！）
$$
f_{p!}\mathcal{F} \longrightarrow f_!\mathcal{F}
$$
；它把定义 \ref{more-etale-definition-f-shriek-separated} 中的层
$f_!\mathcal{F}$ 与上述构造的预层 $f_{p!}\mathcal{F}$ 等同起来。
\end{lemma}

\begin{proof}
设 $V$ 为 $Y_\etale$ 的对象。若 $Z \subset X_V$ 局部闭且在 $V$ 上有限，
则因 $f$ 分离，态射 $Z \to X_V$ 为闭浸入。此外，若 $Z_i$（其中
$i = 1, \ldots, n$）为 $X_V$ 中在 $V$ 上有限的闭子概形，则
$Z_1 \cup \ldots \cup Z_n$（概形论并）是一个在 $V$ 上有限的闭子概形。
因此在此情形，定义 $f_{p!}\mathcal{F}(V)$ 的余极限
(\ref{more-etale-equation-colimit-definition}) 是有向的，并且
$f_{!p}\mathcal{F}(V)$ 恰等于 $\mathcal{F}(X_V)$ 中支撑在 $V$ 上有限的截面之集。
由于 $X_V$ 中在 $V$ 上固有的任意闭子集事实上都在 $V$ 上有限
（因为 $f$ 局部拟有限），由定义即知该集合等于 $f_!\mathcal{F}(V)$。
\end{proof}

\begin{lemma}
\label{more-etale-lemma-finite-support-stalk}
设 $f : X \to Y$ 为局部拟有限的概形态射，
$\overline{y} : \Spec(k) \to Y$ 为几何点。关于
$\textit{Ab}(X_\etale)$ 中的 $\mathcal{F}$ 函子性地有
$$
(f_{p!}\mathcal{F})_{\overline{y}} =
\bigoplus\nolimits_{f(\overline{x}) = \overline{y}} \mathcal{F}_{\overline{x}}
$$
\end{lemma}

\begin{proof}
回忆，预层在 $\overline{y}$ 处的茎定义为取遍 $\overline{y}$ 的 \'etale 邻域
$(V, \overline{v})$ 的通常余极限；见“\'Etale 上同调”定义
\ref{etale-cohomology-definition-stalk}。相应地，设
$s = \sum_{i = 1, \ldots, n} (Z_i, s_i)$ 如
(\ref{more-etale-equation-formal-sum}) 所示，是 $f_{p!}\mathcal{F}(V)$ 的元素，其中
$(V, \overline{v})$ 为 $\overline{y}$ 的 \'etale 邻域。由于
$$
X_{\overline{y}} = (X_V)_{\overline{v}} \supset Z_{i, \overline{v}}
$$
，且 $s_i$ 是 $X_V$ 中 $Z_i$ 的某个开邻域上的 $\mathcal{F}$ 截面，
故可把 $s$ 映到
$$
\sum\nolimits_{i = 1, \ldots, n} 
\sum\nolimits_{\overline{x} \in Z_{i, \overline{v}}}
\left(\text{取 }s_i\text{ 的类，位于 }\mathcal{F}_{\overline{x}}\right)
\quad\in\quad
\bigoplus\nolimits_{f(\overline{x}) = \overline{y}} \mathcal{F}_{\overline{x}}
$$
略去如下验证：该构造与限制映射相容；关系 (\ref{more-etale-item-sum})
$(Z, s) + (Z, s') - (Z, s + s')$ 以及当 $Z \subset Z'$ 时的关系
(\ref{more-etale-item-sub}) $(Z, s) - (Z', s)$ 均映为零。于是得到映射
$$
(f_{p!}\mathcal{F})_{\overline{y}}
\longrightarrow
\bigoplus\nolimits_{f(\overline{x}) = \overline{y}} \mathcal{F}_{\overline{x}}
$$

\medskip\noindent
先证明此箭头满射。只须取满足 $f(\overline{x}) = \overline{y}$ 的
$\overline{x}$，并证明直和项 $\mathcal{F}_{\overline{x}}$ 中的任意元素 $s$
属于其像。设 $s$ 由元素 $s \in \mathcal{F}(U)$ 表示，其中
$(U, \overline{u})$ 是 $\overline{x}$ 的 \'etale 邻域。由于 $f$ 局部拟有限，
态射 $U \to Y$ 也局部拟有限。由“关于态射的进一步结果”引理
\ref{more-morphisms-lemma-etale-makes-quasi-finite-finite-multiple-points-var}
，可以找到 $\overline{y}$ 的 \'etale 邻域 $(V, \overline{v})$、开子概形
$$
W \subset U \times_Y V,
$$
以及映到 $\overline{u}$ 和 $\overline{v}$ 的几何点 $\overline{w}$，使得
$W \to V$ 有限，且 $\overline{w}$ 是 $W$ 中唯一映到 $\overline{v}$ 的几何点。
（这里略去在当前使用的几何点语言与该引理陈述所用的点及剩余域扩张语言之间
作转换。）注意，$W \to X_V = X \times_Y V$ 是 \'etale 的。取
$\overline{w}$ 的像 $\overline{w}'$ 的仿射开邻域
$W' \subset X_V$。由于 $\overline{w}$ 是 $W$ 中位于 $\overline{v}$ 上方的
唯一点，且 $W \to V$ 为闭映射，用 $\overline{v}$ 的一个开邻域替换 $V$ 后，
可设 $W \to X_V$ 的像包含在 $W'$ 中。于是 $W \to W'$ 有限且 \'etale，
并且 $W$ 中恰有一个位于 $\overline{w}'$ 上方的几何点 $\overline{w}$。
因此，在 $W'$ 中 $\overline{w}'$ 的某个开邻域上，$W \to W'$ 是开浸入；
见“\'Etale 态射”引理 \ref{etale-lemma-finite-etale-one-point}。缩小 $V$ 与
$W'$ 后，可设 $W \to W'$ 为同构。于是可把 $s$ 看作 $\mathcal{F}$ 在
$W' \subset X_V$ 上的截面 $s'$，且该开子概形在 $V$ 上有限。因此按定义，
$(W', s')$ 给出 $j_{p!}\mathcal{F}(V)$ 的一个元素，并按所需映到 $s$。

\medskip\noindent
再证明此箭头单射。设
$s = \sum_{i = 1, \ldots, n} (Z_i, s_i)$ 如
(\ref{more-etale-equation-formal-sum}) 所示，是 $f_{p!}\mathcal{F}(V)$ 的元素，其中
$(V, \overline{v})$ 为 $\overline{y}$ 的 \'etale 邻域。假设 $s$ 在上述构造的
映射下映为零。首先，用 $(V, \overline{v})$ 自身的一个 \'etale 邻域替换它后，
可设存在分解
$Z_i = Z_{i, 1} \amalg \ldots \amalg Z_{i, m_i}$，其中各项均为开闭
子概形，使每个 $Z_{i, j}$ 在 $\overline{v}$ 上方恰有一个几何点。在显然的直和分解
$$
H_{Z_i}(\mathcal{F}) = \bigoplus H_{Z_{i, j}}(\mathcal{F})
$$
下，设元素 $s_i$ 对应于 $\sum s_{i, j}$。利用关系 (\ref{more-etale-item-sum}) 与
(\ref{more-etale-item-sub})，可把 $s$ 替换为
$\sum_{i = 1, \ldots, n} \sum_{j = 1, \ldots, m_i} (Z_{i, j}, s_{i, j})$.
换言之，可设 $Z_i$ 在 $\overline{v}$ 上方恰有一个几何点。令
$\overline{x}_1, \ldots, \overline{x}_m$ 为 $X$ 中位于 $\overline{y}$ 上方、
与各 $Z_i$ 中位于 $\overline{v}$ 上方的几何点相对应的几何点；注意，对某个
$j \in \{1, \ldots, m\}$，可能有多个指标 $i$，使 $\overline{x}_j$
对应于 $Z_i$ 的一个点。对 $X_V \to V$ 应用“关于态射的进一步结果”引理
\ref{more-morphisms-lemma-etale-makes-quasi-finite-finite-multiple-points-var}
，并用 $(V, \overline{v})$ 自身的一个 \'etale 邻域替换它后，可设存在开子概形
$$
W_j \subset X \times_Y V,\quad j = 1, \ldots, m
$$
以及 $W_j$ 中映到 $\overline{x}_j$ 和 $\overline{v}$ 的几何点
$\overline{w}_j$，使得 $W_j \to V$ 有限，且 $\overline{w}_j$ 是 $W_j$ 中
唯一映到 $\overline{v}$ 的几何点。缩小 $V$ 后，可设对某个 $j$ 有
$Z_i \subset W_j$，并且有映射
$H_{Z_i}(\mathcal{F}) \to H_{W_j}(\mathcal{F})$。因此，由关系
(\ref{more-etale-item-sub})，我们的元素等价于形如
$$
\sum\nolimits_{j = 1, \ldots, m} (W_j, t_j)
$$
的元素，其中 $t_j \in H_{W_j}(\mathcal{F})$。显然，此元素恰映到直和项
$\mathcal{F}_{\overline{x}_j}$ 中 $t_j$ 的类。由于 $s$ 映为零，$t_j$ 在
$\mathcal{F}_{\overline{x}_j}$ 中映为零。这说明 $t_j$ 在 $W_j$ 中
$\overline{w}_j$ 的某个开邻域上的限制为零；见“\'Etale 上同调”引理
\ref{etale-cohomology-lemma-zero-over-image}。再次缩小 $V$，便按所需得到
对所有 $j$ 都有 $t_j = 0$。
\end{proof}

\begin{lemma}
\label{more-etale-lemma-finite-support-etale-shriek}
设 $f = j : U \to X$ 为概形的 \'etale 态射。以 $j_{p!}$ 表示
“\'Etale 上同调”公式 (\ref{etale-cohomology-equation-j-p-shriek})
中的构造，以 $f_{p!}$ 表示上述构造。关于
$\mathcal{F} \in \textit{Ab}(X_\etale)$ 函子性地存在阿贝尔预层的典范映射
$$
j_{p!}\mathcal{F} \longrightarrow f_{p!}\mathcal{F}
$$
；它把“\'Etale 上同调”定义
\ref{etale-cohomology-definition-extension-zero} 中的层
$j_!\mathcal{F} = (j_{p!}\mathcal{F})^\#$ 与 $(f_{p!}\mathcal{F})^\#$ 等同起来。
\end{lemma}

\begin{proof}
阅读本引理的证明前，请先阅读“\'Etale 上同调”引理
\ref{etale-cohomology-lemma-shriek-into-star-separated-etale} 的证明。
设 $V$ 为 $X_\etale$ 的对象。回忆
$$
j_{p!}\mathcal{F}(V) =
\bigoplus\nolimits_{\varphi : V \to U} \mathcal{F}(V \xrightarrow{\varphi} U)
$$
。给定 $\varphi$，得到开子概形 $Z_\varphi \subset U_V = U \times_X V$，
即 $\varphi$ 的图的像。通过 $\varphi$，得到 $U$ 上的同构
$V \to Z_\varphi$，并可把元素
$$
s_\varphi \in \mathcal{F}(V \xrightarrow{\varphi} U) =
\mathcal{F}(Z_\varphi) = H_{Z_\varphi}(\mathcal{F})
$$
看作 $\mathcal{F}$ 在 $Z_{\varphi}$ 上的截面。由于
$Z_\varphi \subset U_V$ 为开子概形，事实上有
$H_{Z_\varphi}(\mathcal{F}) = \mathcal{F}(Z_\varphi)$，故可把
$s_\varphi$ 看作 $H_{Z_\varphi}(\mathcal{F})$ 的元素。这样，映射
$j_{p!}\mathcal{F} \to f_{p!}\mathcal{F}$ 由规则
$$
\sum\nolimits_{i = 1, \ldots, n} s_{\varphi_i}
\longmapsto
\sum\nolimits_{i = 1, \ldots, n} (Z_{\varphi_i}, s_{\varphi_i})
$$
定义，其中右端是 (\ref{more-etale-equation-formal-sum}) 中的形式和。略去它与限制映射相容
且关于 $\mathcal{F}$ 具有函子性的验证。

\medskip\noindent
为完成证明，我们断言，对给定的几何点
$\overline{y} : \Spec(k) \to Y$，存在交换图
$$
\xymatrix{
(j_{p!}\mathcal{F})_{\overline{y}} \ar[r] \ar[d] &
\bigoplus_{j(\overline{x}) = \overline{y}} \mathcal{F}_{\overline{x}}
\ar@{=}[d] \\
(f_{p!}\mathcal{F})_{\overline{y}} \ar[r] &
\bigoplus_{f(\overline{x}) = \overline{y}} \mathcal{F}_{\overline{x}}
}
$$
其中上方横箭头构造于“\'Etale 上同调”命题
\ref{etale-cohomology-proposition-describe-jshriek} 的证明中，下方横箭头构造于
引理 \ref{more-etale-lemma-finite-support-stalk} 的证明中，右侧竖箭头是显然的等号，
左侧竖箭头是上一段所定义映射在茎上的映射。根据这里及所引文献中涉及的
全部箭头的显式描述，该断言直接成立。由于横箭头均为同构，左侧竖箭头也为同构。
因此，由“\'Etale 上同调”定理
\ref{etale-cohomology-theorem-exactness-stalks}，我们的映射在层化上诱导同构。
\end{proof}

\begin{definition}
\label{more-etale-definition-f-shriek-lqf}
设 $f : X \to Y$ 为局部拟有限的概形态射。定义{\it 紧支撑直像}为函子
$$
f_! : \textit{Ab}(X_\etale) \longrightarrow \textit{Ab}(Y_\etale)
$$
，其定义公式为 $f_!\mathcal{F} = (f_{p!}\mathcal{F})^\#$；换言之，
$f_!\mathcal{F}$ 是上述构造的预层 $f_{p!}\mathcal{F}$ 的层化。
\end{definition}

\noindent
由引理 \ref{more-etale-lemma-finite-support-f-shriek-separated}，它与定义
\ref{more-etale-definition-f-shriek-separated} 不冲突（当两个定义均适用时）；由引理
\ref{more-etale-lemma-finite-support-etale-shriek}，它与“\'Etale 上同调”定义
\ref{etale-cohomology-definition-extension-zero} 也不冲突（当两个定义均适用时）。

\begin{lemma}
\label{more-etale-lemma-lqf-f-shriek-stalk}
设 $f : X \to Y$ 为局部拟有限的概形态射。则
\begin{enumerate}
\item 对 $\textit{Ab}(X_\etale)$ 中的 $\mathcal{F}$ 及几何点
$\overline{y} : \Spec(k) \to Y$，有
$$
(f_!\mathcal{F})_{\overline{y}} =
\bigoplus\nolimits_{f(\overline{x}) = \overline{y}} \mathcal{F}_{\overline{x}}
$$
关于 $\mathcal{F}$ 函子性地成立；
\item 函子 $f_! : \textit{Ab}(X_\etale) \to \textit{Ab}(Y_\etale)$
正合且与直和可交换。
\end{enumerate}
\end{lemma}

\begin{proof}
茎的公式由引理 \ref{more-etale-lemma-finite-support-stalk} 立即得到（事实上二者等价）。
再结合正合性可在茎上检验这一事实，便立即得到函子的正合性；见
“\'Etale 上同调”定理 \ref{etale-cohomology-theorem-exactness-stalks}。
\end{proof}

\begin{remark}[关于开嵌入的协变性]
\label{more-etale-remark-covariance-lqf-f-shriek}
设 $f : X \to Y$ 为局部拟有限的概形态射，$\mathcal{F}$ 为
$X_\etale$ 上的阿贝尔层。设 $X' \subset X$ 为开子概形，并以
$f' : X' \to Y$ 表示 $f$ 的限制。我们断言存在典范映射
$$
f'_!(\mathcal{F}|_{X'}) \longrightarrow f_!\mathcal{F}
$$
。具体地，此映射是典范映射
$$
f'_{p!}(\mathcal{F}|_{X'}) \to f_{p!}\mathcal{F}
$$
的层化；后者构造如下。设 $V \in Y_\etale$，并考虑如
(\ref{more-etale-equation-formal-sum}) 所示的截面
$s' = \sum_{i = 1, \ldots, n} (Z'_i, s'_i)$，它定义了
$f'_{p!}(\mathcal{F}|_{X'})(V)$ 的一个元素。此时
$Z'_i \subset X'_V$ 也可看作 $X_V$ 的局部闭子概形，且有
$H_{Z'_i}(\mathcal{F}|_{X'}) = H_{Z'_i}(\mathcal{F})$。把 $s'$ 映到完全相同的和
$s = \sum_{i = 1, \ldots, n} (Z'_i, s'_i)$
，但现在将它看作 $f_{p!}\mathcal{F}(V)$ 的元素。略去该构造与限制映射相容且
关于 $\mathcal{F}$ 具有函子性的验证。此构造具有以下性质：
\begin{enumerate}
\item 映射 $f'_{p!}\mathcal{F}' \to f_{p!}\mathcal{F}$ 与
$f'_!\mathcal{F}' \to f_!\mathcal{F}$ 同引理
\ref{more-etale-lemma-finite-support-stalk} 和 \ref{more-etale-lemma-lqf-f-shriek-stalk}
给出的茎的描述相容。
\item 若 $f$ 分离，则映射
$f'_{p!}\mathcal{F}' \to f_{p!}\mathcal{F}$ 与评注
\ref{more-etale-remark-covariance-f-shriek-separated} 中构造的映射相同；这里使用引理
\ref{more-etale-lemma-finite-support-f-shriek-separated} 中的同构。
\item 若 $X'' \subset X'$ 是另一个开子概形，则复合
$f''_{p!}(\mathcal{F}|_{X''}) \to f'_{p!}(\mathcal{F}|_{X'}) \to
f_{p!}\mathcal{F}$ 是包含 $X'' \subset X$ 所对应的映射
$f''_{p!}(\mathcal{F}|_{X''}) \to f_{p!}\mathcal{F}$。层化后可知，
对于
$f''_!(\mathcal{F}|_{X''}) \to f'_!(\mathcal{F}|_{X'}) \to f_!\mathcal{F}$.
也有同样的结论。
\item 映射 $f'_!\mathcal{F}' \to f_!\mathcal{F}$ 为单射，因为这可在茎上检验。
\end{enumerate}
把元素表示为上述有限和并考察这些元素的像，即可容易证明所有这些陈述。
\end{remark}

\begin{lemma}
\label{more-etale-lemma-lqf-colimit-f-shriek}
设 $f : X \to Y$ 为局部拟有限的概形态射，
$X = \bigcup_{i \in I} X_i$ 为开覆盖。则存在正合复形
$$
\ldots \to
\bigoplus\nolimits_{i_0, i_1, i_2} f_{i_0i_1i_2, !}
\mathcal{F}|_{X_{i_0i_1i_2}} \to
\bigoplus\nolimits_{i_0, i_1} f_{i_0i_1, !} \mathcal{F}|_{X_{i_0i_1}} \to
\bigoplus\nolimits_{i_0} f_{i_0, !} \mathcal{F}|_{X_{i_0}}
\to f_!\mathcal{F} \to 0
$$
，它关于 $\mathcal{F} \in \textit{Ab}(X_\etale)$ 函子性地成立；细节见证明。
\end{lemma}

\begin{proof}
这里照例令 $X_{i_0 \ldots i_p} = X_{i_0} \cap \ldots \cap X_{i_p}$，
并以 $f_{i_0 \ldots i_p}$ 表示 $f$ 在 $X_{i_0 \ldots i_p}$ 上的限制。
复形中的映射是评注 \ref{more-etale-remark-covariance-lqf-f-shriek} 中构造的映射，
符号规则与 {\v C}ech 复形相同。由引理 \ref{more-etale-lemma-lqf-f-shriek-stalk}
对茎的描述，正合性容易得到。略去细节。
\end{proof}

\begin{remark}[另一种构造]
\label{more-etale-remark-alternative-lqf-f-shriek}
引理 \ref{more-etale-lemma-lqf-colimit-f-shriek} 给出了局部拟有限态射 $f$ 的函子
$f_!$ 的另一种构造。具体地，给定概形的局部拟有限态射
$f : X \to Y$，可选取开覆盖 $X = \bigcup_{i \in I} X_i$，使每个
$f_i : X_i \to Y$ 都分离。例如可取 $X$ 的仿射开覆盖。于是可把
$f_!\mathcal{F}$ 定义为该引理复形倒数第二个映射的余核，即
$$
f_!\mathcal{F} = \Coker\left(
\bigoplus\nolimits_{i_0, i_1} f_{i_0i_1, !} \mathcal{F}|_{X_{i_0i_1}} \to
\bigoplus\nolimits_{i_0} f_{i_0, !} \mathcal{F}|_{X_{i_0}}
\right)
$$
。由于态射 $f_{i_0}$ 与 $f_{i_0 i_1}$ 分离，可使用第
\ref{more-etale-section-compact-support} 节中 $f_{i_0, !}$ 和 $f_{i_0i_1, !}$ 的构造。
随后可计算 $f_!$ 的茎（使用分离情形，即引理
\ref{more-etale-lemma-lqf-f-shriek-separated-colimits}），并得到引理
\ref{more-etale-lemma-lqf-f-shriek-stalk} 的结论。完成这一步后，本节其他全部结果也可由此推出。
\end{remark}

\begin{remark}
\label{more-etale-remark-construct-map-presheaves-downstairs}
设 $g : Y' \to Y$ 为概形态射。对 $Y'_\etale$ 上的阿贝尔预层
$\mathcal{G}'$，以 $g_*\mathcal{G}'$ 表示预层
$V \mapsto \mathcal{G}'(Y' \times_Y V)$。若
$\alpha : \mathcal{G} \to g_*\mathcal{G}'$ 是 $Y_\etale$ 上阿贝尔预层的映射，
则存在唯一的 $Y_\etale$ 上阿贝尔层映射
$\alpha^\# : \mathcal{G}^\# \to g_*((\mathcal{G}')^\#)$
使图
$$
\xymatrix{
\mathcal{G} \ar[d] \ar[r]_\alpha & g_*\mathcal{G}' \ar[d] \\
\mathcal{G}^\# \ar[r]^-{\alpha^\#} & g_*((\mathcal{G}')^\#)
}
$$
交换，其中竖映射来自典范映射 $\mathcal{G} \to \mathcal{G}^\#$ 与
$\mathcal{G}' \to (\mathcal{G}')^\#$。若
$\alpha' : g^{-1}\mathcal{G}^\# \to (\mathcal{G}')^\#$
是 $\alpha^\#$ 的伴随映射，则对几何点
$\overline{y}' : \Spec(k) \to Y'$，其在 $Y$ 中的像为
$\overline{y} = g \circ \overline{y}'$，映射
$$
\alpha'_{\overline{y}'} :
\mathcal{G}_{\overline{y}} =
(\mathcal{G}^\#)_{\overline{y}} =
(g^{-1}\mathcal{G}^\#)_{\overline{y}'}
\longrightarrow
(\mathcal{G}')^\#_{\overline{y}'} =
\mathcal{G}'_{\overline{y}'}
$$
按如下方式给出：把 $\mathcal{G}$ 在 \'etale 邻域
$(V, \overline{v})$ 上的截面 $s$ 在茎中的类，映到截面
$\alpha(s)$ 在 $g_*\mathcal{G}'(V) = \mathcal{G}'(Y' \times_Y V)$ 中的类；
后者取于 \'etale 邻域 $(Y' \times_Y V, (\overline{y}', \overline{v}))$
上，并视为 $\mathcal{G}'$ 在 $\overline{y}'$ 处茎中的类。
\end{remark}

\begin{lemma}
\label{more-etale-lemma-lqf-base-change-f-shriek}
考虑概形的笛卡尔方块
$$
\xymatrix{
X' \ar[r]_{g'} \ar[d]_{f'} & X \ar[d]^f \\
Y' \ar[r]^g & Y
}
$$
，其中 $f$ 局部拟有限。存在关于 $\textit{Ab}(X_\etale)$ 中的
$\mathcal{F}$ 函子性的同构
$g^{-1}f_!\mathcal{F} \to f'_!(g')^{-1}\mathcal{F}$；它与引理
\ref{more-etale-lemma-lqf-f-shriek-stalk} 给出的茎的描述相容（精确陈述见证明）。
\end{lemma}

\begin{proof}
采用评注 \ref{more-etale-remark-construct-map-presheaves-downstairs} 中的约定，我们显式构造
$Y_\etale$ 上阿贝尔预层的映射
$$
c : f_{p!}\mathcal{F} \longrightarrow g_*f'_{p!}(g')^{-1}\mathcal{F}
$$
。由评注 \ref{more-etale-remark-construct-map-presheaves-downstairs} 中的讨论，这确定了典范映射
$g^{-1}f_!\mathcal{F} \to f'_!(g')^{-1}\mathcal{F}$。最后，我们将证明此映射
在茎上诱导同构，再由“\'Etale 上同调”定理
\ref{etale-cohomology-theorem-exactness-stalks} 得出结论。

\medskip\noindent
构造映射 $c$。设 $V \in Y_\etale$，并考虑如
(\ref{more-etale-equation-formal-sum}) 所示的截面
$s = \sum_{i = 1, \ldots, n} (Z_i, s_i)$，它定义了
$f_{p!}\mathcal{F}(V)$ 的一个元素。$g_*f'_{p!}(g')^{-1}\mathcal{F}$ 在 $V$
上的值为 $f'_{p!}(g')^{-1}\mathcal{F}(V')$，其中
$V' = V \times_Y Y'$。以 $Z'_i \subset X'_{V'}$ 表示 $Z_i$ 到 $V'$ 的基变换。
由 (\ref{more-etale-item-pullback})，存在拉回映射
$H_{Z_i}(\mathcal{F}) \to H_{Z'_i}((g')^{-1}\mathcal{F})$.
以 $s'_i \in H_{Z'_i}((g')^{-1}\mathcal{F})$ 表示 $s_i$ 在拉回下的像，并令
$c(s) = \sum_{i = 1, \ldots, n} (Z'_i, s'_i)$；它如
(\ref{more-etale-equation-formal-sum}) 所示，定义了
$f'_{p!}(g')^{-1}\mathcal{F}(V')$ 的一个元素。略去如下验证：此构造与关系
(\ref{more-etale-item-sum})、(\ref{more-etale-item-sub}) 以及限制映射相容。该构造显然关于
$\mathcal{F}$ 具有函子性。

\medskip\noindent
设 $\overline{y}' : \Spec(k) \to Y'$ 为几何点，其在 $Y$ 中的像为
$\overline{y} = g \circ \overline{y}'$。由纤维积的传递性，有
$X'_{\overline{y}'} = X_{\overline{y}}$。因此，$g'$ 给出双射
$\{f'(\overline{x}') = \overline{y}'\} \to \{f(\overline{x}) = \overline{y}\}$
；若 $\overline{x}'$ 映到 $\overline{x}$，则
$((g')^{-1}\mathcal{F})_{\overline{x}'} = \mathcal{F}_{\overline{x}}$
；这由“\'Etale 上同调”引理 \ref{etale-cohomology-lemma-stalk-pullback} 得到。
现在断言图
$$
\xymatrix{
(g^{-1}f_!\mathcal{F})_{\overline{y}'} \ar@{=}[r] \ar[d] &
(f_!\mathcal{F})_{\overline{y}} \ar[r] \ar[ld] &
\bigoplus\nolimits_{f(\overline{x}) = \overline{y}}
\mathcal{F}_{\overline{x}} \ar[d]
\\
(f'_!(g')^{-1}\mathcal{F})_{\overline{y}'} \ar[rr] & &
\bigoplus\nolimits_{f'(\overline{x}') = \overline{y}'}
(g')^{-1}\mathcal{F}_{\overline{x}'}
}
$$
交换，其中横箭头由引理 \ref{more-etale-lemma-finite-support-stalk} 的证明给出，
右侧竖箭头依上述讨论是等号。西南方向的箭头在评注
\ref{more-etale-remark-construct-map-presheaves-downstairs} 中描述为拉回映射，即恰由上述构造
$c$ 给出。引理 \ref{more-etale-lemma-finite-support-stalk} 的证明对和
$\sum (Z_i, z_i)$ 在 $\overline{x}$ 处茎中之像的简单描述，立即说明该图交换。
本引理得证。
\end{proof}

\begin{lemma}
\label{more-etale-lemma-lqf-separated-shriek-composition}
设 $f' : X \to Y'$ 与 $g : Y' \to Y$ 为可复合的概形态射，其中 $f'$ 和
$f = g \circ f'$ 局部拟有限，且 $g$ 分离并局部有限型。则存在函子的典范同构
$g_! \circ f'_! = f_!$。此同构与以下各项相容：
\begin{enumerate}
\item[(a)] 评注 \ref{more-etale-remark-covariance-f-shriek-separated} 与
\ref{more-etale-remark-covariance-lqf-f-shriek} 所述关于开嵌入的协变性；
\item[(b)] 引理 \ref{more-etale-lemma-lqf-base-change-f-shriek} 与
\ref{more-etale-lemma-base-change-f-shriek-separated} 的基变换同构；
\item[(c)] 当 $f'$ 分离时，它等于引理
\ref{more-etale-lemma-f-shriek-composition} 的同构；这里使用引理
\ref{more-etale-lemma-finite-support-f-shriek-separated} 的等同。
\end{enumerate}
\end{lemma}

\begin{proof}
设 $\mathcal{F}$ 为 $X_\etale$ 上的阿贝尔层。采用评注
\ref{more-etale-remark-construct-map-presheaves-downstairs} 中的约定，我们将显式构造
$Y_\etale$ 上阿贝尔预层的映射
$$
c : f_{p!}\mathcal{F} \longrightarrow g_*f'_{p!}\mathcal{F}
$$
。由评注 \ref{more-etale-remark-construct-map-presheaves-downstairs} 中的讨论，这确定了典范映射
$c^\# : f_!\mathcal{F} \to g_*f'_!\mathcal{F}$。我们将证明 $c^\#$ 的像包含在
子层 $g_!f'_!\mathcal{F}$ 中，从而得到映射
$c' : f_!\mathcal{F} \to g_!f'_!\mathcal{F}$。接着证明 (a)、(b) 和 (c)。
最后，(b) 将使我们能够证明 $c'$ 为同构。

\medskip\noindent
构造映射 $c$。设 $V \in Y_\etale$，并设
$s = \sum (Z_i, s_i)$ 为如 (\ref{more-etale-equation-formal-sum}) 所示的和，
它定义了 $f_{p!}\mathcal{F}(V)$ 的一个元素。回忆，
$Z_i \subset X_V = X \times_Y V$ 是在 $V$ 上有限的局部闭子概形。
令 $V' = Y' \times_Y V$，则 $X_{V'} = X \times_{Y'} V' = X_V$。
因此 $Z_i \subset X_{V'}$ 局部闭；又因 $g$ 分离，$Z_i$ 在 $V'$ 上有限
（“态射”引理 \ref{morphisms-lemma-finite-permanence}）。故可令
$c(s) = \sum (Z_i, s_i)$，但现在把它看作
$f'_{p!}\mathcal{F}(V') = (g_*f'_{p!}\mathcal{F})(V)$ 的元素。此构造显然与关系
(\ref{more-etale-item-sum})、(\ref{more-etale-item-sub}) 以及限制映射相容，因而得到映射 $c$。

\medskip\noindent
注意，在上述讨论中，$f'_!\mathcal{F}$ 在 $V'$ 上的截面
$c(s) = \sum (Z_i, s_i)$ 在
$V' \setminus \Im(\coprod Z_i \to V')$ 上限制为零。由于
$\Im(\coprod Z_i \to V')$ 在 $V$ 上固有（例如由“态射”引理
\ref{morphisms-lemma-scheme-theoretic-image-is-proper}），可知 $c(s)$ 定义了
$g_!f'_!\mathcal{F} \subset g_*f'_!\mathcal{F}$ 在 $V$ 上的截面。
由于 $f_!\mathcal{F}$ 的每个局部截面在局部都来自 $f_{p!}\mathcal{F}$
的局部截面，可知 $c^\#$ 的像包含在 $g_!f'_!\mathcal{F}$ 中。因此得到诱导映射
$c' : f_!\mathcal{F} \to g_!f'_!\mathcal{F}$，它按证明首段所述分解 $c^\#$。

\medskip\noindent
证明 (a)。设 $Y'_1 \subset Y'$ 为开子概形，并令
$X_1 = (f')^{-1}(W')$。得到图
$$
\xymatrix{
X_1 \ar[d]_{f'_1} \ar[r]_a \ar@/_2em/[dd]_{f_1} &
X \ar[d]^{f'} \ar@/^2em/[dd]^f \\
Y'_1 \ar[d]_{g_1} \ar[r]_{b'} &
Y' \ar[d]^g \\
Y \ar@{=}[r] &
Y
}
$$
，其中横箭头均为开浸入。我们断言图
$$
\xymatrix{
f_{1, !}\mathcal{F}|_{X_1} \ar[r]_{c'_1} \ar[dd] &
g_{1, !}f'_{1, !}\mathcal{F}|_{X_1} \ar@{=}[d] \\
& g_{1, !}(f'_!\mathcal{F})|_{Y'_1} \ar[d] \\
f_!\mathcal{F} \ar[r]^{c'} &
g_!f'_!\mathcal{F} \ar[r] & g_*f'_!\mathcal{F}
}
$$
交换，其中左侧竖箭头来自评注 \ref{more-etale-remark-covariance-lqf-f-shriek}，
右侧竖箭头来自评注 \ref{more-etale-remark-covariance-f-shriek-separated}。
图中的等号来自如下事实：$f'_1$ 是 $f'$ 在 $Y'_1$ 上的限制，且我们对
$f'_!$ 的构造在基底上是局部的。最后，为证明交换性，取 $Y_\etale$ 的对象
$V$ 以及如 (\ref{more-etale-equation-formal-sum}) 所示的形式和
$s_1 = \sum (Z_{1, i}, s_{1, i})$，它定义了
$f_{1, p!}\mathcal{F}|_{X_1}(V)$ 的一个元素。回忆，这意味着
$Z_{1, i} \subset X_1 \times_Y V$ 局部闭且在 $V$ 上有限，并且
$s_{1, i} \in H_{Z_{1, i}}(\mathcal{F})$。沿所涉及的映射追踪此截面，
只须证明沿图的两条路径最终得到
$g_*f'_!\mathcal{F}(V) = f'_!\mathcal{F}(Y' \times_Y V)$.
中的同一元素。读者立即可见，两条路径均得到元素
$\sum (Z_{1, i}, s_{1, i})$，其中现在把 $Z_{1, i}$ 看作
$X \times_{Y'} (Y' \times_Y V) = X \times_Y V$ 的局部闭子概形，
且它在 $Y' \times_Y V$ 上有限。

\medskip\noindent
证明 (b)。设 $b : Y_1 \to Y$ 为概形态射。构造交换图
$$
\xymatrix{
X_1 \ar[d]_{f'_1} \ar[r]_a \ar@/_2em/[dd]_{f_1} &
X \ar[d]^{f'} \ar@/^2em/[dd]^f \\
Y'_1 \ar[d]_{g_1} \ar[r]_{b'} &
Y' \ar[d]^g \\
Y_1 \ar[r]^b &
Y
}
$$
，其中各方块为笛卡尔方块。我们断言该构造与引理
\ref{more-etale-lemma-lqf-base-change-f-shriek} 和
\ref{more-etale-lemma-base-change-f-shriek-separated} 的基变换映射相容；即下图的上方矩形
$$
\xymatrix{
b^{-1}f_!\mathcal{F} \ar[rr] \ar[d]_{b^{-1}c'} & &
f_{1, !}a^{-1}\mathcal{F} \ar[d]^{c_1'} \\
b^{-1}g_!f'_!\mathcal{F} \ar[r] \ar[d] &
g_{1, !}(b')^{-1}f'_!\mathcal{F} \ar[r] \ar[d] &
g_{1, !}f'_{1, !}a^{-1}\mathcal{F} \ar[d] \\
b^{-1}g_*f'_!\mathcal{F} \ar[r] &
g_{1, *}(b')^{-1}f'_!\mathcal{F} \ar[r] &
g_{1, *}f'_{1, !}a^{-1}\mathcal{F}
}
$$
交换。验证完全例行，我们建议读者略过。由于从中间一行到最下一行的箭头均为
单射，只须证明外框交换。为此，只须取 $b^{-1}f_!\mathcal{F}$ 的局部截面，
并证明沿任一路径最终得到 $g_{1, *}f'_{1, !}a^{-1}\mathcal{F}$ 的同一局部截面。
事实上，只需对如下局部截面作检验：它是上述 $f_{p!}\mathcal{F}(V)$ 的截面
$s = \sum (Z_i, s_i)$ 沿 $b$ 的拉回（因为这些拉回生成阿贝尔层
$b^{-1}f_!\mathcal{F}$）。以 $V_1$、$V'_1$ 和 $Z_{1, i}$ 分别表示
$V$、$V' = Y' \times_Y V$ 和 $Z_i$ 沿 $Y_1 \to Y$ 的基变换。回忆，
$Z_i$ 是 $X_V = X_{V'}$ 的局部闭子概形，因而 $Z_{1, i}$ 是
$(X_1)_{V_1} = (X_1)_{V'_1}$ 的局部闭子概形。此时 $b^{-1}c'$ 把 $s$
的拉回映到局部截面 $c(s) \sum (Z_i, s_i)$ 的拉回，其中后者视为
$f'_{p!}\mathcal{F}(V') = (g_*f'_{p!}\mathcal{F})(V)$ 的元素。
下方两个基变换映射的复合，恰把它映到
$\sum (Z_{i, 1}, s_{1, i})$，后者视为
$f'_{1, p!}a^{-1}\mathcal{F}(V'_1) =
g_{1, *}f'_{1, p!}a^{-1}\mathcal{F}(V_1)$ 的元素。另一方面，图上方的基变换映射
把 $s$ 的拉回映到 $\sum (Z_{1, i}, s_{1, i})$，后者视为
$f_{1, !}a^{-1}\mathcal{F}(V_1)$ 的元素。最后，按 $c'_1$ 的构造，它确实把
该元素映到 $\sum (Z_{i, 1}, s_{1, i})$，后者视为
$f'_{1, p!}a^{-1}\mathcal{F}(V'_1) =
g_{1, *}f'_{1, p!}a^{-1}\mathcal{F}(V_1)$ 的元素。交换性得证。

\medskip\noindent
证明 (c)。比较两个映射的定义即可。略去细节。

\medskip\noindent
为完成证明，只须证明 $c'$ 沿任意几何点
$\overline{y} : \Spec(k) \to Y$ 的拉回为同构。沿 $\overline{y}$ 拉回等同于取
$\overline{y}$ 处的茎（“\'Etale 上同调”评注
\ref{etale-cohomology-remark-stalk-pullback}），因而可应用“\'Etale 上同调”定理
\ref{etale-cohomology-theorem-exactness-stalks}。由刚证明的相容性 (b)，
可设 $Y$ 是 $k$ 的谱，而我们须证明 $c'$ 为同构。只须证明诱导映射
$$
\bigoplus\nolimits_{x \in X} \mathcal{F}_x = H^0(Y, f_!\mathcal{F})
\longrightarrow
H^0(Y, g_!f'_!\mathcal{F}) = H^0_c(Y', f'_!\mathcal{F})
$$
为同构。等号由引理 \ref{more-etale-lemma-lqf-f-shriek-stalk} 与
\ref{more-etale-lemma-stalk-f-shriek-separated} 成立。回忆，$X$ 是剩余域为 $k$ 的
Artin 局部环之谱的不交并；见“簇”引理
\ref{varieties-lemma-algebraic-scheme-dim-0}。由于左右两端都与直和可交换
（略去细节），可设 $\mathcal{F}$ 是支撑在某个 $x \in X$ 处的摩天层
$x_*A$。由引理 \ref{more-etale-lemma-lqf-f-shriek-stalk}，$f'_!\mathcal{F}$ 是
$x$ 在 $Y$ 中的像 $y'$ 处的摩天层。在这种情形下，我们的构造显然给出恒等映射
$A \to H^0_c(Y', y'_*A) = A$，正合所需。
\end{proof}

\begin{lemma}
\label{more-etale-lemma-lqf-shriek-composition}
设 $f : X \to Y$ 与 $g : Y \to Z$ 为可复合的局部拟有限概形态射。
则存在函子的典范同构
$$
(g \circ f)_! \longrightarrow g_! \circ f_!
$$
。这些同构满足以下性质：
\begin{enumerate}
\item 若 $f$ 与 $g$ 分离，则此同构与引理
\ref{more-etale-lemma-f-shriek-composition} 的同构一致。
\item 若 $g$ 分离，则此同构与引理
\ref{more-etale-lemma-lqf-separated-shriek-composition} 的同构一致。
\item 对几何点 $\overline{z} : \Spec(k) \to Z$，图
$$
\xymatrix{
((g \circ f)_!\mathcal{F})_{\overline{z}} \ar[d] \ar[rr] & &
\bigoplus\nolimits_{g(f(\overline{x})) = \overline{z}}
\mathcal{F}_{\overline{x}} \ar@{=}[d] \\
(g_!f_!\mathcal{F})_{\overline{z}} \ar[r] &
\bigoplus\nolimits_{g(\overline{y}) = \overline{z}}
(f_!\mathcal{F})_{\overline{y}} \ar[r] &
\bigoplus\nolimits_{g(f(\overline{x})) = \overline{z}}
\mathcal{F}_{\overline{x}}
}
$$
交换，其中横箭头由引理 \ref{more-etale-lemma-lqf-f-shriek-stalk} 给出。
\item 设 $h : Z \to T$ 为第三个局部拟有限概形态射。则图
$$
\xymatrix{
(h \circ g \circ f)_! \ar[r] \ar[d] &
(h \circ g)_! \circ f_! \ar[d] \\
h_! \circ (g \circ f)_! \ar[r] &
h_! \circ g_! \circ f_!
}
$$
交换。
\item 假设有概形图
$$
\xymatrix{
X' \ar[d]_{f'} \ar[r]_c & X \ar[d]^f \\
Y' \ar[d]_{g'} \ar[r]_b & Y \ar[d]^g \\
Z' \ar[r]^a & Z
}
$$
，其中两个方块均为笛卡尔方块，且 $f$ 与 $g$ 局部拟有限。则图
$$
\xymatrix{
a^{-1} \circ (g \circ f)_! \ar[d] \ar[rr] & &
(g' \circ f')_! \circ c^{-1} \ar[d] \\
a^{-1} \circ g_! \circ f_! \ar[r] &
g'_! \circ b^{-1} \circ f_! \ar[r] &
g'_! \circ f'_! \circ c^{-1}
}
$$
交换，其中横箭头来自引理 \ref{more-etale-lemma-lqf-base-change-f-shriek}。
\end{enumerate}
\end{lemma}

\begin{proof}
若 $f$ 与 $g$ 分离，则这是引理 \ref{more-etale-lemma-f-shriek-composition} 的特例。
若 $g$ 分离，则这是引理 \ref{more-etale-lemma-lqf-separated-shriek-composition}
的特例；而且它与 $f$、$g$ 均分离的情形一致。

\medskip\noindent
一般情形下的构造。选取开覆盖 $Y = \bigcup Y_i$，使 $g$ 的限制
$g_i : Y_i \to Z$ 分离。令 $X_i = f^{-1}(Y_i)$，并以
$f_i : X_i \to Y_i$ 表示 $f$ 的限制。再令 $h = g \circ f$，并以
$h_i : X_i \to Z$ 表示 $h$ 的限制。考虑图
$$
\xymatrix{
\bigoplus\nolimits_{i_0, i_1}
h_{i_0i_1, !}\mathcal{F}|_{X_{i_0i_1}} \ar[r] \ar[d] &
\bigoplus\nolimits_{i_0} h_{i_0, !}\mathcal{F}|_{X_{i_0}} \ar[r] \ar[d] &
h_!\mathcal{F} \ar[r] \ar@{..>}[dd] &
0 \\
\bigoplus\nolimits_{i_0, i_1}
g_{i_0i_1, !} f_{i_0i_1, !}\mathcal{F}|_{X_{i_0i_1}} \ar[r] \ar[d] &
\bigoplus\nolimits_{i_0}
g_{i_0, !} f_{i_0, !}\mathcal{F}|_{X_{i_0}} \ar[d] \\
\bigoplus\nolimits_{i_0, i_1}
g_{i_0i_1, !} (f_!\mathcal{F})|_{Y_{i_0i_1}} \ar[r] &
\bigoplus\nolimits_{i_0}
g_{i_0, !} (f_!\mathcal{F})|_{Y_{i_0}} \ar[r] &
g_!f_!\mathcal{F} \ar[r] &
0
}
$$
由引理 \ref{more-etale-lemma-lqf-colimit-f-shriek}，图的最上行与最下行均正合。
由引理 \ref{more-etale-lemma-lqf-separated-shriek-composition}，左上方块交换。
左下方块中的竖箭头来自等式
$(f_!\mathcal{F})|_{Y_{i_0i_1}} = f_{i_0i_1, !}\mathcal{F}|_{X_{i_0i_1}}$ 与
$(f_!\mathcal{F})|_{Y_{i_0}} = f_{i_0, !}\mathcal{F}|_{X_{i_0}}$
，因为 $f_!$ 的构造在基底上是局部的。此外，这些等式当然与以下等同相容：
$((f_!\mathcal{F})|_{Y_{i_0}})|_{Y_{i_0i_1}} =
(f_!\mathcal{F})|_{Y_{i_0i_1}}$ 以及
$(f_{i_0, !}\mathcal{F}|_{X_{i_0}})|_{Y_{i_0i_1}} =
f_{i_0i_1, !}\mathcal{F}|_{X_{i_0i_1}}$
。这些等同连同关于开嵌入 $Y_{i_0i_1} \subset Y_{i_0}$ 的协变性一起，
用于定义左下方块的横映射。因此这个方块也交换。由此可知，存在图中所示的
唯一虚线箭头，而且该箭头是同构。

\medskip\noindent
证明性质 (1)--(5)。固定开覆盖 $Y = \bigcup Y_i$。注意，若
$Y \to Z$ 恰好分离，则使用引理
\ref{more-etale-lemma-lqf-separated-shriek-composition} 的映射，便得到适配于上述大图的
虚线箭头（由该引理自身的性质）。这证明了 (2)；再由引理
\ref{more-etale-lemma-lqf-separated-shriek-composition} 与引理
\ref{more-etale-lemma-f-shriek-composition} 的映射之相容性，也证明了 (1)。接着，对 $Z$
上的任意概形 $Z'$，采用开覆盖 $Y' = \bigcup b^{-1}(Y_i)$ 构造的映射
$(g' \circ f')_! \to g'_! \circ f'_!$ 具有 (5) 中的相容性。这由引理
\ref{more-etale-lemma-lqf-separated-shriek-composition} 中所构造映射的相应相容性立即得到。
特别地，可考虑几何点 $\overline{z} : \Spec(k) \to Z$。由于
$X_{\overline{z}} \to Y_{\overline{z}} \to \Spec(k)$ 均为分离态射，
映射
$(g \circ f)_!\mathcal{F} \to g_! f_! \mathcal{F}$
沿 $\overline{z}$ 的基变换等于引理 \ref{more-etale-lemma-f-shriek-composition} 的映射。
读者立即可见，由此得到性质 (3)。当然，性质 (3) 保证采用开覆盖
$Y = \bigcup Y_i$ 构造的函子变换
$(g \circ f)_! \to g_! \circ f_!$ 不依赖于该开覆盖的选择。最后，利用已证明的
性质 (3) 考察茎，即得性质 (4)。
\end{proof}








\section{局部拟有限态射的权函数与迹映射}
\label{more-etale-section-weightings}

\noindent
本节可参考 \cite[报告 XVII，命题 6.2.5]{SGA4}。

\medskip\noindent
设 $f : X \to Y$ 为局部拟有限的概形态射。设
$w : X \to \mathbf{Z}$ 为 $f$ 的权函数；见“关于态射的进一步结果”定义
\ref{more-morphisms-definition-weighting}。设 $\mathcal{F}$ 为
$Y_\etale$ 上的阿贝尔层。本节将证明存在 $Y_\etale$ 上阿贝尔层的映射
$$
\text{Tr}_{f, w, \mathcal{F}} :
f_!f^{-1}\mathcal{F}
\longrightarrow
\mathcal{F}
$$
，它由以下性质刻画：在 $Y$ 的几何点 $\overline{y}$ 处的茎上，得到映射
$$
\bigoplus\nolimits_{f(\overline{x}) = \overline{y}} w(\overline{x}) :
(f_!f^{-1}\mathcal{F})_{\overline{y}} =
\bigoplus\nolimits_{f(\overline{x}) = \overline{y}}
\mathcal{F}_{\overline{y}}
\longrightarrow
\mathcal{F}_{\overline{y}}
$$
如此处所示，在对应于 $\overline{x}$ 的直和项上，该箭头由乘以整数
$w(\overline{x})$ 给出。箭头左侧的等号由引理
\ref{more-etale-lemma-lqf-f-shriek-stalk} 与“\'Etale 上同调”引理
\ref{etale-cohomology-lemma-stalk-pullback} 得到。

\medskip\noindent
若态射 $f : X \to Y$ 平坦、局部拟有限且局部有限表示，则存在典范权函数，
并得到其构造与基变换相容的典范迹映射；见例
\ref{more-etale-example-trace-for-flat-quasi-finite}。若 $Y$ 为局部 Noether 单枝概形，
且 $f : X \to Y$ 局部拟有限，则也可为 $f$ 定义一个（自然的）权函数，
从而在此情形下也有迹映射；见例
\ref{more-etale-example-trace-for-quasi-finite-over-normal}。

\begin{lemma}
\label{more-etale-lemma-f-shriek-projection}
设 $f : X \to Y$ 为局部拟有限的概形态射，$\Lambda$ 为环。
设 $\mathcal{F}$ 为 $X_\etale$ 上的 $\Lambda$-模层，$\mathcal{G}$ 为
$Y_\etale$ 上的 $\Lambda$-模层。存在 $Y_\etale$ 上 $\Lambda$-模层的典范同构
$$
can :
f_!\mathcal{F} \otimes_\Lambda \mathcal{G}
\longrightarrow
f_!(\mathcal{F} \otimes_\Lambda f^{-1}\mathcal{G})
$$
\end{lemma}

\begin{proof}
回忆，由定义 \ref{more-etale-definition-f-shriek-lqf}，有
$f_!\mathcal{F} = (f_{p!}\mathcal{F})^\#$，其中 $f_{p!}\mathcal{F}$ 是第
\ref{more-etale-section-finite-support} 节构造的预层。因此，为构造该箭头，只须构造
$Y_\etale$ 上预层的映射
$$
f_{p!}\mathcal{F} \otimes_{p, \Lambda} \mathcal{G}
\longrightarrow
f_{p!}(\mathcal{F} \otimes_\Lambda f^{-1}\mathcal{G})
$$
。这里符号 $\otimes_{p, \Lambda}$ 表示预层张量积；见“位点上的模”第
\ref{sites-modules-section-tensor-product} 节。设 $V$ 为 $Y_\etale$ 的对象。
回忆
$$
f_{p!}\mathcal{F}(V) = \colim_Z H_Z(\mathcal{F})
\quad\text{且}\quad
f_{p!}(\mathcal{F} \otimes_\Lambda f^{-1}\mathcal{G})(V) =
\colim_Z H_Z(\mathcal{F} \otimes_\Lambda f^{-1}\mathcal{G})
$$
；见第 \ref{more-etale-section-finite-support} 节。我们的映射在纯张量上由规则
$$
(Z, s) \otimes t \longmapsto (Z, s \otimes f^{-1}t)
$$
定义（记号见下文），再线性延拓到整个
$(f_{p!}\mathcal{F} \otimes_{p, \Lambda} \mathcal{G})(V) =
f_{p!}\mathcal{F}(V) \otimes_\Lambda \mathcal{G}(V)$.
所用记号如下：
\begin{enumerate}
\item $Z \subset X_V$ 是在 $V$ 上有限的局部闭子概形；
\item $s \in H_Z(\mathcal{F})$，这意味着对某个开子概形
$U \subset X_V$，其中 $Z \subset U$ 闭，有 $s \in \mathcal{F}(U)$ 且
$\text{Supp}(s) \subset Z$；
\item $t \in \mathcal{G}(V)$，其像为
$f^{-1}t \in f^{-1}\mathcal{G}(U)$。
\end{enumerate}
由于 $s \in \mathcal{F}(U)$ 的支撑包含在 $Z$ 中，
$s \otimes f^{-1}t$ 的支撑显然也包含在 $Z$ 中。因此，有序对
$(Z, s \otimes f^{-1}t)$ 确有意义。立即可见，该构造与余极限
$\colim_Z H_Z(\mathcal{F})$ 的转移映射可交换，且与限制映射相容。
最后，同样显然的是，该构造与引理 \ref{more-etale-lemma-lqf-f-shriek-stalk}
中 $f_!$ 的茎的等同相容。换言之，我们构造的映射 $can$ 在几何点
$\overline{y}$ 处的茎上适配于交换图
$$
\xymatrix{
(f_!\mathcal{F} \otimes_\Lambda \mathcal{G})_{\overline{y}}
\ar[r]_-{can_{\overline{y}}} \ar[d] &
f_!(\mathcal{F} \otimes_\Lambda f^{-1}\mathcal{G})_{\overline{y}} \ar[d] \\
(\bigoplus \mathcal{F}_{\overline{x}})
\otimes_\Lambda \mathcal{G}_{\overline{y}} \ar[r] &
\bigoplus
(\mathcal{F}_{\overline{x}} \otimes_\Lambda \mathcal{G}_{\overline{y}})
}
$$
，其中直和取遍位于 $\overline{y}$ 上方的几何点 $\overline{x}$，竖箭头是引理
\ref{more-etale-lemma-lqf-f-shriek-stalk} 的等同，而下方横箭头是显然的同构。
因此 $can$ 是同构，正合所需。
\end{proof}

\begin{lemma}
\label{more-etale-lemma-trace-map-exists}
设 $f : X \to Y$ 为局部拟有限的概形态射，
$w : X \to \mathbf{Z}$ 为 $f$ 的权函数。对 $Y$ 上任意阿贝尔层
$\mathcal{F}$，存在唯一的迹映射
$\text{Tr}_{f, w, \mathcal{F}} : f_!f^{-1}\mathcal{F} \to \mathcal{F}$
，它在茎上具有预先指定的行为。
\end{lemma}

\begin{proof}
由引理 \ref{more-etale-lemma-f-shriek-projection}，有等同
$f_!f^{-1}\mathcal{F} = f_!\underline{\mathbf{Z}} \otimes \mathcal{F}$
，它与这些层在几何点处茎的描述相容。因此，只须构造映射
$$
\text{Tr}_{f, w, \underline{\mathbf{Z}}} :
f_!\underline{\mathbf{Z}}
\longrightarrow
\underline{\mathbf{Z}}
$$
，使其在茎上具有预先指定的行为。由定义 \ref{more-etale-definition-f-shriek-lqf}，有
$f_!\underline{\mathbf{Z}} = (f_{p!}\underline{\mathbf{Z}})^\#$
，其中 $f_{p!}\underline{\mathbf{Z}}$ 是第 \ref{more-etale-section-finite-support} 节
构造的预层。因此只须构造 $Y_\etale$ 上预层的映射
$$
f_{p!}\underline{\mathbf{Z}} \longrightarrow \underline{\mathbf{Z}}
$$
。设 $V$ 为 $Y_\etale$ 的对象。由第 \ref{more-etale-section-finite-support} 节回忆
$$
f_{p!}\underline{\mathbf{Z}}(V) = \colim_Z H_Z(\underline{\mathbf{Z}})
$$
。这里余极限取遍所有在 $V$ 上有限的局部闭子概形 $Z \subset X_V$
所成的（偏序）集合。对每个这样的 $Z$，定义映射
$$
H_Z(\underline{\mathbf{Z}}) \longrightarrow \underline{\mathbf{Z}}(V)
$$
，使其与定义余极限的映射相容。

\medskip\noindent
设 $Z \subset X_V$ 局部闭且在 $V$ 上有限。取开子集
$U \subset X_V$，使它把 $Z$ 包含为闭子集。$H_Z(\underline{\mathbf{Z}})$
的元素 $s$ 是截面 $s \in \underline{\mathbf{Z}}(U)$，其支撑包含在 $Z$ 中。
令 $U_n \subset U$ 为 $s$ 取值 $n \in \mathbf{Z}$ 的开闭子集。由支撑条件，
当 $n \not = 0$ 时有 $Z \cap U_n = U_n$。因此，对 $n \not = 0$，开集
$U_n$ 在 $Z$ 中也闭（作为其余各项之并的补集）；又因 $Z$ 在 $V$ 上有限，
可知 $U_n \to V$ 有限。由权函数的定义，这意味着函数
$\int_{U_n \to V} w|_{U_n}$ 在 $V$ 上局部常值，故可把它看作
$\underline{\mathbf{Z}}(V)$ 的元素。我们的构造把 $(Z, s)$ 映到元素
$$
\sum\nolimits_{n \in \mathbf{Z},\ n \not = 0}
n \left(\int_{U_n \to V} w|_{U_n}\right)
\quad \in \quad \underline{\mathbf{Z}}(V)
$$
。该和在 $V$ 上局部有限，因而确有意义；略去细节（在整个讨论中，读者可先取
仿射开集，并保证论证中出现的所有概形都拟紧，从而使该和为有限和）。略去如下
验证：此构造与余极限中的映射以及定义
$f_{p!}\underline{\mathbf{Z}}$ 的限制映射相容。

\medskip\noindent
设 $\overline{y}$ 为 $Y$ 的几何点，位于点 $y \in Y$ 上方。在
$\overline{y}$ 处取茎，上述构造确定映射
$$
(f_!\underline{\mathbf{Z}})_{\overline{y}}
= \bigoplus\nolimits_{f(\overline{x}) = \overline{y}} \mathbf{Z}
\longrightarrow
\mathbf{Z} = \underline{\mathbf{Z}}_{\overline{y}}
$$
。为完成证明，将说明在对应于 $\overline{x}$ 的直和项上，此映射由乘以
$w(\overline{x})$ 给出。具体地，取位于 $\overline{y}$ 上方的
$\overline{x}$。可找到 \'etale 邻域
$(V, \overline{v}) \to (Y, \overline{y})$，使 $X_V$ 含有在 $V$ 上有限的
开子概形 $U$，并且 $U$ 只包含几何点 $\overline{x}$，不包含 $X$ 中位于
$\overline{y}$ 上方的其他几何点。这由“关于态射的进一步结果”引理
\ref{more-morphisms-lemma-etale-makes-quasi-finite-finite-multiple-points-var};
得出；略去一些细节。于是 $(U, 1)$ 定义了 $f_!\underline{\mathbf{Z}}$ 在
$V$ 上的截面；它在对应于 $\overline{x}$ 的直和项中映到 $1$，在其他直和项中
映到零（见引理 \ref{more-etale-lemma-finite-support-stalk} 的证明）。上述构造把
$(U, 1)$ 映到 $\int_{U \to V} w|_U$，后者在 $\overline{v}$ 的某个邻域上
恒取值 $w(\overline{x})$，正合所需。
\end{proof}

\begin{lemma}
\label{more-etale-lemma-properties-trace-map}
设 $f : X \to Y$ 为局部拟有限的概形态射，
$w : X \to \mathbf{Z}$ 为 $f$ 的权函数。上述构造的迹映射具有以下性质：
\begin{enumerate}
\item $\text{Tr}_{f, w, \mathcal{F}}$ 关于 $\mathcal{F}$ 具有函子性；
\item $\text{Tr}_{f, w, \mathcal{F}}$ 与任意基变换相容；
\item 给定环 $\Lambda$ 及 $D(Y_\etale, \Lambda)$ 中的 $K$，得到
$\text{Tr}_{f, w, K} : f_!f^{-1}K \to K$；它关于 $K$ 具有函子性，
且与任意基变换相容。
\end{enumerate}
\end{lemma}

\begin{proof}
(1) 可由引理 \ref{more-etale-lemma-trace-map-exists} 的证明中迹映射的构造得到；
也可更简单地注意到，该映射的刻画迫使这一性质在所有茎上成立。设
$$
\xymatrix{
X' \ar[r]_{g'} \ar[d]_{f'} & X \ar[d]^f \\
Y' \ar[r]^g & Y
}
$$
为概形的笛卡尔图。由“关于态射的进一步结果”引理
\ref{more-morphisms-lemma-weighting-base-change}，函数
$w' = w \circ g' : X' \to \mathbf{Z}$ 是 $f'$ 的权函数。陈述 (2) 意味着图
$$
\xymatrix{
g^{-1}f_!f^{-1}\mathcal{F}
\ar[rr]_-{g^{-1}\text{Tr}_{f, w, \mathcal{F}}} \ar@{=}[d] & &
g^{-1}\mathcal{F} \ar@{=}[d] \\
f'_!(f')^{-1}g^{-1}\mathcal{F}
\ar[rr]^-{\text{Tr}_{f', w', g^{-1}\mathcal{F}}} & &
g^{-1}\mathcal{F}
}
$$
交换，其中左侧竖等号由
$$
g^{-1}f_!f^{-1}\mathcal{F} =
f'_!(g')^{-1}f^{-1}\mathcal{F} =
f'_!(f')^{-1}g^{-1}\mathcal{F}
$$
给出；其中第一个等号来自引理 \ref{more-etale-lemma-lqf-base-change-f-shriek}
（下叹号的基变换）。该图的交换性来自迹映射在茎上作用的刻画，以及引理
\ref{more-etale-lemma-lqf-base-change-f-shriek} 的基变换映射保持茎的描述这一事实。

\medskip\noindent
由 (1) 与 (2)，(3) 随即得到，因为函子
$f^{-1} : D(Y_\etale, \Lambda) \to D(X_\etale, \Lambda)$ 与
$f_! : D(X_\etale, \Lambda) \to D(Y_\etale, \Lambda)$
是通过对表示所论对象的任意模复形逐项应用 $f^{-1}$ 与 $f_!$ 得到的。
\end{proof}

\begin{lemma}
\label{more-etale-lemma-trace-composition}
设 $f : X \to Y$ 与 $g : Y \to Z$ 为局部拟有限态射。设
$w_f : X \to \mathbf{Z}$ 为 $f$ 的权函数，
$w_g : Y \to \mathbf{Z}$ 为 $g$ 的权函数。对
$K \in D(Z_\etale, \Lambda)$，复合
$$
(g \circ f)_!(g \circ f)^{-1}K =
g_! f_! f^{-1} g^{-1}K
\xrightarrow{g_! \text{Tr}_{f, w_f, g^{-1}K}}
g_!g^{-1}K
\xrightarrow{\text{Tr}_{g, w_g, K}}
K
$$
等于 $\text{Tr}_{g \circ f, w_{g \circ f}, K}$，其中
$w_{g \circ f}(x) = w_f(x) w_g(f(x))$。
\end{lemma}

\begin{proof}
由引理 \ref{more-etale-lemma-lqf-shriek-composition}，有
$(g \circ f)_! = g_! \circ f_!$。在“关于态射的进一步结果”引理
\ref{more-morphisms-lemma-weighting-composition} 中已知 $w_{g \circ f}$
是 $g \circ f$ 的权函数，故该陈述确有意义。在茎上计算即可检验等式。
略去细节。
\end{proof}

\begin{example}[平坦拟有限态射的迹]
\label{more-etale-example-trace-for-flat-quasi-finite}
设 $f : X \to Y$ 为平坦、局部拟有限且局部有限表示的概形态射。
令
$$
w(x) = \text{length}_{\mathcal{O}_{X, x}}
(\mathcal{O}_{X, x}/\mathfrak m_{f(x)} \mathcal{O}_{X, x})
[\kappa(x) : \kappa(f(x))]_i
$$
，便得到典范正权函数 $w : X \to \mathbf{Z}$；见“关于态射的进一步结果”引理
\ref{more-morphisms-lemma-weighting-flat-quasi-finite}。因此，由引理
\ref{more-etale-lemma-trace-map-exists} 与 \ref{more-etale-lemma-properties-trace-map}，对 $f$ 得到迹映射
$$
\text{Tr}_{f, K} : f_!f^{-1}K \longrightarrow K
$$
；它关于 $D(Y_\etale, \Lambda)$ 中的 $K$ 具有函子性，且与任意基变换相容。
注意，$f$ 的任意基变换 $f' : X' \to Y'$ 满足相同性质，而且 $w$
限制为 $f'$ 的典范权函数。
\end{example}

\begin{remark}
\label{more-etale-remark-trace-etale-counit}
设 $j : U \to X$ 为概形的 \'etale 态射。则例
\ref{more-etale-example-trace-for-flat-quasi-finite} 的迹映射
$\text{Tr} : j_!j^{-1}K \to K$ 等于 $j_!$ 与 $j^{-1}$ 之间伴随的余单位。
在“\'Etale 上同调”第 \ref{etale-cohomology-section-trace-method} 节中，
我们已对这个余单位使用术语“迹”。
\end{remark}

\begin{example}[正规基底上拟有限态射的迹]
\label{more-etale-example-trace-for-quasi-finite-over-normal}
设 $Y$ 为几何单枝且局部 Noether 的概形；例如 $Y$ 可以是正规簇。
设 $f : X \to Y$ 为局部拟有限的概形态射。则存在 $f$ 的正权函数
$w : X \to \mathbf{Z}$；粗略地说，它把 $x$ 映到
$\mathcal{O}_{X, x}^{sh}$ 在 $\mathcal{O}_{Y, f(x)}^{sh}$ 上的“泛可分次数”。
见“关于态射的进一步结果”引理
\ref{more-morphisms-lemma-weighting-quasi-finite-Noetherian}。因此，由引理
\ref{more-etale-lemma-trace-map-exists} 与 \ref{more-etale-lemma-properties-trace-map}，对 $f$ 和 $w$
得到迹映射
$$
\text{Tr}_{f, w, K} : f_!f^{-1}K \longrightarrow K
$$
；它关于 $D(Y_\etale, \Lambda)$ 中的 $K$ 具有函子性，且与任意基变换相容。
不过，在此情形下，给定 $f$ 的基变换 $f' : X' \to Y'$，$w$ 在 $X'$
上的限制一般没有关于态射 $f'$ 的“自然”解释。
\end{example}
















\section{局部拟有限态射的上叹号}
\label{more-etale-section-duality-locally-quasi-finite}

\noindent
对概形的局部拟有限态射 $f : X \to Y$，函子
$f_! : \textit{Ab}(X_\etale) \to \textit{Ab}(Y_\etale)$ 与直和可交换且正合；
见引理 \ref{more-etale-lemma-lqf-f-shriek-stalk}。这表明它应有一个右伴随，记为 $f^!$。

\medskip\noindent
警告：这个函子是非导出版本！

\begin{lemma}
\label{more-etale-lemma-lqf-f-upper-shriek}
设 $f : X \to Y$ 为局部拟有限的概形态射。
\begin{enumerate}
\item 函子 $f_! : \textit{Ab}(X_\etale) \to \textit{Ab}(Y_\etale)$
有右伴随 $f^! : \textit{Ab}(Y_\etale) \to \textit{Ab}(X_\etale)$。
\item 有
$f^!(\overline{y}_*A) = \prod_{f(\overline{x}) = \overline{y}} \overline{x}_*A$.
\item 若 $\Lambda$ 为环，则函子
$f_! : \textit{Mod}(X_\etale, \Lambda) \to \textit{Mod}(Y_\etale, \Lambda)$
有右伴随
$f^! : \textit{Mod}(Y_\etale, \Lambda) \to \textit{Mod}(X_\etale, \Lambda)$
，且在底阿贝尔层上它与 $f^!$ 一致。
\end{enumerate}
\end{lemma}

\begin{proof}
证明 (1)。令 $E \subset \Ob(\textit{Ab}(Y_\etale))$ 为摩天层之积所成的类。
我们断言：
\begin{enumerate}
\item[(a)] $\textit{Ab}(Y_\etale)$ 中的每个 $\mathcal{G}$ 都是 $E$
中某个元素的子层；
\item[(b)] 对每个 $\mathcal{G} \in E$，存在
$\textit{Ab}(X_\etale)$ 的对象 $\mathcal{H}$，使得
$\Hom(f_!\mathcal{F}, \mathcal{G}) = \Hom(\mathcal{F}, \mathcal{H})$
关于 $\mathcal{F}$ 函子性地成立。
\end{enumerate}
一旦验证此断言，“同调”引理 \ref{homology-lemma-partially-defined-adjoint}
的对偶便给出伴随函子 $f^!$。

\medskip\noindent
(a) 成立，因为可把 $\mathcal{G}$ 映到层
$\prod \overline{y}_*\mathcal{G}_{\overline{y}}$，其中积取遍 $Y$ 的所有几何点。
由“\'Etale 上同调”定理 \ref{etale-cohomology-theorem-exactness-stalks}，
这是单射。（在拓扑空间上的阿贝尔层范畴中，这正是 Godement 分解的第一步。）

\medskip\noindent
如下可见 (b) 与引理的 (2)。假设对某些阿贝尔群 $A_{\overline{y}}$，有
$\mathcal{G} = \prod \overline{y}_*A_{\overline{y}}$。则
$$
\Hom(f_!\mathcal{F}, \mathcal{G}) =
\prod \Hom(f_!\mathcal{F}, \overline{y}_*A_{\overline{y}})
$$
。因此只须找到 $X_\etale$ 上的阿贝尔层 $\mathcal{H}_{\overline{y}}$，
它们表示函子
$\mathcal{F} \mapsto \Hom(f_!\mathcal{F}, \overline{y}_*A_{\overline{y}})$
，再取 $\mathcal{H} = \prod \mathcal{H}_{\overline{y}}$。这把问题化归到
$\mathcal{H} = \overline{y}_*A$ 的情形，其中
$\overline{y} : \Spec(k) \to Y$ 为某个固定几何点，$A$ 为某个固定阿贝尔群。
我们断言，此时取
$\mathcal{H} = \prod_{f(\overline{x}) = \overline{y}} \overline{x}_*A$ 即可。
这将完成引理 (1) 与 (2) 的证明。事实上，由引理
\ref{more-etale-lemma-lqf-f-shriek-stalk} 对茎的描述，一方面有
$$
\Hom(f_!\mathcal{F}, \overline{y}_*A) =
\Hom_{\textit{Ab}}((f_!\mathcal{F})_{\overline{y}}, A) =
\Hom_{\textit{Ab}}(\bigoplus\nolimits_{f(\overline{x}) = \overline{y}}
\mathcal{F}_{\overline{x}}, A)
$$
，另一方面有
$$
\Hom(\mathcal{F}, \mathcal{H}) =
\prod\nolimits_{f(\overline{x}) = \overline{y}}
\Hom(\mathcal{F}, \overline{x}_*A) =
\prod\nolimits_{f(\overline{x}) = \overline{y}}
\Hom_{\textit{Ab}}(\mathcal{F}_{\overline{x}}, A)
$$
。留给读者验证它们作为 $\mathcal{F}$ 的函子彼此等同。

\medskip\noindent
证明 (3)。注意，$\textit{Mod}(X_\etale, \Lambda)$ 的对象等同于
$\textit{Ab}(X_\etale)$ 的对象 $\mathcal{F}$ 连同映射
$\Lambda \to \text{End}(\mathcal{F})$。因此，(1) 中的函子 $f_!$ 与 $f^!$
定义了 (3) 中的函子 $f_!$ 与 $f^!$。直接计算可知它们互为伴随。
\end{proof}

\begin{lemma}
\label{more-etale-lemma-etale-upper-shriek}
设 $j : U \to X$ 为 \'etale 态射。则 $j^! = j^{-1}$。
\end{lemma}

\begin{proof}
这是因为第 \ref{more-etale-section-finite-support} 节定义的 $j_!$ 与“\'Etale 上同调”
第 \ref{etale-cohomology-section-extension-by-zero} 节定义的 $j_!$ 一致；
见引理 \ref{more-etale-lemma-finite-support-etale-shriek}。最后，在“\'Etale 上同调”
第 \ref{etale-cohomology-section-extension-by-zero} 节中，函子 $j_!$
定义为 $j^{-1}$ 的左伴随，故由伴随函子的唯一性得出结论。
\end{proof}

\begin{lemma}
\label{more-etale-lemma-upper-shriek-restriction}
设 $f : X \to Y$ 与 $g : Y \to Z$ 为分离且局部拟有限的态射。
存在典范同构 $(g \circ f)^! \to f^! \circ g^!$。给定第三个局部拟有限态射
$h : Z \to T$，图
$$
\xymatrix{
(h \circ g \circ f)^! \ar[r] \ar[d] &
f^! \circ (h \circ g)^! \ar[d] \\
(g \circ f)^! \circ h^! \ar[r] & f^! \circ g^! \circ h^!
}
$$
交换。
\end{lemma}

\begin{proof}
由伴随函子的唯一性，这立即转化为函子 $f_!$ 的相应（对偶）陈述。
见引理 \ref{more-etale-lemma-lqf-shriek-composition}。
\end{proof}

\begin{lemma}
\label{more-etale-lemma-upper-shriek-restriction-etale}
设 $j : U \to X$ 与 $j' : V \to U$ 为 \'etale 态射。同构
$(j \circ j')^{-1} = (j')^{-1} \circ j^{-1}$ 与引理
\ref{more-etale-lemma-upper-shriek-restriction} 的同构
$(j \circ j')^! = (j')^! \circ j^!$，通过引理
\ref{more-etale-lemma-etale-upper-shriek} 的同构彼此一致。
\end{lemma}

\begin{proof}
略。
\end{proof}

\begin{lemma}
\label{more-etale-lemma-lqf-base-change-upper-shriek}
考虑概形的笛卡尔方块
$$
\xymatrix{
X' \ar[r]_{g'} \ar[d]_{f'} & X \ar[d]^f \\
Y' \ar[r]^g & Y
}
$$
，其中 $f$ 局部拟有限。对 $Y'_\etale$ 上任意阿贝尔层 $\mathcal{F}$，有
$(g')_*(f')^!\mathcal{F} = f^!g_*\mathcal{F}$。
\end{lemma}

\begin{proof}
由伴随函子的唯一性，这由函子 $f_!$ 的相应（对偶）陈述得出。
见引理 \ref{more-etale-lemma-lqf-base-change-f-shriek}。
\end{proof}

\begin{remark}
\label{more-etale-remark-pointed-sets}
本节内容可推广到带基点集合的层。具体地，对位点 $\mathcal{C}$，以
$\Sh^*(\mathcal{C})$ 表示带基点集合的层之范畴。本节及前一节的构造经相应修改后
适用于带基点集合的层。因此，给定概形的局部拟有限态射 $f : X \to Y$，
得到一对伴随函子
$$
f_! : \Sh^*(X_\etale) \longrightarrow \Sh^*(Y_\etale)
\quad\text{且}\quad
f^! : \Sh^*(Y_\etale) \longrightarrow \Sh^*(X_\etale)
$$
，使得对 $Y$ 的每个几何点 $\overline{y}$，都有同构
$$
(f_!\mathcal{F})_{\overline{y}} =
\coprod\nolimits_{f(\overline{x}) = \overline{y}}
\mathcal{F}_{\overline{x}}
$$
（余积取于带基点集合的范畴中），它关于
$\mathcal{F} \in \Sh^*(X_\etale)$ 具有函子性；还有同构
$$
f^!(\overline{y}_*S) =
\prod\nolimits_{f(\overline{x}) = \overline{y}}
\overline{x}_*S
$$
，它关于带基点集合 $S$ 具有函子性。若
$F : \textit{Ab}(X_\etale) \to \Sh^*(X_\etale)$ 与
$F : \textit{Ab}(Y_\etale) \to \Sh^*(Y_\etale)$
表示遗忘函子，则各构造之间的相容性保证存在典范映射
$$
f_!F(\mathcal{F}) \longrightarrow F(f_!\mathcal{F})
$$
，它关于 $\mathcal{F} \in \textit{Ab}(X_\etale)$ 具有函子性；以及
$$
F(f^!\mathcal{G}) \longrightarrow f^!F(\mathcal{G})
$$
，它关于 $\mathcal{G} \in \textit{Ab}(Y_\etale)$ 具有函子性；二者分别在茎上、
在摩天层上给出显然的映射。事实上，变换
$F \circ f^! \to f^! \circ F$ 是同构（因为 $f^!$ 与积可交换）。
\end{remark}







\section{局部拟有限态射的导出上叹号}
\label{more-etale-section-derived-duality-locally-quasi-finite}

\noindent
取第 \ref{more-etale-section-duality-locally-quasi-finite} 节中各函子的导出版本，
可得如下结论。

\begin{lemma}
\label{more-etale-lemma-lqf-shriek-derived}
设 $f : X \to Y$ 是概形的局部拟有限态射，$\Lambda$ 是环。
定义 \ref{more-etale-definition-f-shriek-lqf} 与引理
\ref{more-etale-lemma-lqf-f-upper-shriek} 中的函子 $f_!$ 和 $f^!$
在导出范畴上诱导出一对伴随函子
$f_! : D(X_\etale, \Lambda) \to D(Y_\etale, \Lambda)$ 与
$Rf^! : D(Y_\etale, \Lambda) \to D(X_\etale, \Lambda)$。
\end{lemma}

\noindent
在分离情形，函子 $f_!$ 定义于第 \ref{more-etale-section-compact-support} 节。

\begin{proof}
这立即由“导出范畴”引理
\ref{derived-lemma-derived-adjoint-functors}、$f_!$ 的正合性
（引理 \ref{more-etale-lemma-lqf-f-shriek-stalk}，因而 $Lf_! = f_!$）以及如下事实得出：
$Y_\etale$ 上有足够多的 $\Lambda$-模 K-内射复形，故 $Rf^!$ 有定义。
\end{proof}

\begin{remark}
\label{more-etale-remark-lqf-shriek-derived-torsion}
设 $f : X \to Y$ 是概形的局部拟有限态射，$\Lambda$ 是环。
引理 \ref{more-etale-lemma-lqf-shriek-derived} 中的函子
$f_! : D(X_\etale, \Lambda) \to D(Y_\etale, \Lambda)$
把上同调层为挠层的复形送到上同调层为挠层的复形。
由 $f_!$ 的茎的描述，此事立即成立；见引理
\ref{more-etale-lemma-lqf-f-shriek-stalk}。
\end{remark}

\begin{lemma}
\label{more-etale-lemma-mayer-vietoris-shriek}
设 $X$ 是概形，$X = U \cup V$，其中 $U$ 与 $V$ 均为开集。
设 $\Lambda$ 是环且 $K \in D(X_\etale, \Lambda)$。则有特异三角
$$
j_{U \cap V!}K|_{U \cap V} \to
j_{U!}K|_U \oplus j_{V!}K|_V \to K \to 
j_{U \cap V!}K|_{U \cap V}[1]
$$
成立于 $D(X_\etale, \Lambda)$ 中，其中记号含义显然。
\end{lemma}

\begin{proof}
由于这里所用的限制函子与下叹号函子均正合，只需证明：对
$X_\etale$ 上任意阿贝尔层 $\mathcal{F}$，序列
$$
0 \to j_{U \cap V!}\mathcal{F}|_{U \cap V} \to
j_{U!}\mathcal{F}|_U \oplus j_{V!}\mathcal{F}|_V \to \mathcal{F} \to 0
$$
正合。这可在茎上验证。
\end{proof}

\begin{lemma}
\label{more-etale-lemma-triangle-associated-to-open}
设 $X$ 是概形，$Z \subset X$ 是闭子概形，$U \subset X$ 是其补集。
以 $i : Z \to X$ 和 $j : U \to X$ 表示嵌入态射。设 $\Lambda$ 是环且
$K \in D(X_\etale, \Lambda)$。则有特异三角
$$
j_!j^{-1}K \to K \to i_*i^{-1}K \to j_!j^{-1}K[1]
$$
成立于 $D(X_\etale, \Lambda)$ 中。
\end{lemma}

\begin{proof}
这立即由“\'Etale 上同调”引理
\ref{etale-cohomology-lemma-ses-associated-to-open} 以及如下事实得出：
函子 $j_!$、$j^{-1}$、$i_*$、$i^{-1}$ 均正合，故其导出版本可通过将这些函子
作用于任意一个表示 $K$ 的层复形来计算。
\end{proof}





\section{通过紧化构造导出下叹号的预备结果}
\label{more-etale-section-prelim}

\noindent
本节证明若干关于函子之间某些自然同构之存在性的引理；这些同构
立即来自固有基变换。

\begin{lemma}
\label{more-etale-lemma-shriek-proper-and-open}
考虑概形的交换图
$$
\xymatrix{
X' \ar[r]_{g'} \ar[d]_{f'} & X \ar[d]^f \\
Y' \ar[r]^g & Y
}
$$
其中 $f$ 与 $f'$ 固有，而 $g$ 与 $g'$ 分离且局部拟有限。
设 $\Lambda$ 是环。关于 $K \in D(X'_\etale, \Lambda)$ 函子性地，
在 $D(Y_\etale, \Lambda)$ 中有典范映射
$$
g_!Rf'_*K \longrightarrow Rf_*(g'_!K)
$$
。在以下任一情形，该映射是同构：
(a) $K$ 下有界且其上同调层为挠层；或
(b) $\Lambda$ 是挠环。
\end{lemma}

\begin{proof}
以 $X'_\etale$ 上的 $\Lambda$-模层 K-内射复形
$\mathcal{J}^\bullet$ 表示 $K$。选取到 $X_\etale$ 上的
$\Lambda$-模层 K-内射复形 $\mathcal{I}^\bullet$ 的拟同构
$g'_!\mathcal{J}^\bullet \to \mathcal{I}^\bullet$。于是可考虑映射
$$
g_!f'_*\mathcal{J}^\bullet =
g_!f'_!\mathcal{J}^\bullet =
f_!g'_!\mathcal{J}^\bullet =
f_*g'_!\mathcal{J}^\bullet \to
f_*\mathcal{I}^\bullet
$$
其中第一与第三个等号来自引理 \ref{more-etale-lemma-proper-f-shriek}，
第二个等号来自引理 \ref{more-etale-lemma-f-shriek-composition}；该引理告诉我们，
$g_! \circ f'_!$ 与 $f_! \circ g'_!$ 作为
$(g \circ f')_* = (f \circ g')_*$ 的子层，均等于
$(g \circ f')_! = (f \circ g')_!$。

\medskip\noindent
设 $\Lambda$ 为挠环，即处于情形 (b)。沿用上述记号，只需证明
$f_*g'_!\mathcal{J}^\bullet \to f_*\mathcal{I}^\bullet$ 是同构。
该问题在 $Y$ 上是局部的，故可设 $f$ 的纤维维数有界；见“态射”引理
\ref{morphisms-lemma-morphism-finite-type-bounded-dimension}。
于是 $Rf_*$ 的上同调维数有限；见“\'Etale 上同调”引理
\ref{etale-cohomology-lemma-cohomological-dimension-proper}。
因此由“导出范畴”引理 \ref{derived-lemma-unbounded-right-derived}，
若能证明对 $q > 0$ 以及 $X'_\etale$ 上任意内射 $\Lambda$-模层
$\mathcal{J}$，均有 $R^qf_*(g'_!\mathcal{J}) = 0$，结论即得。

\medskip\noindent
由“\'Etale 上同调”引理
\ref{etale-cohomology-lemma-proper-base-change-stalk}，
$R^qf_*(g'_!\mathcal{J})$ 在几何点 $\overline{y}$ 处的茎等于
$H^q(X_{\overline{y}}, (g'_!\mathcal{J})|_{X_{\overline{y}}})$。
由于 $g'_!$ 的构造与基变换交换
（引理 \ref{more-etale-lemma-base-change-f-shriek-separated}），它等于
$$
H^q(X_{\overline{y}}, g'_{\overline{y}, !}(\mathcal{J}|_{X'_{\overline{y}}}))
$$
，其中 $g'_{\overline{y}} : X'_{\overline{y}} \to X_{\overline{y}}$
是几何纤维之间所诱导的态射。由于 $Y' \to Y$ 局部拟有限，
$X'_{\overline{y}}$ 是各纤维 $X'_{\overline{y}'}$ 的不交并；这里
$\overline{y}'$ 遍历 $Y'$ 中位于 $\overline{y}$ 上方的几何点。
记 $g'_{\overline{y}'} : X'_{\overline{y}'} \to X_{\overline{y}}$
为 $g'_{\overline{y}}$ 在 $X'_{\overline{y}'}$ 上的限制。于是前述上同调群等于
$$
H^q(X_{\overline{y}},
\bigoplus\nolimits_{\overline{y}'/\overline{y}}
g'_{\overline{y}', !}(\mathcal{J}|_{X'_{\overline{y}'}}))
$$
；这例如可由引理 \ref{more-etale-lemma-colim-f-shriek-separated} 得到（但从第
\ref{more-etale-section-compact-support} 节中 $g'_{\overline{y}, !}$ 的定义也显然可见）。
由于在 $X_{\overline{y}}$ 上取 \'etale 上同调与直和交换
（“\'Etale 上同调”定理 \ref{etale-cohomology-theorem-colimit}），
我们归结为证明
$$
H^q(X_{\overline{y}}, g'_{\overline{y}', !}(\mathcal{J}|_{X'_{\overline{y}'}}))
$$
为零。注意，$g_{\overline{y}'} : X'_{\overline{y}'} \to X_{\overline{y}}$
是 $\overline{y}$ 上固有概形之间的态射，故其本身固有。它又局部拟有限，
因而 $g_{\overline{y}'}$ 有限。于是
$g'_{\overline{y}', !} = g'_{\overline{y}', *} = Rg'_{\overline{y}', *}$。
由 Leray，所需证明化为
$$
H^q(X'_{\overline{y}'}, \mathcal{J}|_{X'_{\overline{y}'}})
$$
为零。由于 $\Lambda$ 为挠环，这由固有基变换（“\'Etale 上同调”引理
\ref{etale-cohomology-lemma-proper-base-change-stalk}）得出，因为
$\mathcal{J}$ 在 $f'$ 下的高阶直像为零。

\medskip\noindent
现证情形 (a)。我们用标准论证从情形 (b) 推出它。我们将证明，对所有
$p \in \mathbf{Z}$，所诱导的映射
$g_! R^pf'_* K \to R^pf_*(g'_!K)$ 均为同构。固定整数
$p_0 \in \mathbf{Z}$，并取整数 $a$，使得当 $j < a$ 时 $H^j(K) = 0$。
我们对 $a$ 作降序归纳，证明当 $p \leq p_0$ 时
$g_! R^pf'_* K \to R^pf_*(g'_!K)$ 为同构。若 $a > p_0$，则由平凡消没，
上述映射的左右两端在 $p \leq p_0$ 时均为零；见“导出范畴”引理
\ref{derived-lemma-negative-vanishing}（并使用 $g_!$ 与 $g'_!$ 均为正合函子）。
设 $a \leq p_0$。考虑特异三角
$$
H^a(K)[-a] \to K \to \tau_{\geq a + 1}K
$$
由归纳假设，结论对 $\tau_{\geq a + 1}K$ 成立。下一段将证明结论对
$H^a(K)[-a]$ 成立。随后，对如下特异三角映射所伴的上同调层长正合列之间的
映射应用五引理
$$
\xymatrix{
g_! Rf'_*(H^a(K)[-a]) \ar[d] \ar[r] &
g_! Rf'_* K \ar[r] \ar[d] &
g_! Rf'_* \tau_{\geq a + 1} K \ar[d] \\
Rf_*(g'_!(H^a(K)[-a])) \ar[r] &
Rf_*(g'_!K) \ar[r] &
Rf_*(g'_!\tau_{|geq a + 1}K)
}
$$
，即得 $K$ 的结论。略去一些细节。

\medskip\noindent
设 $\mathcal{F}$ 是 $X'_\etale$ 上的挠阿贝尔层。为完成证明，我们证明对所有
$p$，映射 $g_! Rf'_*\mathcal{F} \to R^pf_*(g'_!\mathcal{F})$ 均为同构。
可写成 $\mathcal{F} = \bigcup \mathcal{F}[n]$，其中
$\mathcal{F}[n] = \Ker(n : \mathcal{F} \to \mathcal{F})$。
由情形 (b)，同构对 $\mathcal{F}[n]$ 成立。函子 $g_!$、$g'_!$、
$R^pf_*$、$R^pf'_*$ 均与滤过余极限交换（这由引理
\ref{more-etale-lemma-lqf-f-shriek-separated-colimits} 和“\'Etale 上同调”引理
\ref{etale-cohomology-lemma-relative-colimit-general} 得出），故证明完成。
\end{proof}

\begin{lemma}
\label{more-etale-lemma-shriek-proper-and-open-compose}
考虑概形的交换图
$$
\xymatrix{
X' \ar[r]_k \ar[d]_{f'} & X \ar[d]^f \\
Y' \ar[r]_l \ar[d]_{g'} & Y \ar[d]^g \\
Z' \ar[r]^m & Z
}
$$
其中 $f$、$f'$、$g$、$g'$ 固有，而 $k$、$l$、$m$ 分离且局部拟有限。
则引理 \ref{more-etale-lemma-shriek-proper-and-open} 对两个方块给出的同构复合后，
给出外侧矩形的同构（精确陈述见证明）。
\end{lemma}

\begin{proof}
该陈述的含义如下。若写成
$R(g \circ f)_* = Rg_* \circ Rf_*$ 与
$R(g' \circ f')_* = Rg'_* \circ Rf'_*$,
，则外侧矩形的同构
$m_! \circ Rg'_* \circ Rf'_* \to Rg_* \circ Rf_* \circ k_!$
等于两个方块的映射之复合
$$
m_! \circ Rg'_* \circ Rf'_* \to
Rg_* \circ l_! \circ Rf'_* \to
Rg_* \circ Rf_* \circ k_!
$$
。为证明此事，选取 $X'_\etale$ 上的 $\Lambda$-模 K-内射复形
$\mathcal{J}^\bullet$，以及到 $X_\etale$ 上的 $\Lambda$-模 K-内射复形
$\mathcal{I}^\bullet$ 的拟同构
$k_!\mathcal{J}^\bullet \to \mathcal{I}^\bullet$。
引理 \ref{more-etale-lemma-shriek-proper-and-open} 的证明表明，典范映射
$$
a : l_!f'_*\mathcal{J}^\bullet \to f_*\mathcal{I}^\bullet
$$
是拟同构，对它应用 $Rg_*$ 即产生第二个箭头。由“位点上同调”引理
\ref{sites-cohomology-lemma-K-injective-flat}，复形
$f_*\mathcal{I}^\bullet$（相应地，$f'_*\mathcal{J}^\bullet$）是
$Y_\etale$（相应地，$Y'_\etale$）上的 $\Lambda$-模 K-内射复形。
（使用这一点算是取巧，本可避免。）特别地，同一论证给出典范映射
$$
b : m_!g'_*f'_*\mathcal{J}^\bullet \to g_*f_*\mathcal{I}^\bullet
$$
也是拟同构，且该拟同构表示第一个箭头。最后，由于
$f'_*\mathcal{J}^\bullet$ 是 K-内射的，引理
\ref{more-etale-lemma-shriek-proper-and-open} 的证明表明
$g_*l_!f'_!\mathcal{J}^\bullet$ 表示
$Rg_*(l_!f'_*\mathcal{J}^\bullet)$。因此 $Rg_*(a) = g_*(a)$，且复合
$g_*(a) \circ b$ 正是引理 \ref{more-etale-lemma-shriek-proper-and-open}
对该矩形给出的箭头。
\end{proof}

\begin{lemma}
\label{more-etale-lemma-shriek-proper-and-open-compose-horizontal}
考虑概形的交换图
$$
\xymatrix{
X'' \ar[r]_{g'} \ar[d]_{f''} &
X' \ar[r]_g \ar[d]_{f'} &
X \ar[d]^f \\
Y'' \ar[r]^{h'} &
Y' \ar[r]^h &
Y
}
$$
其中 $f$、$f'$、$f''$ 固有，而 $g$、$g'$、$h$、$h'$ 分离且局部拟有限。
则引理 \ref{more-etale-lemma-shriek-proper-and-open} 对两个方块给出的同构复合后，
给出外侧矩形的同构（精确陈述见证明）。
\end{lemma}

\begin{proof}
该陈述的含义如下。若用引理 \ref{more-etale-lemma-f-shriek-composition} 的等式写成
$(h \circ h')_! = h_! \circ h'_!$ 与
$(g \circ g')_! = g_! \circ g'_!$
，则外侧矩形的同构
$h_! \circ h'_! \circ Rf''_* \to Rf_* \circ g_! \circ g'_!$
等于两个方块的映射之复合
$$
h_! \circ h'_! \circ Rf''_* \to
h_! \circ Rf'_* \circ g'_! \to
Rf_* \circ g_! \circ g'_!
$$
。为证明此事，选取 $X''_\etale$ 上的 $\Lambda$-模 K-内射复形
$\mathcal{I}^\bullet$，以及到 $X'_\etale$ 上的 $\Lambda$-模 K-内射复形
$\mathcal{J}^\bullet$ 的拟同构
$g'_!\mathcal{I}^\bullet \to \mathcal{J}^\bullet$。继而选取到
$X_\etale$ 上的 $\Lambda$-模 K-内射复形 $\mathcal{K}^\bullet$ 的拟同构
$g_!\mathcal{J}^\bullet \to \mathcal{K}^\bullet$。
引理 \ref{more-etale-lemma-shriek-proper-and-open} 的证明表明，典范映射
$$
h'_!f''_*\mathcal{I}^\bullet \to f'_*\mathcal{J}^\bullet
\quad\text{且}\quad
h_!f'_*\mathcal{J}^\bullet \to f_*\mathcal{K}^\bullet
$$
均为拟同构，且分别定义上述第一与第二个箭头。由于 $g_!$ 是正合函子
（引理 \ref{more-etale-lemma-lqf-f-shriek-separated-colimits}），
$g_!g'_!\mathcal{I}^\bullet \to \mathcal{K}^\bullet$ 是拟同构，
因而典范映射
$$
h_!h'_!f''_*\mathcal{I}^\bullet \to f_*\mathcal{K}^\bullet
$$
也是拟同构，并在导出范畴中表示外侧矩形的映射。显然，该映射是另两个映射的
复合，证明完成。
\end{proof}

\begin{remark}
\label{more-etale-remark-going-around}
考虑交换图
$$
\xymatrix{
X'' \ar[r]_{k'} \ar[d]_{f''} & X' \ar[r]_k \ar[d]_{f'} & X \ar[d]^f \\
Y'' \ar[r]^{l'} \ar[d]_{g''} & Y' \ar[r]^l \ar[d]_{g'} & Y \ar[d]^g \\
Z'' \ar[r]^{m'} & Z' \ar[r]^m & Z
}
$$
，其中竖箭头均固有，横箭头均分离且局部拟有限。如下把图中的方块标记为
$A$、$B$、$C$、$D$：
$$
\begin{matrix}
A & B \\
C & D
\end{matrix}
$$
于是，引理 \ref{more-etale-lemma-shriek-proper-and-open} 对各方块给出的映射为
（这里使用 $Rf_* = f_*$ 等）
$$
\begin{matrix}
\gamma_A : l'_! \circ f''_* \to f'_* \circ k'_! &
\gamma_B : l_! \circ f'_* \to f_* \circ k_! \\
\gamma_C : m'_! \circ g''_* \to g'_* \circ l'_! &
\gamma_D : m_! \circ g'_* \to g_* \circ l_!
\end{matrix}
$$
对各 $2 \times 1$ 与 $1 \times 2$ 矩形，另有四个映射
$$
\begin{matrix}
\gamma_{A + B} :
(l \circ l')_! \circ f''_* \to f_* \circ (k \circ k')_*  \\
\gamma_{C + D} :
(m \circ m')_! \circ g''_* \to g_* \circ (l \circ l')_! \\
\gamma_{A + C} :
m'_! \circ (g'' \circ f'')_* \to (g' \circ f')_* \circ k'_! \\
\gamma_{B + D} :
m_! \circ (g' \circ f')_* \to (g \circ f)_* \circ k_!
\end{matrix}
$$
由引理 \ref{more-etale-lemma-shriek-proper-and-open-compose-horizontal}，有
$$
\gamma_{A + B} = \gamma_B \circ \gamma_A, \quad
\gamma_{C + D} = \gamma_D \circ \gamma_C
$$
；由引理 \ref{more-etale-lemma-shriek-proper-and-open-compose}，有
$$
\gamma_{A + C} = \gamma_A \circ \gamma_C, \quad
\gamma_{B + D} = \gamma_B \circ \gamma_D
$$
更准确地说，此处应写成
$\gamma_{A + B} = (\gamma_B \star \text{id}_{k'_!}) \circ
(\text{id}_{l_!} \star \gamma_A)$，其中采用
“范畴”第 \ref{categories-section-formal-cat-cat} 节的记号
；其余情形亦同。综上，我们（先验地）得到两个变换
$$
m_! \circ m'_! \circ g''_* \circ f''_*
\longrightarrow
g_* \circ f_* \circ k_! \circ k'_!
$$
，即
$$
\gamma_B \circ \gamma_D \circ \gamma_A \circ \gamma_C =
\gamma_{B + D} \circ \gamma_{A + C}
$$
以及
$$
\gamma_B \circ \gamma_A \circ \gamma_D \circ \gamma_C =
\gamma_{A + B} \circ \gamma_{C + D}
$$
本评注的要点是指出这两个变换相等。为此，只需证明
$$
\xymatrix{
m_! \circ g'_* \circ l'_! \circ f''_* \ar[r]_{\gamma_D} \ar[d]_{\gamma_A} &
g_* \circ l_! \circ l'_! \circ f''_* \ar[d]^{\gamma_A} \\
m_! \circ g'_* \circ f'_* \circ k'_! \ar[r]^{\gamma_D} &
g_* \circ l_! \circ f'_* \circ k'_!
}
$$
交换。其原因是方块 $A$ 与 $D$ 仅在一点相交；更准确地说，这由“范畴”引理
\ref{categories-lemma-properties-2-cat-cats} 得出，或者更直接地，由“范畴”定义
\ref{categories-definition-horizontal-composition} 之前的讨论得出。
\end{remark}

\begin{lemma}
\label{more-etale-lemma-shriek-proper-and-open-base-change}
设 $b : Y_1 \to Y$ 是概形的态射。考虑概形的交换图
$$
\vcenter{
\xymatrix{
X' \ar[r]_{g'} \ar[d]_{f'} & X \ar[d]^f \\
Y' \ar[r]^g & Y
}
}
\quad\text{且令}\quad
\vcenter{
\xymatrix{
X'_1 \ar[r]_{g'_1} \ar[d]_{f'_1} & X_1 \ar[d]^{f_1} \\
Y'_1 \ar[r]^{g_1} & Y_1
}
}
$$
为沿 $b$ 的基变换。设 $f$ 与 $f'$ 固有，而 $g$ 与 $g'$ 分离且局部拟有限。
对环 $\Lambda$ 以及 $D(X'_\etale, \Lambda)$ 中的 $K$，有交换图
$$
\xymatrix{
b^{-1}g_!Rf'_*K \ar[d] \ar[r] &
g_{1, !}(b')^{-1}Rf'_*K \ar[r] &
g_{1, !}Rf'_{1, *}(a')^{-1}K \ar[d] \\
b^{-1}Rf_*g'_!K \ar[r] &
Rf_{1, *}a^{-1}g'_!K \ar[r] &
Rf_{1, *}g'_{1, !}(a')^{-1}K
}
$$
成立于 $D(Y_{1, \etale}, \Lambda)$ 中；其中
$a : X_1 \to X$、$a' : X'_1 \to X'$、$b' : Y'_1 \to Y'$ 是投影，
竖映射是引理 \ref{more-etale-lemma-shriek-proper-and-open} 的箭头，横箭头是基变换映射
（来自“\'Etale 上同调”第
\ref{etale-cohomology-section-base-change-preliminaries} 节）与引理
\ref{more-etale-lemma-base-change-f-shriek-separated} 的基变换映射。
\end{lemma}

\begin{proof}
以 $X'_\etale$ 上的 $\Lambda$-模层 K-内射复形
$\mathcal{J}^\bullet$ 表示 $K$。选取到 $X_\etale$ 上的
$\Lambda$-模层 K-内射复形 $\mathcal{I}^\bullet$ 的拟同构
$g'_!\mathcal{J}^\bullet \to \mathcal{I}^\bullet$。
引理 \ref{more-etale-lemma-shriek-proper-and-open} 的证明把
$g_!Rf'_*K \to Rf_*g'_!K$ 构造为
$$
g_!f'_*\mathcal{J}^\bullet =
g_!f'_!\mathcal{J}^\bullet =
f_!g'_!\mathcal{J}^\bullet =
f_*g'_!\mathcal{J}^\bullet \to
f_*\mathcal{I}^\bullet
$$
选取到 $X'_{1, \etale}$ 上的 $\Lambda$-模层 K-内射复形
$\mathcal{J}_1^\bullet$ 的拟同构
$(a')^{-1}\mathcal{J}^\bullet \to \mathcal{J}_1^\bullet$。
于是可选取复形图
$$
\xymatrix{
g'_{1, !}\mathcal{J}_1^\bullet \ar[rr] & &
\mathcal{I}_1^\bullet \\
g'_{1, !}(a')^{-1}\mathcal{J}^\bullet \ar[u] \ar@{=}[r] &
a^{-1}g'_!\mathcal{J}^\bullet \ar[r] &
a^{-1}\mathcal{I}^\bullet \ar[u]
}
$$
，它在同伦意义下交换，且所有箭头均为拟同构；等号来自引理
\ref{more-etale-lemma-proper-f-shriek}，而 $\mathcal{I}_1^\bullet$ 是
$X_{1, \etale}$ 上的 $\Lambda$-模层 K-内射复形。映射
$g_{1, !}Rf'_{1, *}(a')^{-1}K \to Rf_{1, *}g'_{1, !}(a')^{-1}K$ 由下式给出：
$$
g_{1, !}f'_{1, *}\mathcal{J}_1^\bullet =
g_{1, !}f'_{1, !}\mathcal{J}_1^\bullet =
f_{1, !}g'_{1, !}\mathcal{J}_1^\bullet =
f_{1, *}g'_{1, !}\mathcal{J}_1^\bullet \to
f_{1, *}\mathcal{I}_1^\bullet
$$
两个箭头中跨越这 $3$ 个等号的等同均与拉回映射相容；换言之，
阿贝尔层复形的图
$$
\xymatrix{
b^{-1}g_!f'_*\mathcal{J}^\bullet \ar@{=}[d] \ar[r] &
g_{1, !}(b')^{-1}f'_*\mathcal{J}^\bullet \ar[r] &
g_{1, !}f'_{1, *}(a')^{-1}\mathcal{J}^\bullet \ar@{=}[d] \\
b^{-1}f_*g'_!\mathcal{J}^\bullet \ar[r] &
f_{1, *}a^{-1}g'_!\mathcal{J}^\bullet \ar[r] &
f_{1, *}g'_{1, !}(a')^{-1}\mathcal{J}^\bullet
}
$$
交换。要证明这一点，只需把图中的 $g_!, g_{1, !}, g'_!, g'_{1, !}$
分别替换为 $g_*, g_{1, *}, g'_*, g'_{1, *}$ 后证明其交换；这是因为下叹号
函子被定义为 $*$ 函子的子函子，而基变换映射的定义与此相容，见引理
\ref{more-etale-lemma-base-change-f-shriek-separated} 的证明。对这个新图，交换性来自
拉回映射与图的横向、纵向叠合之相容性；见“位点”评注
\ref{sites-remark-compose-base-change} 与
\ref{sites-remark-compose-base-change-horizontal}。因此沿图任一方向绕行，
得到的都是 $f \circ g' = g \circ f'$ 沿 $b$ 基变换的拉回映射。
当然，下图交换：
$$
\xymatrix{
g_{1, !}f'_{1, *}(a')^{-1}\mathcal{J}^\bullet \ar@{=}[d] \ar[r] &
g_{1, !}f'_{1, *}\mathcal{J}_1^\bullet \ar@{=}[d] \\
f_{1, *}g'_{1, !}(a')^{-1}\mathcal{J}^\bullet \ar[r] &
f_{1, *}g'_{1, !}\mathcal{J}_1^\bullet 
}
$$
故要完成证明，只需证明
$$
\xymatrix{
b^{-1}f_*g'_!\mathcal{J}^\bullet \ar[r] \ar[d] &
f_{1, *}a^{-1}g'_!\mathcal{J}^\bullet \ar[r] \ar[d] &
f_{1, *}g'_{1, !}(a')^{-1}\mathcal{J}^\bullet \ar[r] &
f_{1, *}g'_{1, !}\mathcal{J}_1^\bullet \ar[d] \\
b^{-1}f_*\mathcal{I}^\bullet \ar[r] &
f_{1, *}a^{-1}\mathcal{I}^\bullet \ar[rr] & &
f_{1, *}\mathcal{I}_1^\bullet
}
$$
在导出范畴中交换；这由我们先前对映射的选取成立。
\end{proof}

\begin{lemma}
\label{more-etale-lemma-shriek-lqf-and-proper}
考虑概形的交换图
$$
\xymatrix{
X \ar[r]_f \ar[rd]_g & Y \ar[d]^h \\
& Z
}
$$
，其中 $f$ 与 $g$ 局部拟有限，而 $h$ 固有。设 $\Lambda$ 是环。
关于 $K \in D(X_\etale, \Lambda)$ 函子性地，在
$D(Z_\etale, \Lambda)$ 中有典范映射
$$
g_!K \longrightarrow Rh_*(f_!K)
$$
。在以下任一情形，该映射是同构：
(a) $K$ 下有界且其上同调层为挠层；或
(b) $\Lambda$ 是挠环。
\end{lemma}

\begin{proof}
若 $f$ 与 $g$ 分离，这是引理 \ref{more-etale-lemma-shriek-proper-and-open} 的特例。
我们建议读者跳过一般情形的证明，因为我们主要使用 $f$ 与 $g$ 分离的情形。

\medskip\noindent
以 $X_\etale$ 上的 $\Lambda$-模层复形 $\mathcal{K}^\bullet$ 表示 $K$。
选取到 $Y_\etale$ 上的 $\Lambda$-模层 K-内射复形
$\mathcal{I}^\bullet$ 的拟同构
$f_!\mathcal{K}^\bullet \to \mathcal{I}^\bullet$。考虑映射
$$
g_!\mathcal{K}^\bullet =
h_!f_!\mathcal{K}^\bullet =
h_*f_!\mathcal{K}^\bullet
\longrightarrow
h_*\mathcal{I}^\bullet
$$
，其中各等号由引理 \ref{more-etale-lemma-lqf-separated-shriek-composition}
和 \ref{more-etale-lemma-proper-f-shriek} 给出。这个复形映射确定引理陈述中的映射
$g_!K \to Rh_*(f_!K)$。

\medskip\noindent
设 $\Lambda$ 为挠环，即处于情形 (b)。为检验该映射是同构，可在 $Z$ 上
局部地工作。因此可设 $h$ 的纤维维数有界；见“态射”引理
\ref{morphisms-lemma-morphism-finite-type-bounded-dimension}。
于是 $Rh_*$ 的上同调维数有限；见“\'Etale 上同调”引理
\ref{etale-cohomology-lemma-cohomological-dimension-proper}。
故由“导出范畴”引理 \ref{derived-lemma-unbounded-right-derived}，
若能证明对 $q > 0$ 以及 $X_\etale$ 上任意 $\Lambda$-模层
$\mathcal{F}$，均有 $R^qh_*(f_!\mathcal{F}) = 0$，则
$h_*f_!\mathcal{K}^\bullet \to h_*\mathcal{I}^\bullet$ 是拟同构。

\medskip\noindent
注意，$\mathcal{G} = f_!\mathcal{F}$ 是 $Y$ 上的 $\Lambda$-模层；其茎仅可能在
满足 $\kappa(y)/\kappa(h(y))$ 为有限扩张的点 $y \in Y$ 处非零。这由引理
\ref{more-etale-lemma-lqf-f-shriek-stalk} 中对 $f_!\mathcal{F}$ 的茎的描述，以及
$f$ 与 $g$ 均局部拟有限这一事实得出。因此，由固有基变换定理
（“\'Etale 上同调”引理
\ref{etale-cohomology-lemma-proper-base-change-stalk}），只需证明
$H^q(Y_{\overline{z}}, \mathcal{H}) = 0$；这里 $\mathcal{H}$ 是
$\kappa(\overline{z})$ 上固有概形 $Y_{\overline{z}}$ 的一个层，且其支撑包含于
闭点集。所需消没于是由“\'Etale 上同调”引理
\ref{etale-cohomology-lemma-supported-in-closed-points} 得出。

\medskip\noindent
情形 (a) 可用引理 \ref{more-etale-lemma-shriek-proper-and-open} 的证明中完全相同的论证
从情形 (b) 推出（以引理 \ref{more-etale-lemma-lqf-f-shriek-stalk} 代替引理
\ref{more-etale-lemma-lqf-f-shriek-separated-colimits}）。
\end{proof}







\section{通过紧化构造导出下叹号}
\label{more-etale-section-derived-lower-shriek-compactification}

\noindent
设 $f : X \to Y$ 是概形的有限型分离态射，其中 $Y$ 拟紧且拟分离。
选取在 $Y$ 上的紧化 $j : X \to \overline{X}$；见“平坦性的进一步结果”
定理 \ref{flat-theorem-nagata}。设 $\Lambda$ 是环。以
$D^+_{tors}(X_\etale, \Lambda)$ 表示 $D(X_\etale, \Lambda)$ 的严格全且
饱和的三角子范畴，其对象为下有界且上同调层为挠层的对象 $K$。
我们将考虑函子
$$
Rf_! = R\overline{f}_* \circ j_! :
D^+_{tors}(X_\etale, \Lambda)
\longrightarrow
D^+_{tors}(Y_\etale, \Lambda)
$$
，其中 $\overline{f} : \overline{X} \to Y$ 是结构态射。该定义是有意义的：
由评注 \ref{more-etale-remark-lqf-shriek-derived-torsion}，函子 $j_!$ 把
$D^+_{tors}(X_\etale, \Lambda)$ 送入
$D^+_{tors}(\overline{X}_\etale, \Lambda)$；由“\'Etale 上同调”引理
\ref{etale-cohomology-lemma-torsion-direct-image}，$R\overline{f}_*$ 把
$D^+_{tors}(\overline{X}_\etale, \Lambda)$ 送入
$D^+_{tors}(Y_\etale, \Lambda)$。若 $\Lambda$ 是挠环，则定义
$$
Rf_! = R\overline{f}_* \circ j_! :
D(X_\etale, \Lambda)
\longrightarrow
D(Y_\etale, \Lambda)
$$
下面是必不可少的引理。

\begin{lemma}
\label{more-etale-lemma-shriek-well-defined}
设 $f : X \to Y$ 是拟紧拟分离概形之间的有限型分离态射。上述构造的函子
$Rf_!$ 在典范同构意义下与紧化的选取无关。
\end{lemma}

\begin{proof}
我们将在 $\Lambda$ 为挠环时对函子
$Rf_! : D(X_\etale, \Lambda) \to D(Y_\etale, \Lambda)$ 证明此事；
函子 $Rf_! : D^+_{tors}(X_\etale, \Lambda) \to
D^+_{tors}(Y_\etale, \Lambda)$ 的情形完全同样地证明。

\medskip\noindent
考虑 $X$ 在 $Y$ 上的紧化所成的范畴。由“平坦性的进一步结果”定理
\ref{flat-theorem-nagata} 以及引理
\ref{flat-lemma-compactifications-cofiltered} 和
\ref{flat-lemma-compactifyable}，该范畴是余滤过的。对紧化的每个选取
$$
j : X \to \overline{X},\quad \overline{f} : \overline{X} \to Y
$$
，上述构造配以函子 $R\overline{f}_* \circ j_! :
D(X_\etale, \Lambda) \to D(Y_\etale, \Lambda)$。更具体一些：给定
$X_\etale$ 上的 $\Lambda$-模层复形 $\mathcal{K}^\bullet$，我们选取一个
到 $\overline{X}_\etale$ 上某个 $\Lambda$-模层 K-内射复形的拟同构
$j_!\mathcal{K}^\bullet \to \mathcal{I}^\bullet$。我们的函子遂把
$\mathcal{K}^\bullet$ 送到 $\overline{f}_*\mathcal{I}^\bullet$。

\medskip\noindent
设给定 $Y$ 上紧化 $j_i : X \to \overline{X}_i$ 之间的态射
$g : \overline{X}_1 \to \overline{X}_2$。则得到同构
$$
R\overline{f}_{2, *} \circ j_{2, !} =
R\overline{f}_{2, *} \circ Rg_* \circ j_{1, !} =
R\overline{f}_{1, *} \circ j_{1, !}
$$
；第一个等号使用引理 \ref{more-etale-lemma-shriek-lqf-and-proper}。

\medskip\noindent
为完成证明，因 $X$ 在 $Y$ 上的紧化范畴余滤过，只需证明：$X$ 在 $Y$ 上
各紧化态射的复合，被送到函子同构的复合\footnote{具体地，若
$\alpha, \beta : F \to G$ 是函子态射，而 $\gamma : G \to H$ 是函子同构，
且 $\gamma \circ \alpha = \gamma \circ \beta$，则可推出
$\alpha = \beta$。}。为此，设 $j_3 : X \to \overline{X}_3$ 是第三个紧化，
而 $h : \overline{X}_2 \to \overline{X}_3$ 是紧化态射。我们必须证明复合
$$
R\overline{f}_{3, *} \circ j_{3, !} =
R\overline{f}_{3, *} \circ Rh_* \circ j_{2, !} =
R\overline{f}_{2, *} \circ j_{2, !} =
R\overline{f}_{2, *} \circ Rg_* \circ j_{1, !} =
R\overline{f}_{1, *} \circ j_{1, !}
$$
等于仅使用 $j_3$、$g \circ h$ 与 $j_1$ 构造出的函子同构。计算表明，只需证明
引理 \ref{more-etale-lemma-shriek-lqf-and-proper} 的映射之复合
$$
j_{3, !} \to Rh_* \circ j_{2, !} \to Rh_* \circ Rg_* \circ j_{1, !}
$$
通过等同 $R(h \circ g)_* = Rh_* \circ Rg_*$，与相应映射
$j_{3, !} \to R(h \circ g)_* \circ j_{1, !}$ 一致。
由于引理 \ref{more-etale-lemma-shriek-lqf-and-proper} 的映射是引理
\ref{more-etale-lemma-shriek-proper-and-open} 的映射的特例（因为 $j_1$ 与 $j_2$ 分离），
这立即由引理 \ref{more-etale-lemma-shriek-proper-and-open-compose} 得出。
\end{proof}

\begin{lemma}
\label{more-etale-lemma-shriek-composition}
设 $f : X \to Y$ 与 $g : Y \to Z$ 是拟紧拟分离概形之间的有限型分离态射。
则有典范同构 $Rg_! \circ Rf_! \to R(g \circ f)_!$。
\end{lemma}

\begin{proof}
选取 $Y$ 在 $Z$ 上的紧化 $i : Y \to \overline{Y}$，再选取 $X$ 在
$\overline{Y}$ 上的紧化 $X \to \overline{X}$。这里两次使用“平坦性的进一步结果”
定理 \ref{flat-theorem-nagata} 和引理 \ref{flat-lemma-compactifyable}。
令 $U$ 为 $Y$ 在 $\overline{X}$ 中的逆像，于是得到交换图
$$
\xymatrix{
X \ar[r]_j \ar[d]_f &
U \ar[dl]^{f'} \ar[r]_{j'} &
\overline{X} \ar[dl]^{\overline{f}} \\
Y \ar[r]_i \ar[d]_g &
\overline{Y} \ar[dl]^{\overline{g}} \\
Z
}
$$
于是有
\begin{align*}
R(g \circ f)_!
& =
R(\overline{g} \circ \overline{f})_* \circ (j' \circ j)_! \\
& =
R\overline{g}_* \circ R\overline{f}_* \circ j'_! \circ j_! \\
& =
R\overline{g}_* \circ i_! \circ Rf'_* \circ j_! \\
& =
Rg_! \circ Rf_!
\end{align*}
第一个等号是 $R(g \circ f)_!$ 的定义。第二个等号使用等同
$R(\overline{g} \circ \overline{f})_* = R\overline{g}_* \circ R\overline{f}_*$
以及引理 \ref{more-etale-lemma-f-shriek-composition} 的等同
$(j' \circ j)_! = j'_! \circ j_!$。第三个等号所用的等同
$i_! \circ Rf'_* \to R\overline{f}_* \circ j_!$ 来自引理
\ref{more-etale-lemma-shriek-proper-and-open}。最后的第四个等号是 $Rg_!$ 与 $Rf_!$ 的定义。
为完成证明，我们要说明该同构与所作选取无关。

\medskip\noindent
设有两个图
$$
\vcenter{
\xymatrix{
X \ar[r]_{j_1} \ar[d] &
U_1 \ar[dl]^{f_1} \ar[r]_{j'_1} &
\overline{X}_1 \ar[dl]^{\overline{f}_1} \\
Y \ar[r]_{i_1} \ar[d] &
\overline{Y}_1 \ar[dl]^{\overline{g}_1} \\
Z
}
}
\quad\text{且}\quad
\vcenter{
\xymatrix{
X \ar[r]_{j_2} \ar[d] &
U_2 \ar[dl]^{f_2} \ar[r]_{j'_2} &
\overline{X}_2 \ar[dl]^{\overline{f}_2} \\
Y \ar[r]_{i_2} \ar[d] &
\overline{Y}_2 \ar[dl]^{\overline{g}_2} \\
Z
}
}
$$
首先可选取 $Y$ 在 $Z$ 上的紧化 $i : Y \to \overline{Y}$，使其同时支配
$\overline{Y}_1$ 与 $\overline{Y}_2$；见“平坦性的进一步结果”引理
\ref{flat-lemma-compactifications-cofiltered}。由“平坦性的进一步结果”引理
\ref{flat-lemma-right-multiplicative-system} 以及“范畴”引理
\ref{categories-lemma-morphisms-right-localization} 和
\ref{categories-lemma-equality-morphisms-right-localization}，可选取 $X$ 在
$\overline{Y}$ 上的紧化 $X \to \overline{X}$，并配有态射
$\overline{X} \to \overline{X}_1$ 与 $\overline{X} \to \overline{X}_2$，
使复合 $\overline{X} \to \overline{Y} \to \overline{Y}_1$ 等于复合
$\overline{X} \to \overline{X}_1 \to \overline{Y}_1$，且复合
$\overline{X} \to \overline{Y} \to \overline{Y}_2$ 等于复合
$\overline{X} \to \overline{X}_2 \to \overline{Y}_2$。因此，只需在有如下
交换图时，比较这两个图所确定的映射：
$$
\xymatrix{
X \ar[rr]_{j_1} \ar@{=}[d] & &
U_1 \ar[d]^{h'} \ar[ddll] \ar[rr]_{j'_1} & &
\overline{X}_1 \ar[d]^h \ar[ddll] \\
X \ar'[r][rr]^-{j_2} \ar[d] & &
U_2 \ar'[dl][ddll] \ar'[r][rr]^-{j'_2} & &
\overline{X}_2 \ar[ddll] \\
Y \ar[rr]^{i_1} \ar@{=}[d] & & \overline{Y}_1 \ar[d]^k \\
Y \ar[rr]^{i_2} \ar[d] & & \overline{Y}_2 \ar[dll] \\
Z
}
$$
下列每个方块
$$
\xymatrix{
X \ar[r]_{j_1} \ar[d]_{\text{id}} \ar@{}[dr]|A &
U_1 \ar[d]^{h'} \\
X \ar[r]^{j_2} &
U_2
}
\quad
\xymatrix{
U_2 \ar[r]_{j_2'} \ar[d]_{f_2} \ar@{}[dr]|B &
\overline{X}_2 \ar[d]^{\overline{f}_2} \\
Y \ar[r]^{i_2} &
\overline{Y}_2
}
\quad
\xymatrix{
U_1 \ar[r]_{j_1'} \ar[d]_{f_1} \ar@{}[dr]|C &
\overline{X}_1 \ar[d]^{\overline{f}_1} \\
Y \ar[r]^{i_1} &
\overline{Y}_1
}
\quad
\xymatrix{
Y \ar[r]_{i_1} \ar[d]_{\text{id}} \ar@{}[dr]|D &
\overline{Y}_1 \ar[d]^k \\
Y \ar[r]^{i_2} &
\overline{Y}_2
}
\quad
\xymatrix{
X \ar[r]_{j_1' \circ j_1} \ar[d]_{\text{id}} \ar@{}[dr]|E &
\overline{X}_1 \ar[d]^h \\
X \ar[r]^{j_2} &
\overline{X}_2
}
$$
均产生如下同构：
\begin{align*}
\gamma_A & :
j_{2, !} \to Rh'_* \circ j_{1, !}  \\
\gamma_B & :
i_{2, !} \circ Rf_{2, *} \to R\overline{f}_{2, *} \circ j'_{2, !} \\
\gamma_C & :
i_{1, !} \circ Rf_{1, *} \to R\overline{f}_{1, *} \circ j'_{1, !} \\
\gamma_D & :
i_{2, !} \to Rk_* \circ i_{1, !} \\
\gamma_E & :
j_{2, !} \to Rh_* \circ (j'_1 \circ j_1)_!
\end{align*}
；这是通过应用引理 \ref{more-etale-lemma-shriek-proper-and-open} 的映射得到的
（当左侧竖箭头为恒等时，它与引理 \ref{more-etale-lemma-shriek-lqf-and-proper}
的映射相同）。记
\begin{align*}
F_1 & = Rf_{1, *} \circ j_{1, !} \\
F_2 & = Rf_{2, *} \circ j_{2, !} \\
G_1 & = R\overline{g}_{1, *} \circ i_{1, !} \\
G_2 & = R\overline{g}_{2, *} \circ i_{2, !} \\
C_1 & = R(\overline{g}_1 \circ \overline{f}_1)_* \circ (j'_1 \circ j_1)_! \\
C_2 & = R(\overline{g}_2 \circ \overline{f}_2)_* \circ (j'_2 \circ j_2)_!
\end{align*}
本证明第一段以及引理 \ref{more-etale-lemma-shriek-well-defined} 中给出的构造使用
\begin{enumerate}
\item $\gamma_C$ 作为映射 $G_1 \circ F_1 \to C_1$；
\item $\gamma_B$ 作为映射 $G_2 \circ F_2 \to C_2 $；
\item $\gamma_A$ 作为映射 $F_2 \to F_1$；
\item $\gamma_D$ 作为映射 $G_2 \to G_1$；以及
\item $\gamma_E$ 作为映射 $C_2 \to C_1$。
\end{enumerate}
因此必须证明图
$$
\xymatrix{
C_2 \ar[rr]_{\gamma_E} & &
C_1 \\
G_2 \circ F_2 \ar[rr]^{\gamma_D \circ \gamma_A} \ar[u]^{\gamma_B} & &
G_1 \circ F_1 \ar[u]_{\gamma_C}
}
$$
交换。我们将使用引理 \ref{more-etale-lemma-shriek-proper-and-open-compose} 与
\ref{more-etale-lemma-shriek-proper-and-open-compose-horizontal}，并沿用评注
\ref{more-etale-remark-going-around} 中（略有滥用的）记号；特别地，从记号中略去与恒等变换的
$\star$ 乘积。可写成 $\gamma_E = \gamma_F \circ \gamma_A$，其中
$$
\xymatrix{
U_1 \ar[r]_{j'_1} \ar[d]_{h'} \ar@{}[rd]|F &
\overline{X}_1 \ar[d]^h \\
U_2 \ar[r]^{j'_2} &
\overline{X}_2
}
$$
于是
$$
\gamma_E \circ \gamma_B = \gamma_F \circ \gamma_A  \circ \gamma_B
= \gamma_F \circ \gamma_B \circ \gamma_A
$$
；最后一个等号成立，是因为方块 $A$ 与 $B$ 仅在一点相交（类似于评注
\ref{more-etale-remark-going-around} 末尾的论证）。因此只需证明
$\gamma_C \circ \gamma_D = \gamma_F \circ \gamma_B$。两者都等于方块
$$
\xymatrix{
U_1 \ar[r] \ar[d] & \overline{X}_1 \ar[d] \\
Y \ar[r] & \overline{Y}_2
}
$$
的映射，故结论成立。
\end{proof}

\begin{lemma}
\label{more-etale-lemma-pseudo-functor}
设 $f : X \to Y$、$g : Y \to Z$、$h : Z \to T$ 是拟紧拟分离概形之间的
有限型分离态射。则引理 \ref{more-etale-lemma-shriek-composition} 的各同构所成的图
$$
\xymatrix{
Rh_! \circ Rg_! \circ Rf_! \ar[r]_{\gamma_C} \ar[d]^{\gamma_A} &
R(h \circ g)_! \circ Rf_! \ar[d]_{\gamma_{A + B}} \\
Rh_! \circ R(g \circ f)_! \ar[r]^{\gamma_{B + C}} &
R(h \circ g \circ f)_!
}
$$
交换（各 $\gamma$ 的含义见证明）。
\end{lemma}

\begin{proof}
为此，先选取 $Z$ 在 $T$ 上的紧化 $\overline{Z}$，再选取 $Y$ 在
$\overline{Z}$ 上的紧化 $\overline{Y}$，继而选取 $X$ 在 $\overline{Y}$ 上的
紧化 $\overline{X}$。这里使用“平坦性的进一步结果”定理
\ref{flat-theorem-nagata} 和引理 \ref{flat-lemma-compactifyable}。
令 $W \subset \overline{Y}$ 为 $Z$ 在 $\overline{Y} \to \overline{Z}$ 下的逆像，
并令 $U \subset V \subset \overline{X}$ 为 $Y \subset W$ 在
$\overline{X} \to \overline{Y}$ 下的逆像。由此得到图
$$
\xymatrix{
X \ar[d]_f \ar[r] & U \ar[r] \ar[d] \ar@{}[dr]|A &
V \ar[d] \ar[r] \ar@{}[rd]|B & \overline{X} \ar[d] \\
Y \ar[d]_g \ar[r] & Y \ar[r] \ar[d] & W \ar[r] \ar[d] \ar@{}[rd]|C &
\overline{Y} \ar[d] \\
Z \ar[d]_h \ar[r] & Z \ar[d] \ar[r] & Z \ar[d] \ar[r] & \overline{Z} \ar[d] \\
T \ar[r] & T \ar[r] & T \ar[r] & T
}
$$
不引入大量记号而完全仿照引理 \ref{more-etale-lemma-shriek-composition} 的证明论证，
可见第一个显示图中的各映射，按引理陈述的图所示，分别使用引理
\ref{more-etale-lemma-shriek-proper-and-open} 对矩形 $A + B$、$B + C$、$A$ 与 $C$
给出的映射。由引理 \ref{more-etale-lemma-shriek-proper-and-open-compose} 和
\ref{more-etale-lemma-shriek-proper-and-open-compose-horizontal}，有
$\gamma_{A + B} = \gamma_B \circ \gamma_A$ 以及
$\gamma_{B + C} = \gamma_B \circ \gamma_C$。因此，只要
$\gamma_A \circ \gamma_C = \gamma_C \circ \gamma_A$，所需等式即成立。
后一个等式确实成立，因为方块 $A$ 与 $C$ 仅在一点相交（类似于评注
\ref{more-etale-remark-going-around} 末尾的论证）。
\end{proof}

\begin{lemma}
\label{more-etale-lemma-base-change-shriek}
考虑笛卡尔方块
$$
\xymatrix{
X' \ar[r]_{g'} \ar[d]_{f'} & X \ar[d]^f \\
Y' \ar[r]^g & Y
}
$$
，其中各概形拟紧且拟分离，而 $f$ 分离且有限型。则有典范同构
$$
g^{-1} \circ Rf_! \to Rf'_! \circ (g')^{-1}
$$
此外，这些同构与引理 \ref{more-etale-lemma-shriek-composition} 的同构相容。
\end{lemma}

\begin{proof}
选取在 $Y$ 上的紧化 $j : X \to \overline{X}$，并以
$\overline{f} : \overline{X} \to Y$ 表示结构态射。以
$j' : X' \to \overline{X}'$ 与 $\overline{f}' : \overline{X}' \to Y'$
表示 $j$ 与 $\overline{f}$ 的基变换。由于
$Rf_! = R\overline{f}_* \circ j_!$ 且
$Rf'_! = R\overline{f}'_* \circ j'_!$，可通过下式构造该同构：
$$
g^{-1} \circ R\overline{f}_* \circ j_! \to
R\overline{f}'_* \circ (\overline{g}')^{-1} \circ j_! \to
R\overline{f}'_* \circ j'_! \circ (g')^{-1} 
$$
其中第一个箭头是固有基变换定理给出的同构（在下有界挠情形，使用
“\'Etale 上同调”引理 \ref{etale-cohomology-lemma-proper-base-change}；
在 $\Lambda$ 为挠环的情形，使用“\'Etale 上同调”引理
\ref{etale-cohomology-lemma-proper-base-change-mod-n}），第二个箭头是引理
\ref{more-etale-lemma-base-change-f-shriek-separated} 的同构。

\medskip\noindent
为完成证明，必须说明两点：第一，如此得到的函子同构与紧化的选取无关；
第二，若把引理中的两个基变换图竖向叠合，则这些基变换同构与引理
\ref{more-etale-lemma-shriek-composition} 的同构相容。一个略去的直接论证表明，
若能证明以下各同构与基变换相容，则两点均成立：
\begin{enumerate}
\item 当 $f : X \to Y$ 与 $g : Y \to Z$ 固有时，
$Rg_* \circ Rf_* = R(g \circ f)_*$；
\item 当 $f : X \to Y$ 与 $g : Y \to Z$ 分离且拟有限时，
$g_! \circ f_! = (g \circ f)_!$；以及
\item 当 $f : X \to Y$ 与 $f' : X' \to Y'$ 固有，而
$g : Y' \to Y$ 与 $g' : X' \to X$ 分离且拟有限，并满足
$f \circ g' = g \circ f'$ 时，$g_! \circ Rf'_* = Rf_* \circ g'_!$。
\end{enumerate}
对 (1)，这由“位点上同调”评注
\ref{sites-cohomology-remark-compose-base-change} 成立；对 (2)，由评注
\ref{more-etale-remark-f-shriek-base-change-composition} 成立；对 (3)，由引理
\ref{more-etale-lemma-shriek-proper-and-open-base-change} 成立。
\end{proof}

\begin{remark}
\label{more-etale-remark-shriek-to-star}
设 $f : X \to Y$ 是概形的有限型分离态射，其中 $Y$ 拟紧且拟分离。
下面将构造映射
$$
Rf_!K \longrightarrow Rf_*K
$$
；它关于 $D^+_{tors}(X_\etale, \Lambda)$ 中的 $K$ 函子性地成立；若
$\Lambda$ 为挠环，则同样关于 $D(X_\etale, \Lambda)$ 中的对象函子性地成立。
两种情形下，该函子变换均与下列同构相容：
\begin{enumerate}
\item 引理 \ref{more-etale-lemma-shriek-composition} 的同构
$Rg_! \circ Rf_! \to R(g \circ f)_!$，以及“位点上同调”引理
\ref{sites-cohomology-lemma-derived-pushforward-composition} 的同构
$Rg_* \circ Rf_* \to R(g \circ f)_*$；以及
\item 引理 \ref{more-etale-lemma-base-change-shriek} 的同构
$g^{-1} \circ Rf_! \to Rf'_! \circ (g')^{-1}$，以及“位点上同调”评注
\ref{sites-cohomology-remark-base-change} 的基变换映射。
\end{enumerate}
具体地，选取在 $Y$ 上的紧化 $j : X \to \overline{X}$，并以
$\overline{f} : \overline{X} \to Y$ 表示结构态射。由于
$Rf_! = R\overline{f}_* \circ j_!$ 且
$Rf_* = R\overline{f}_* \circ Rj_*$，只需构造函子变换 $j_! \to Rj_*$。
为此，使用“\'Etale 上同调”引理
\ref{etale-cohomology-lemma-shriek-into-star-separated-etale} 的典范变换
$j_! \to j_*$。我们略去所得变换与紧化选取无关的证明，也略去相容性
(1) 与 (2) 的证明。
\end{remark}













\section{导出下叹号的性质}
\label{more-etale-section-derived-lower-shriek-properties}

\noindent
下面列出导出下叹号的一些性质。

\begin{lemma}
\label{more-etale-lemma-derived-lower-shriek-commute-direct-sums}
设 $f : X \to Y$ 是拟紧拟分离概形之间的有限型分离态射。设 $\Lambda$ 是环。
\begin{enumerate}
\item 设 $K_i \in D^+_{tors}(X_\etale, \Lambda)$，$i \in I$，是一族对象。
假设给定 $a \in \mathbf{Z}$，使得当 $n < a$ 且 $i \in I$ 时
$H^n(K_i) = 0$。则 $Rf_!(\bigoplus_i K_i) = \bigoplus_i Rf_!K_i$。
\item 若 $\Lambda$ 为挠环，则函子
$Rf_! : D(X_\etale, \Lambda) \to D(Y_\etale, \Lambda)$
与直和交换。
\end{enumerate}
\end{lemma}

\begin{proof}
由构造，只需在 $f$ 为开浸入以及 $f$ 为固有态射时证明。对概形的任意开浸入
$j : U \to X$，函子 $j_! : D(U_\etale) \to D(X_\etale)$ 是拉回
$j^{-1} : D(X_\etale) \to D(U_\etale)$ 的左伴随，因而与直和交换；
见“位点上同调”引理 \ref{sites-cohomology-lemma-adjoint-lower-shriek-restrict}。
在固有情形有 $Rf_! = Rf_*$；在下有界情形，结论由“\'Etale 上同调”引理
\ref{etale-cohomology-lemma-direct-sum-bounded-below-pushforward} 得出；
当系数环 $\Lambda$ 为挠环时，结论由“\'Etale 上同调”引理
\ref{etale-cohomology-lemma-proper-mod-n-direct-sums} 得出。
\end{proof}

\begin{lemma}
\label{more-etale-lemma-derived-lower-shriek-bounded}
设 $f : X \to Y$ 是拟紧拟分离概形之间的有限型分离态射，$\Lambda$ 是环。
第 \ref{more-etale-section-derived-lower-shriek-compactification} 节中构造的函子 $Rf_!$
在如下意义下有界：存在整数 $N$，使得对
$E \in D^+_{tors}(X_\etale, \Lambda)$，或者在 $\Lambda$ 为挠环时对
$E \in D(X_\etale, \Lambda)$，有
\begin{enumerate}
\item 当 $i \leq a$ 时，$H^i(Rf_!(\tau_{\leq a}E) \to H^i(Rf_!(E))$ 是同构；
\item 当 $i \geq b$ 时，$H^i(Rf_!(E)) \to H^i(Rf_!(\tau_{\geq b - N}E))$ 是同构；
\item 若对某个 $-\infty \leq a \leq b \leq \infty$，当
$i \not \in [a, b]$ 时 $H^i(E) = 0$，则当 $i \not \in [a, b + N]$ 时
$H^i(Rf_!(E)) = 0$。
\end{enumerate}
\end{lemma}

\begin{proof}
设 $\Lambda$ 为挠环，并考虑函子
$Rf_! : D(X_\etale, \Lambda) \to D(Y_\etale, \Lambda)$。
由构造，只需在 $f$ 为开浸入以及 $f$ 为固有态射时证明。对概形的任意开浸入
$j : U \to X$，函子 $j_! : D(U_\etale) \to D(X_\etale)$ 正合，
故在此情形陈述以 $N = 0$ 成立。若 $f$ 固有，则 $Rf_! = Rf_*$，
即它是右导出函子。因此左侧的界由“导出范畴”引理
\ref{derived-lemma-negative-vanishing} 得出。此外，在此情形，
由“态射”引理 \ref{morphisms-lemma-morphism-finite-type-bounded-dimension}
和“\'Etale 上同调”引理
\ref{etale-cohomology-lemma-cohomological-dimension-proper}，
$f_* : \textit{Mod}(X_\etale, \Lambda) \to
\textit{Mod}(Y_\etale, \Lambda)$ 的上同调维数有界。因此由“导出范畴”引理
\ref{derived-lemma-unbounded-right-derived} 得到结论。

\medskip\noindent
其次，设 $\Lambda$ 任意，并考虑函子
$Rf_! : D^+_{tors}(X_\etale, \Lambda) \to D^+_{tors}(Y_\etale, \Lambda)$。
仍然立即化归到 $f$ 固有且 $Rf_! = Rf_*$ 的情形。(1) 仍是立即的。
为证明 (3)，可对 $b - a$ 作归纳，并使用截断的特异三角以及“\'Etale 上同调”引理
\ref{etale-cohomology-lemma-cohomological-dimension-proper}。
(2) 由 (3) 得出。略去细节。
\end{proof}

\begin{lemma}
\label{more-etale-lemma-Rf-shriek-for-quasi-finite}
设 $f : X \to Y$ 是拟紧拟分离概形之间的拟有限分离态射。则第
\ref{more-etale-section-derived-lower-shriek-compactification} 节中构造的函子 $Rf_!$
与第 \ref{more-etale-section-derived-duality-locally-quasi-finite} 节中构造的函子
$f_! : D(X_\etale, \Lambda) \to D(Y_\etale, \Lambda)$ 在二者共同定义域上的
限制一致。
\end{lemma}

\begin{proof}
由 Zariski 主定理（“态射的进一步结果”引理
\ref{more-morphisms-lemma-quasi-finite-separated-pass-through-finite}），
可找到开浸入 $j : X \to \overline{X}$ 与有限态射
$\overline{f} : \overline{X} \to Y$，使 $f = \overline{f} \circ j$。
由构造有 $Rf_! = R\overline{f}_* \circ j_!$。由于 $\overline{f}$ 有限，
由“\'Etale 上同调”命题
\ref{etale-cohomology-proposition-finite-higher-direct-image-zero}，有
$R\overline{f}_* = \overline{f}_*$。又例如由引理
\ref{more-etale-lemma-compactify-f-shriek-separated}，有
$\overline{f}_* \circ j_! = f_!$，故引理成立。
\end{proof}

\begin{lemma}
\label{more-etale-lemma-relative-mayer-vietoris}
设 $f : X \to Y$ 是拟紧拟分离概形之间的有限型分离态射。设 $U$ 与 $V$
是 $X$ 的拟紧开集，且 $X = U \cup V$。以 $a : U \to Y$、
$b : V \to Y$ 与 $c : U \cap V \to Y$ 表示 $f$ 的限制。设 $\Lambda$ 是环。
对 $D^+_{tors}(X_\etale, \Lambda)$ 中的 $K$，或者当 $\Lambda$ 为挠环时，
对 $K \in D(X_\etale, \Lambda)$，有特异三角
$$
Rc_!(K|_{U \cap V}) \to
Ra_!(K|_U) \oplus Rb_!(K|_V) \to
Rf_!K \to 
Rc_!(K|_{U \cap V})[1]
$$
成立于 $D(Y_\etale, \Lambda)$ 中。
\end{lemma}

\begin{proof}
这由引理 \ref{more-etale-lemma-mayer-vietoris-shriek}、引理
\ref{more-etale-lemma-shriek-composition} 给出的 $Rf_! \circ Rj_{U!} = Ra_!$，以及引理
\ref{more-etale-lemma-Rf-shriek-for-quasi-finite} 给出的 $Rj_{U!} = j_{U!}$ 得出。
\end{proof}

\begin{lemma}
\label{more-etale-lemma-relative-triangle-associated-to-open}
设 $f : X \to Y$ 是拟紧拟分离概形之间的有限型分离态射。设 $U$ 是 $X$
的拟紧开集，其补集为 $Z \subset X$。以 $g : U \to Y$ 和 $h : Z \to Y$
表示 $f$ 的限制。设 $\Lambda$ 是环。对 $D^+_{tors}(X_\etale, \Lambda)$
中的 $K$，或者当 $\Lambda$ 为挠环时，对 $K \in D(X_\etale, \Lambda)$，
有特异三角
$$
Rg_!(K|_U) \to
Rf_!K \to
Rh_!(K|_Z) \to
Rg_!(K|_U)[1]
$$
成立于 $D(Y_\etale, \Lambda)$ 中。
\end{lemma}

\begin{proof}
这由引理 \ref{more-etale-lemma-triangle-associated-to-open}、引理
\ref{more-etale-lemma-shriek-composition} 给出的 $Rf_! \circ Rj_! = Rg_!$ 与
$Rf_! \circ Ri_!$，以及引理 \ref{more-etale-lemma-Rf-shriek-for-quasi-finite} 给出的
$Rj_! = j_!$ 和 $Ri_! = i_! = i_*$ 得出。
\end{proof}

\begin{lemma}
\label{more-etale-lemma-shriek-and-thickening}
设 $f' : X' \to Y$ 是拟紧拟分离概形之间的有限型分离态射。设
$i : X \to X'$ 是增厚，并记 $f = f' \circ i$。设 $\Lambda$ 是环。
对 $D^+_{tors}(X'_\etale, \Lambda)$ 中的 $K'$，或者当 $\Lambda$ 为挠环时，
对 $K' \in D(X'_\etale, \Lambda)$，有 $Rf_!i^{-1}K' = Rf'_!K'$。
\end{lemma}

\begin{proof}
这成立是因为，由小 \'etale 拓扑斯的拓扑不变性（“\'Etale 上同调”定理
\ref{etale-cohomology-theorem-topological-invariance}），$i^{-1}$ 与
$i_* = i_!$ 是范畴的互逆等价，并且可以应用引理
\ref{more-etale-lemma-shriek-composition}。
\end{proof}

\begin{lemma}
\label{more-etale-lemma-projection-formula-Rf-lower-shriek}
设 $f : X \to Y$ 是拟紧拟分离概形之间的有限型分离态射。设 $\Lambda$ 是挠环，
$E \in D(X_\etale, \Lambda)$ 且 $K \in D(Y_\etale, \Lambda)$。则
$$
Rf_!E \otimes_\Lambda^\mathbf{L} K =
Rf_!(E \otimes_\Lambda^\mathbf{L} f^{-1}K)
$$
成立于 $D(Y_\etale, \Lambda)$ 中。
\end{lemma}

\begin{proof}
按 $Rf_!$ 的构造选取 $j : X \to \overline{X}$ 与
$\overline{f} : \overline{X} \to Y$。由“位点上同调”引理
\ref{sites-cohomology-lemma-j-shriek-and-tensor}，有
$j_!E \otimes_\Lambda^\mathbf{L} \overline{f}^{-1}K =
j_!(E \otimes_\Lambda^\mathbf{L} f^{-1}K)$。再应用“\'Etale 上同调”引理
\ref{etale-cohomology-lemma-projection-formula-proper-mod-n}，并使用
$f^{-1} = j^{-1} \circ \overline{f}^{-1}$ 与
$Rf_! = R\overline{f}_*j_!$，即得结论。
\end{proof}

\begin{remark}
\label{more-etale-remark-Rf-lower-shriek-change-of-rings}
设 $\Lambda_1 \to \Lambda_2$ 是挠环同态，$f : X \to Y$ 是拟紧拟分离概形之间的
有限型分离态射。图
$$
\xymatrix{
D(X_\etale, \Lambda_2) \ar[r]_{res} \ar[d]_{Rf_!} &
D(X_\etale, \Lambda_1) \ar[d]^{Rf_!} \\
D(Y_\etale, \Lambda_2) \ar[r]^{res} &
D(Y_\etale, \Lambda_1)
}
$$
交换；其中 $res$ 是“限制”函子，它通过给定环映射把 $\Lambda_2$-模视为
$\Lambda_1$-模。对第 \ref{more-etale-section-derived-lower-shriek-compactification} 节中
那样的分解 $f = \overline{f} \circ j$，写成
$Rf_! = R\overline{f}_* \circ j_!$。直接检验可见结论对 $j_!$ 成立，
而由“位点上同调”引理
\ref{sites-cohomology-lemma-modules-abelian-unbounded}，结论对
$R\overline{f}_*$ 成立。另一方面，图
$$
\xymatrix{
D(X_\etale, \Lambda_1)
\ar[r]_{- \otimes_{\Lambda_1}^\mathbf{L} \Lambda_2} \ar[d]_{Rf_!} &
D(X_\etale, \Lambda_2) \ar[d]^{Rf_!} \\
D(Y_\etale, \Lambda_1)
\ar[r]^{- \otimes_{\Lambda_1}^\mathbf{L} \Lambda_2} &
D(Y_\etale, \Lambda_2)
}
$$
也交换；这由引理 \ref{more-etale-lemma-projection-formula-Rf-lower-shriek} 得出。
\end{remark}

\begin{remark}
\label{more-etale-remark-projection-formula-for-internal-hom}
设 $f : X \to Y$ 是拟紧拟分离概形之间的有限型分离态射，$\Lambda$ 是挠系数环，
而 $K$ 与 $L$ 是 $D(X_\etale, \Lambda)$ 的对象。我们断言有典范映射
$$
\alpha :
Rf_*R\SheafHom_\Lambda(K, L)
\longrightarrow
R\SheafHom_\Lambda(Rf_!K, Rf_!L)
$$
关于 $K$ 与 $L$ 函子性地成立。具体地，按 $Rf_!$ 的构造选取
$j : X \to \overline{X}$ 与 $\overline{f} : \overline{X} \to Y$。
先定义映射
$$
\beta :
Rj_*R\SheafHom_\Lambda(K, L)
\longrightarrow
R\SheafHom_\Lambda(j_!K, j_!L)
$$
由导出范畴中内部 Hom 的构造，这等价于定义映射
$$
\beta' :
Rj_*R\SheafHom_\Lambda(K, L) \otimes_\Lambda^\mathbf{L} j_!K
\longrightarrow
j_!L
$$
见“位点上同调”第 \ref{sites-cohomology-section-internal-hom} 节。
$\beta'$ 的源等于
$$
j_!\left(R\SheafHom_\Lambda(K, L) \otimes_\Lambda^\mathbf{L} K\right)
$$
；这由“位点上同调”引理
\ref{sites-cohomology-lemma-j-shriek-and-tensor} 得出。因此可令
$\beta' = j_!\beta''$，其中
$\beta'' : R\SheafHom_\Lambda(K, L) \otimes_\Lambda^\mathbf{L} K \to L$
通过上述内部 Hom 的泛性质对应于 $R\SheafHom_\Lambda(K, L)$ 上的恒等态射。
由“位点上同调”评注
\ref{sites-cohomology-remark-projection-formula-for-internal-hom}，有典范映射
$$
\gamma :
R\overline{f}_*R\SheafHom_\Lambda(j_!K, j_!L)
\longrightarrow
R\SheafHom_\Lambda(R\overline{f}_*j_!K, R\overline{f}_*j_!L)
$$
由于 $Rf_! = R\overline{f}_*j_!$，并且（由 Leray）
$Rf_* = R\overline{f}_* Rj_*$，得到所需映射
$\alpha = \gamma \circ R\overline{f}_*\beta$。
\end{remark}









\section{导出上叹号}
\label{more-etale-section-derived-upper-shriek}

\noindent
我们用一个 Brown 可表示性定理得到 $Rf^!$。

\begin{lemma}
\label{more-etale-lemma-upper-shriek-derived}
设 $f : X \to Y$ 是拟紧拟分离概形之间的有限型分离态射，$\Lambda$ 是挠系数环。
函子
$Rf_! : D(X_\etale, \Lambda) \to D(Y_\etale, \Lambda)$
有右伴随
$Rf^! : D(Y_\etale, \Lambda) \to D(X_\etale, \Lambda)$.
\end{lemma}

\begin{proof}
这由“内射对象”命题 \ref{injectives-proposition-brown} 以及上面的引理
\ref{more-etale-lemma-derived-lower-shriek-commute-direct-sums} 得出。
\end{proof}

\begin{lemma}
\label{more-etale-lemma-Rf-upper-shriek-for-quasi-finite}
设 $f : X \to Y$ 是拟紧拟分离概形之间的拟有限分离态射，$\Lambda$ 是挠系数环。
引理 \ref{more-etale-lemma-upper-shriek-derived} 的函子
$Rf^! : D(Y_\etale, \Lambda) \to D(X_\etale, \Lambda)$ 与引理
\ref{more-etale-lemma-lqf-shriek-derived} 的函子 $Rf^!$ 相同。
\end{lemma}

\begin{proof}
由引理 \ref{more-etale-lemma-Rf-shriek-for-quasi-finite}，有 $Rf_! = f_!$；
结论遂由伴随的唯一性得出。
\end{proof}

\begin{lemma}
\label{more-etale-lemma-Rf-upper-shriek-for-etale}
设 $j : U \to X$ 是拟紧拟分离概形之间的分离 \'etale 态射，$\Lambda$ 是挠系数环。
函子 $Rj^! : D(X_\etale, \Lambda) \to D(U_\etale, \Lambda)$ 等于
$j^{-1}$。
\end{lemma}

\begin{proof}
这是因为 $Rj^!$ 与 $j^{-1}$ 都是 $Rj_! = j_!$ 的右伴随。例如见引理
\ref{more-etale-lemma-Rf-upper-shriek-for-quasi-finite} 与
\ref{more-etale-lemma-etale-upper-shriek}。
\end{proof}

\begin{lemma}
\label{more-etale-lemma-twisted-inverse-image-bounded-below}
设 $f : X \to Y$ 是拟紧拟分离概形之间的有限型分离态射，$\Lambda$ 是挠环。
函子 $Rf^!$ 把 $D^+(Y_\etale, \Lambda)$ 送入
$D^+(X_\etale, \Lambda)$。更准确地，存在整数 $N \geq 0$，使得若
$K \in D(Y_\etale, \Lambda)$ 在 $i < a$ 时满足 $H^i(K) = 0$，则在
$i < a - N$ 时 $H^i(Rf^!K) = 0$。
\end{lemma}

\begin{proof}
令 $N$ 为引理 \ref{more-etale-lemma-derived-lower-shriek-bounded} 中得到的整数。
由构造，对 $K \in D(Y_\etale, \Lambda)$ 与
$L \in \in D(X_\etale, \Lambda)$，有
$\Hom_X(L, Rf^!K) = \Hom_Y(Rf_!L, K)$。设当 $i < a$ 时
$H^i(K) = 0$。取 $L = \tau_{\leq a - N - 1}Rf^!K$。由引理
\ref{more-etale-lemma-derived-lower-shriek-bounded}，复形 $Rf_!L$ 的上同调层在次数
$\leq a - 1$ 消没。故由“导出范畴”引理
\ref{derived-lemma-negative-exts}，有 $\Hom_Y(Rf_!L, K) = 0$。
因此典范映射 $\tau_{\leq a - N - 1}Rf^!K \to Rf^!K$ 为零，
这蕴含当 $i \leq a - N - 1$ 时 $H^i(Rf^!K) = 0$。
\end{proof}

\noindent
设 $f : X \to Y$ 是拟分离拟紧概形之间的有限型分离态射，$\Lambda$ 是挠系数环。
对每个 $K \in D(Y_\etale, \Lambda)$ 与 $L \in D(X_\etale, \Lambda)$，
得到典范映射
\begin{equation}
\label{more-etale-equation-sheafy-trace}
Rf_*R\SheafHom_\Lambda(L, Rf^!K) \longrightarrow R\SheafHom_\Lambda(Rf_!L, K)
\end{equation}
具体地，该映射构造为复合
$$
Rf_*R\SheafHom_\Lambda(L, Rf^!K) \to
R\SheafHom_\Lambda(Rf_!L, Rf_!Rf^!K) \to
R\SheafHom_\Lambda(Rf_!L, K)
$$
，其中第一个箭头来自评注 \ref{more-etale-remark-projection-formula-for-internal-hom}，
第二个箭头是该伴随的余单位 $Rf_!Rf^!K \to K$。

\begin{lemma}
\label{more-etale-lemma-iso-on-RSheafHom}
设 $f : X \to Y$ 是拟紧拟分离概形之间的有限型分离态射，$\Lambda$ 是挠环。
对每个 $K \in D(Y_\etale, \Lambda)$ 与 $L \in D(X_\etale, \Lambda)$，
映射 (\ref{more-etale-equation-sheafy-trace})
$$
Rf_*R\SheafHom_\Lambda(L, Rf^!K) \longrightarrow R\SheafHom_\Lambda(Rf_!L, K)
$$
是同构。
\end{lemma}

\begin{proof}
为证明该引理，必须说明：对任意 $M \in D(Y_\etale, \Lambda)$，映射
(\ref{more-etale-equation-sheafy-trace}) 诱导双射
$$
\Hom_Y(M, Rf_*R\SheafHom_\Lambda(L, Rf^!K))
\longrightarrow
\Hom_Y(M, R\SheafHom_\Lambda(Rf_!L, K))
$$
为此使用如下等式链：
\begin{align*}
\Hom_Y(M, Rf_*R\SheafHom_\Lambda(L, Rf^!K))
& =
\Hom_X(f^{-1}M, R\SheafHom_\Lambda(L, Rf^!K)) \\
& =
\Hom_X(f^{-1}M \otimes_\Lambda^\mathbf{L} L, Rf^!K) \\
& =
\Hom_Y(Rf_!(f^{-1}M \otimes_\Lambda^\mathbf{L} L), K) \\
& =
\Hom_Y(M \otimes_\Lambda^\mathbf{L} Rf_!L, K) \\
& =
\Hom_Y(M, R\SheafHom_\Lambda(Rf_!L, K))
\end{align*}
第一个等号由“位点上同调”引理 \ref{sites-cohomology-lemma-adjoint} 成立。
第二个等号由“位点上同调”引理
\ref{sites-cohomology-lemma-internal-hom} 成立。第三个等号由 $Rf^!$ 的构造成立。
第四个等号由引理 \ref{more-etale-lemma-projection-formula-Rf-lower-shriek} 成立
（这是关键一步）。第五个等号由“位点上同调”引理
\ref{sites-cohomology-lemma-internal-hom} 成立。
\end{proof}

\begin{lemma}
\label{more-etale-lemma-iso-global-hom}
设 $f : X \to Y$ 是拟分离拟紧概形之间的有限型分离态射，$\Lambda$ 是挠环。
对每个 $K \in D(Y_\etale, \Lambda)$ 与 $L \in D(X_\etale, \Lambda)$，
映射 (\ref{more-etale-equation-sheafy-trace}) 诱导同构
$$
R\Hom_X(L, Rf^!K) \longrightarrow R\Hom_Y(Rf_!L, K)
$$
，其中两端为全局导出 Hom。
\end{lemma}

\begin{proof}
由“位点上同调”第 \ref{sites-cohomology-section-global-RHom} 节中的构造，有
$$
R\Hom_X(L, Rf^!K) =
R\Gamma(X, R\SheafHom_\Lambda(L, Rf^!K)) =
R\Gamma(Y, Rf_*R\SheafHom_\Lambda(L, Rf^!K))
$$
（第二个等号由 Leray 成立），以及
$$
R\Hom_Y(Rf_!L, K) = R\Gamma(Y, R\SheafHom_\Lambda(Rf_!L, K))
$$
因此该引理是引理 \ref{more-etale-lemma-iso-on-RSheafHom} 的推论。
\end{proof}

\begin{lemma}
\label{more-etale-lemma-cartesian-square-Rf-upper-shriek}
考虑笛卡尔方块
$$
\xymatrix{
X' \ar[r]_{g'} \ar[d]_{f'} & X \ar[d]^f \\
Y' \ar[r]^g & Y
}
$$
，其中各概形拟紧且拟分离，而 $f$ 分离且有限型。则有
$Rf^! \circ Rg_* = Rg'_* \circ R(f')^!$。
\end{lemma}

\begin{proof}
由伴随函子的唯一性，这从导出下叹号的基变换得出：由引理
\ref{more-etale-lemma-base-change-shriek}，有
$g^{-1} \circ Rf_! = Rf'_! \circ (g')^{-1}$。
\end{proof}

\begin{remark}
\label{more-etale-remark-Rf-upper-shriek-change-of-rings}
设 $\Lambda_1 \to \Lambda_2$ 是挠环同态，$f : X \to Y$ 是拟紧拟分离概形之间的
有限型分离态射。图
$$
\xymatrix{
D(X_\etale, \Lambda_2) \ar[r]_{res} &
D(X_\etale, \Lambda_1) \\
D(Y_\etale, \Lambda_2) \ar[r]^{res} \ar[u]^{Rf^!} &
D(Y_\etale, \Lambda_1) \ar[u]_{Rf^!}
}
$$
交换；其中 $res$ 是“限制”函子，它通过给定环映射把 $\Lambda_2$-模视为
$\Lambda_1$-模。这由伴随的唯一性、评注
\ref{more-etale-remark-Rf-lower-shriek-change-of-rings} 中第二个交换图，以及等式
$$
\Hom_{\Lambda_2}(K_1 \otimes_{\Lambda_1}^\mathbf{L} \Lambda_2, K_2) =
\Hom_{\Lambda_1}(K_1, res(K_2))
$$
成立这一事实得出。无论对象位于 $X_\etale$ 上还是 $Y_\etale$ 上，该等式都是
“位点上同调”引理 \ref{sites-cohomology-lemma-adjoint} 的一个非常特殊的情形。
\end{remark}











\section{紧支撑上同调}
\label{more-etale-section-compactly-supported-cohomology}

\noindent
设 $k$ 是域，$\Lambda$ 是环，$X$ 是 $k$ 上的有限型分离概形，其结构态射为
$f : X \to \Spec(k)$。在第
\ref{more-etale-section-derived-lower-shriek-compactification} 节中，我们定义了函子
$Rf_! : D^+_{tors}(X_\etale, \Lambda) \to D^+_{tors}(\Spec(k), \Lambda)$
；若 $\Lambda$ 是挠环，还定义了函子
$Rf_! : D(X_\etale, \Lambda) \to D(\Spec(k), \Lambda)$。
再与 $\Spec(k)$ 上的全局截面函子复合，便得到我们所谓的紧支撑上同调。

\begin{definition}
\label{more-etale-definition-cohomology-compact-support}
设 $X$ 是域 $k$ 上的有限型分离概形，$\Lambda$ 是环。设 $K$ 是
$D^+_{tors}(X_\etale, \Lambda)$ 的对象；若 $\Lambda$ 为挠环，也可设它是
$D(X_\etale, \Lambda)$ 的对象。$K$ 的 {\it 带紧支撑上同调}，亦称 $K$ 的
{\it 紧支撑上同调}，是
$$
R\Gamma_c(X, K) = R\Gamma(\Spec(k), Rf_!K)
$$
，其中 $f : X \to \Spec(k)$ 是结构态射。我们记
$H^i_c(X, K) = H^i(R\Gamma_c(X, K))$。
\end{definition}

\noindent
该定义与定义 \ref{more-etale-definition-compact-support} 不冲突；这将在引理
\ref{more-etale-lemma-compact-support-h0} 中检验。该定义的效用在于如下结果。

\begin{lemma}
\label{more-etale-lemma-stalk-R-f-shriek}
设 $f : X \to Y$ 是概形的有限型分离态射，其中 $Y$ 拟紧且拟分离。
设 $K$ 是 $D^+_{tors}(X_\etale, \Lambda)$ 的对象；若 $\Lambda$ 为挠环，
也可设它是 $D(X_\etale, \Lambda)$ 的对象。则对任意几何点
$\overline{y} : \Spec(k) \to Y$，在 $D(\Lambda)$ 中有典范同构
$$
(Rf_!K)_{\overline{y}}
\longrightarrow
R\Gamma_c(X_{\overline{y}}, K|_{X_{\overline{y}}})
$$
\end{lemma}

\begin{proof}
这是引理 \ref{more-etale-lemma-base-change-shriek} 与各定义的直接推论。
\end{proof}

\begin{lemma}
\label{more-etale-lemma-compact-support-h0}
设 $X$ 是域 $k$ 上的有限型分离概形。若 $\mathcal{F}$ 是挠阿贝尔层，
则定义 \ref{more-etale-definition-compact-support} 中定义的阿贝尔群
$H^0_c(X, \mathcal{F})$ 与定义
\ref{more-etale-definition-cohomology-compact-support} 中定义的阿贝尔群
$H^0_c(X, \mathcal{F})$ 一致。
\end{lemma}

\begin{proof}
选取在 $k$ 上的紧化 $j : X \to \overline{X}$。两种情形中，该群都被定义为
$H^0(\overline{X}, j_!\mathcal{F})$。对第一个版本，这是引理
\ref{more-etale-lemma-compactify-compact-support} 的结论；对第二个版本，则由构造成立。
\end{proof}

\begin{lemma}
\label{more-etale-lemma-finiteness-curves}
设 $k$ 是代数闭域，$X$ 是 $k$ 上维数 $\leq 1$ 的有限型有限型分离概形，
$\Lambda$ 是 Noether 环。设 $\mathcal{F}$ 是 $X$ 上的可构造挠
$\Lambda$-模层。则 $H^q_c(X, \mathcal{F})$ 是有限 $\Lambda$-模。
\end{lemma}

\begin{proof}
这是“\'Etale 上同调”定理
\ref{etale-cohomology-theorem-vanishing-affine-curves-coefficients} 的推论。
具体地，选取紧化 $j : X \to \overline{X}$。以 $X$ 的概形论闭包替换
$\overline{X}$ 后，可设 $\dim(\overline{X}) \leq 1$。于是
$H^q_c(X, \mathcal{F}) = H^q(\overline{X}, j_!\mathcal{F})$，可应用该定理。
\end{proof}

\begin{remark}[紧支撑上同调的协变性]
\label{more-etale-remark-covariance-compactly-supported}
设 $k$ 是域，$f : X \to Y$ 是 $k$ 上有限型分离概形之间的态射。若
$X$、$Y$ 与 $f$ 满足下列条件之一：
\begin{enumerate}
\item $f$ 为 \'etale；或
\item $f$ 平坦且拟有限；或
\item $f$ 拟有限且 $Y$ 几何单枝；或
\item $f$ 拟有限，且存在 $f$ 的加权 $w : X \to \mathbf{Z}$，
\end{enumerate}
则紧支撑上同调关于 $f$ 协变。更准确地，设 $\Lambda$ 是环，$K$ 是
$D^+_{tors}(Y_\etale, \Lambda)$ 的对象；若 $\Lambda$ 为挠环，也可设它是
$D(Y_\etale, \Lambda)$ 的对象。在假设 (1)--(4) 之一成立时，有典范映射
$$
\text{Tr}_{f, w, K} : f_!f^{-1}K \longrightarrow K
$$
迹映射的存在性见第 \ref{more-etale-section-weightings} 节；情形 (2) 与 (3) 分别见例
\ref{more-etale-example-trace-for-flat-quasi-finite} 和
\ref{more-etale-example-trace-for-quasi-finite-over-normal}。若
$p : X \to \Spec(k)$ 与 $q : Y \to \Spec(k)$ 表示结构态射，则由引理
\ref{more-etale-lemma-shriek-composition} 以及引理
\ref{more-etale-lemma-Rf-shriek-for-quasi-finite} 对拟有限分离态射 $f$ 给出的
$Rf_! = f_!$，有 $Rq_! \circ f_! = Rp_!$。因此可考虑映射
\begin{align*}
R\Gamma_c(X, f^{-1}K)
& =
R\Gamma(\Spec(k), Rp_!f^{-1}K) \\
& =
R\Gamma(\Spec(k), Rq_!f_!f^{-1}K) \\
& \xrightarrow{Rq_!\text{Tr}_{f, w, K}}
R\Gamma(\Spec(k), Rq_!K) \\
& =
R\Gamma_c(Y, K)
\end{align*}
特别地，若 $\Lambda$ 是挠环，则得到箭头
$$
\text{Tr}_f : R\Gamma_c(X, \Lambda) \longrightarrow R\Gamma_c(Y, \Lambda)
$$
该映射还有许多附加性质；例如，它与基域扩张相容。
\end{remark}










\section{一个可构造性结果}
\label{more-etale-section-family-smooth-curves-cohomology}

\noindent
我们“计算”常系数光滑射影曲线族的上同调。

\begin{lemma}
\label{more-etale-lemma-frobenius-linear-on-vb}
设 $p$ 是素数，$S$ 是 $\mathbf{F}_p$ 上的概形。设 $\mathcal{E}$ 是有限局部自由
$\mathcal{O}_S$-模，并把它视为 $S_\etale$ 上的 $\mathcal{O}_S$-模。
设 $F : \mathcal{E} \to \mathcal{E}$ 是 $S_\etale$ 上阿贝尔层的同态，
且对 $S_\etale$ 上 $\mathcal{O}_S$、$\mathcal{E}$ 的局部截面 $a$、$e$，满足
$F(a e) = a^pF(e)$。则
$$
\Coker(F - 1 : \mathcal{E} \to \mathcal{E})
$$
为零，且
$$
\Ker(F - 1 : \mathcal{E} \to \mathcal{E})
$$
是 $S_\etale$ 上的可构造阿贝尔层。
\end{lemma}

\noindent
该引理是“\'Etale 上同调”引理
\ref{etale-cohomology-lemma-F-1} 的推广。

\begin{proof}
可设 $S = \Spec(A)$，其中 $A$ 是 $\mathbf{F}_p$-代数，并设 $\mathcal{E}$ 是
自由 $A$-模 $Ae_1 \oplus \ldots \oplus Ae_n$ 所伴的拟凝聚模。写成
$F(e_i) = \sum a_{ij} e_j$。

\medskip\noindent
$F - 1$ 的满射性。只需证明：以一个忠实平坦 \'etale 扩张替换 $A$ 后，
任意元素 $\sum a_i e_i$（$a_i \in A$）都属于 $F - 1$ 的像。注意
$$
F(\sum x_ie_i) - \sum x_i e_i = \sum x_i^p a_{ij} e_j - \sum x_i e_i
$$
考虑 $A$-代数
$$
A' = A[x_1, \ldots, x_n]/(a_i + x_i - \sum\nolimits_j a_{ji} x_j^p)
$$
计算表明，$\text{d}x_i$ 在 $\Omega_{A'/A}$ 中为零，因而
$\Omega_{A'/A} = 0$。由于 $A'$ 是 $A$ 上有限型的，这蕴含
$\Spec(A') \to \Spec(A)$ 非分歧，因而拟有限。又因 $A'$ 由 $n$ 个元素生成，
且由 $n$ 个方程截出，故 $A'$ 是 $A$ 上的全局相对完全交。因此 $A'$ 在 $A$
上平坦，遂知 $A \to A'$ 是 \'etale 的（作为平坦且非分歧的环映射）。最后，
读者可通过直接验证 $\Spec(A') \to \Spec(A)$ 的所有几何纤维非空来证明
$A \to A'$ 忠实平坦；不过这也由“\'Etale 上同调”引理
\ref{etale-cohomology-lemma-F-1} 得出。最后，元素
$\sum x_i e_i \in A'e_1 \oplus \ldots \oplus A'e_n$ 在 $F - 1$ 下映到
$\sum a_i e_i$。

\medskip\noindent
核的可构造性。上述计算表明，$\Ker(F - 1)$ 由概形
$$
\Spec(A[x_1, \ldots, x_n]/(x_i - \sum\nolimits_j a_{ji} x_j^p))
$$
在 $S = \Spec(A)$ 上表示。由于这是 $S$ 上仿射且 \'etale 的概形，结论由
“\'Etale 上同调”引理 \ref{etale-cohomology-lemma-jshriek-constructible} 得出。
\end{proof}

\begin{lemma}
\label{more-etale-lemma-proper-smooth-family-curves}
设 $f : X \to S$ 是概形的固有光滑态射，其纤维几何连通且维数为 $1$。
设 $\ell$ 是素数。则 $R^qf_*\underline{\mathbf{Z}/\ell\mathbf{Z}}$
是一个可构造对象。
\end{lemma}

\begin{proof}
可设 $S$ 仿射，写成 $S = \Spec(A)$。若把 $A = \bigcup A_i$ 写成其有限型
$\mathbf{Z}$-子代数之并，则可找到某个 $i$ 以及有限型态射
$f_i : X_i \to S_i = \Spec(A_i)$，使其到 $S$ 的基变换为
$f : X \to S$；见“极限”引理 \ref{limits-lemma-descend-finite-presentation}。
增大 $i$ 后，可设 $f_i : X_i \to S_i$ 光滑、固有且相对维数为 $1$；
见“极限”引理
\ref{limits-lemma-eventually-proper}
\ref{limits-lemma-descend-smooth} 与
\ref{limits-lemma-descend-dimension-d}。由“态射的进一步结果”引理
\ref{more-morphisms-lemma-proper-flat-geom-red}，得到开子概形
$U_i \subset S_i$，使 $f_i : X_i \to S_i$ 在 $U_i$ 上的纤维几何连通。
于是 $S \to S_i$ 映入 $U_i$。可以用
$f_i : f_i^{-1}(U_i) \to U_i$ 替换 $X \to S$，从而化归到下一段讨论的情形。

\medskip\noindent
设 $S$ 为 Noether 概形。可写成 $S = U \cup Z$，其中 $U$ 是由 $\ell$
不消没所定义的开子概形，而 $Z = V(\ell) \subset S$。由于
$R^qf_*\underline{\mathbf{Z}/\ell\mathbf{Z}}$ 的构造与任意基变换交换
（“\'Etale 上同调”定理 \ref{etale-cohomology-theorem-proper-base-change}），
只需分别在 $U$ 与 $Z$ 上证明。于是化归到两个情形：
(a) $\ell$ 在 $S$ 上可逆；(b) $\ell$ 在 $S$ 上为零。

\medskip\noindent
情形 (a)。我们断言，此时层 $R^qf_*\underline{\mathbf{Z}/\ell\mathbf{Z}}$
在 $S$ 上有限且局部常值。首先，由固有基变换（采用“\'Etale 上同调”引理
\ref{etale-cohomology-lemma-proper-base-change-stalk} 的形式）以及有限性
（“\'Etale 上同调”定理
\ref{etale-cohomology-theorem-vanishing-affine-curves}），可知
$R^qf_*\underline{\mathbf{Z}/\ell\mathbf{Z}}$ 的茎有限。由“\'Etale 上同调”引理
\ref{etale-cohomology-lemma-sp-isom-proper-loc-cst-torsion}，
所有特殊化映射均为同构。故由“\'Etale 上同调”引理
\ref{etale-cohomology-lemma-characterize-locally-constant-module}，断言成立。

\medskip\noindent
情形 (b)。这里 $\ell = p$ 是素数，且 $S$ 是 $\Spec(\mathbf{F}_p)$ 上的概形。
由上述相同文献，已知 $R^qf_*\underline{\mathbf{Z}/p\mathbf{Z}}$ 的茎有限，
且在 $q \geq 2$ 时为零。由“\'Etale 上同调”引理
\ref{etale-cohomology-lemma-sections-upstairs}，有
$f_*\underline{\mathbf{Z}/p\mathbf{Z}} =
\underline{\mathbf{Z}/p\mathbf{Z}}$。只剩证明
$R^1f_*\underline{\mathbf{Z}/p\mathbf{Z}}$ 可构造。考虑 Artin--Schreier 序列
$$
0 \to \underline{\mathbf{Z}/p\mathbf{Z}} \to \mathcal{O}_X
\xrightarrow{F - 1} \mathcal{O}_X
\to 0
$$
见“\'Etale 上同调”第 \ref{etale-cohomology-section-artin-schreier} 节。
回忆，$f_*\mathcal{O}_X = \mathcal{O}_S$，且
$R^1f_*\mathcal{O}_X$ 是有限局部自由 $\mathcal{O}_S$-模，其秩等于
$X \to S$ 各纤维的亏格；见“代数曲线”引理
\ref{curves-lemma-genus-in-nodal-family-of-curves}。故有短正合列
$$
0 \to
\Coker(F - 1 : \mathcal{O}_S \to \mathcal{O}_S)
\to R^1f_*\underline{\mathbf{Z}/p\mathbf{Z}} \to
\Ker(F - 1 : R^1f_*\mathcal{O}_X \to R^1f_*\mathcal{O}_X)
\to 0
$$
应用引理 \ref{more-etale-lemma-frobenius-linear-on-vb} 即得结论。
\end{proof}

\begin{lemma}
\label{more-etale-lemma-proper-smooth-family-curves-modules}
设 $f : X \to S$ 是概形的固有光滑态射，其纤维几何连通且维数为 $1$。
设 $\Lambda$ 是 Noether 环，$M$ 是被整数 $n > 0$ 消去的有限 $\Lambda$-模。
则 $R^qf_*\underline{M}$ 是 $S$ 上的可构造 $\Lambda$-模层。
\end{lemma}

\begin{proof}
若对某个素数 $\ell$ 有 $n = \ell n'$，则得到有限 $\Lambda$-模的短正合列
$0 \to M[\ell] \to M \to M' \to 0$，且 $M'$ 被 $n'$ 消去。
这产生相应的常值层短正合列，继而产生正合列
$$
R^{q - 1}f_*\underline{M'} \to
R^qf_*\underline{M[n]} \to
R^qf_*\underline{M} \to
R^qf_*\underline{M'} \to
R^{q + 1}f_*\underline{M[n]}
$$
因此，若能在 $M$ 被某个素数消去时证明该结果，则对 $n$ 作归纳，并应用
“\'Etale 上同调”引理 \ref{etale-cohomology-lemma-constructible-abelian}，
即得一般情形。

\medskip\noindent
设 $\ell$ 是素数，且 $\ell$ 消去 $M$。于是可用 $\mathbf{F}_\ell$-代数
$\Lambda/\ell \Lambda$ 替换 $\Lambda$。具体地，把 $\underline{M}$ 视为
$\Lambda$-模层时所得的层 $R^qf_*\underline{M}$，与把 $\underline{M}$ 视为
$\Lambda/\ell \Lambda$-模层时算得的层 $R^qf_*\underline{M}$ 相同；
见“位点上同调”引理 \ref{sites-cohomology-lemma-modules-abelian-unbounded}。

\medskip\noindent
设 $\ell$ 是素数，且 $\ell$ 消去 $M$ 与 $\Lambda$。我们化归到 $M$ 为有限自由
$\Lambda$-模的情形。具体地，选取消解
$$
\ldots \to
\Lambda^{\oplus m_2} \to
\Lambda^{\oplus m_1} \to
\Lambda^{\oplus m_0} \to
M \to 0
$$
回忆，$f_*$ 在 $\Lambda$-模层上的上同调维数有限；见“\'Etale 上同调”引理
\ref{etale-cohomology-lemma-cohomological-dimension-proper} 以及“导出范畴”引理
\ref{derived-lemma-unbounded-right-derived}。因此，对某个充分大的整数 $a$，
$R^qf_*\underline{M}$ 是对象
$$
Rf_*(\underline{\Lambda^{\oplus m_a}} \to \ldots \to
\underline{\Lambda^{\oplus m_0}})
$$
在 $D(S_\etale, \Lambda)$ 中的第 $q$ 个上同调层。使用“导出范畴”引理
\ref{derived-lemma-two-ss-complex-functor} 的第一个谱序列（或者使用截断论证），
可知只需证明 $R^qf_*\underline{\Lambda})$ 可构造。

\medskip\noindent
此时终于可以使用等式
$$
(Rf_*\underline{\mathbf{Z}/\ell\mathbf{Z}})
\otimes_{\mathbf{Z}/\ell\mathbf{Z}}^\mathbf{L} \underline{\Lambda} =
Rf_*\underline{\Lambda}
$$
；它由“\'Etale 上同调”引理
\ref{etale-cohomology-lemma-projection-formula-proper-mod-n} 成立。
由于域 $\mathbf{Z}/\ell\mathbf{Z}$ 上任意模都平坦，得到
$$
(R^qf_*\underline{\mathbf{Z}/\ell\mathbf{Z}})
\otimes_{\mathbf{Z}/\ell\mathbf{Z}} \underline{\Lambda} =
R^qf_*\underline{\Lambda}
$$
故由“\'Etale 上同调”引理
\ref{etale-cohomology-lemma-tensor-constructible}，只需对
$R^qf_*\underline{\mathbf{Z}/\ell\mathbf{Z}}$ 证明结论。这正是引理
\ref{more-etale-lemma-proper-smooth-family-curves}。
\end{proof}






\section{上同调可构造的复形}
\label{more-etale-section-Dc}

\noindent
我们继续“\'Etale 上同调”第 \ref{etale-cohomology-section-Dc} 节开始的讨论。
特别地，对概形 $X$ 与 Noether 环 $\Lambda$，以 $D_c(X_\etale, \Lambda)$
表示 $D(X_\etale, \Lambda)$ 的严格全且饱和的三角子范畴；其中对象的
上同调层均为可构造 $\Lambda$-模层。

\begin{lemma}
\label{more-etale-lemma-qf-f-shriek-constructible}
设 $f : X \to Y$ 是局部拟有限且有限表示的概形态射。引理
\ref{more-etale-lemma-lqf-shriek-derived} 的函子
$f_! : D(X_\etale, \Lambda) \to D(Y_\etale, \Lambda)$ 把
$D_c(X_\etale, \Lambda)$ 送入 $D_c(Y_\etale, \Lambda)$。
\end{lemma}

\begin{proof}
由于函子 $f_!$ 正合，只需证明：对 $X_\etale$ 上任意可构造
$\Lambda$-模层 $\mathcal{F}$，$f_!\mathcal{F}$ 可构造。该问题在 $Y$ 上局部，
故可且确实设 $Y$ 仿射。于是 $X$ 拟紧且拟分离；见“态射”定义
\ref{morphisms-definition-finite-presentation}。设
$X = \bigcup_{i = 1, \ldots, n} X_i$ 是有限仿射开覆盖。由引理
\ref{more-etale-lemma-lqf-colimit-f-shriek}，只需证明
$f_{i, !}\mathcal{F}|_{X_i}$ 与 $f_{ii', !}\mathcal{F}|_{X_i \cap X_{i'}}$
可构造；其中 $f_i : X_i \to Y$ 与
$f_{ii'} : X_i \cap X_{i'} \to Y$ 是 $f$ 的限制。由于 $X_i$ 和
$X_i \cap X_{i'}$ 拟紧且分离，这意味着可设 $f$ 分离。由 Zariski 主定理
（采用“态射的进一步结果”引理
\ref{more-morphisms-lemma-quasi-finite-separated-pass-through-finite-addendum}
的形式），可选取分解 $f = g \circ j$，其中 $j : X \to X'$ 是开浸入，
$g : X' \to Y$ 有限且有限表示。由引理
\ref{more-etale-lemma-f-shriek-composition}，有 $f_! = g_! \circ j_!$。
由“\'Etale 上同调”引理 \ref{etale-cohomology-lemma-jshriek-constructible}，
$j_!\mathcal{F}$ 在 $X'$ 上可构造。态射 $g$ 有限，故由引理
\ref{more-etale-lemma-proper-f-shriek}，有 $g_! = g_*$。因此由“\'Etale 上同调”引理
\ref{etale-cohomology-lemma-finite-pushforward-constructible}，
$f_!\mathcal{F} = g_!j_!\mathcal{F} = g_*j_!\mathcal{F}$ 可构造。
\end{proof}

\begin{lemma}
\label{more-etale-lemma-constant-shriek-rel-dim-1}
设 $S$ 是有限维 Noether 仿射概形，$f : X \to S$ 是相对维数为 $1$ 的
分离、仿射、光滑态射。设 $\Lambda$ 是挠 Noether 环，$M$ 是有限
$\Lambda$-模。则 $Rf_!\underline{M}$ 的上同调层可构造。
\end{lemma}

\begin{proof}
我们对 $d = \dim(S)$ 作归纳来证明该结果。

\medskip\noindent
起始情形。若 $d = 0$，则只需证明 $R^qf_!\underline{M}$ 的茎是有限
$\Lambda$-模。若 $\overline{s}$ 是 $S$ 的几何点，则由引理
\ref{more-etale-lemma-stalk-R-f-shriek}，有
$(R^qf_!\underline{M})_{\overline{s}} = H^q_c(X_{\overline{s}}, \underline{M})$。
由引理 \ref{more-etale-lemma-finiteness-curves}，这是有限 $\Lambda$-模。

\medskip\noindent
归纳步骤。只需找到稠密开集 $U \subset S$，使
$Rf_!\underline{M}|_U$ 的上同调层可构造。事实上，由归纳假设以及
$Rf_!\underline{M}$ 的构造与所有基变换交换这一事实（引理
\ref{more-etale-lemma-base-change-shriek}），$Rf_!\underline{M}$ 在补集
$S \setminus U$ 上的限制将有可构造上同调层。具体地，设
$\eta \in S$ 是 $S$ 某个不可约分支的泛点。只需找到 $\eta$ 的开邻域 $U$，
使 $Rf_!\underline{M}$ 在 $U$ 上的限制可构造。下一段将完成此事。

\medskip\noindent
给定泛点 $\eta \in S$，选取图
$$
\xymatrix{
\overline{Y}_1 \amalg \ldots \amalg \overline{Y}_n \ar[rd] &
Y_1 \amalg \ldots \amalg Y_n \ar[r]_-\nu \ar[d] \ar[l]^j &
X_V \ar[r] \ar[d] &
X_U \ar[r] \ar[d] &
X \ar[d]^f \\
& T_1 \amalg \ldots \amalg T_n \ar[r] &
V \ar[r] &
U \ar[r] &
S
}
$$
，如“态射的进一步结果”引理
\ref{more-morphisms-lemma-make-good-curves}。我们将证明
$Rf_!\underline{M}|_U$ 可构造。首先，由于 $V \to U$ 有限且满，只需证明
到 $V$ 的拉回可构造；见“\'Etale 上同调”引理
\ref{etale-cohomology-lemma-check-constructible}。由于 $Rf_!$ 的构造与基变换
交换，只需证明 $R(X_V \to V)_!\underline{M}$ 可构造。令 $W \subset X_V$
为“态射的进一步结果”引理
\ref{more-morphisms-lemma-make-good-curves} (4) 给出的开子概形。令
$Z \subset X_V$ 为 $W$ 在 $X_V$ 中的补集所带的既约诱导概形结构。
则 $Z \to V$ 的纤维维数为 $0$（因为 $W$ 在各纤维中稠密），故
$Z \to V$ 拟有限。由特异三角
$$
R(W \to V)_!\underline{M} \to
R(X_V \to V)_!\underline{M} \to
R(Z \to V)_!\underline{M} \to
\ldots
$$
以及引理 \ref{more-etale-lemma-relative-triangle-associated-to-open} 和
\ref{more-etale-lemma-qf-f-shriek-constructible}，可知只需证明
$R(W \to V)_!\underline{M}$ 的上同调层可构造。其次，有
$$
R(W \to V)_!\underline{M} = R(\nu^{-1}(W) \to V)_!\underline{M}
$$
，因为态射 $\nu : \nu^{-1}(W) \to W$ 是增厚，且可应用引理
\ref{more-etale-lemma-shriek-and-thickening}。接着，以
$Z' \subset \coprod \overline{Y}_i$ 表示开集 $j(\nu^{-1}(W))$ 的补集。
$Z' \to V$ 仍拟有限。再次使用特异三角
$$
R(\nu^{-1}(W) \to V)_!\underline{M} \to
R(\coprod \overline{Y}_i \to V)_!\underline{M} \to
R(Z' \to V)_!\underline{M} \to
\ldots
$$
，可知只需证明
$$
R(\coprod \overline{Y}_i \to V)_!\underline{M} =
\bigoplus\nolimits_i R(\overline{Y}_i \to V)_!\underline{M} =
\bigoplus\nolimits_i R(T_i \to V)_!R(\overline{Y}_i \to T_i)_!\underline{M}
$$
的上同调层可构造（第二个等号由引理
\ref{more-etale-lemma-shriek-composition} 成立）。对
$R(\overline{Y}_i \to T_i)_!\underline{M}$，结论是引理
\ref{more-etale-lemma-proper-smooth-family-curves-modules}。又因 $T_i \to V$ 有限 \'etale，
可应用引理 \ref{more-etale-lemma-qf-f-shriek-constructible}，故证明完成。
\end{proof}

\begin{lemma}
\label{more-etale-lemma-loc-constant-shriek-rel-dim-1}
设 $Y$ 是有限维 Noether 仿射概形，$\Lambda$ 是挠 Noether 环。
设 $\mathcal{F}$ 是开子概形 $U \subset \mathbf{A}^1_Y$ 上有限型局部常值的
$\Lambda$-模层。则 $Rf_!\mathcal{F}$ 的上同调层可构造，其中
$f : U \to Y$ 是结构态射。
\end{lemma}

\begin{proof}
可把 $\Lambda$ 分解为乘积
$\Lambda = \Lambda_1 \times \ldots \times \Lambda_r$
，其中对某个素数 $\ell_i$，$\Lambda_i$ 是 $\ell_i$-准素的。因此可设存在素数
$\ell$ 与整数 $n > 0$，使 $\ell^n$ 消去 $\Lambda$（因而也消去
$\mathcal{F}$）。

\medskip\noindent
由于 $U$ 为 Noether 概形，可见 $U$ 只有有限多个连通分支。因此可设 $U$ 连通。
令 $g : U' \to U$ 为“\'Etale 上同调”引理
\ref{etale-cohomology-lemma-pullback-filtered-modules} 中构造的有限 \'etale 覆盖。
“\'Etale 上同调”第 \ref{etale-cohomology-section-trace-method} 节的讨论给出映射
$$
\mathcal{F} \to g_*g^{-1}\mathcal{F} \to \mathcal{F}
$$
，其复合为同构。因此只需对 $g_*g^{-1}\mathcal{F}$ 证明结论。另一方面，
由引理 \ref{more-etale-lemma-shriek-composition}，有
$Rf_!g_*g^{-1}\mathcal{F} = R(f \circ g)_!g^{-1}\mathcal{F}$。
由于（根据我们对 $g$ 的选取）$g^{-1}\mathcal{F}$ 有有限滤过，其分次商是形如
$\underline{M}$ 的常值 $\Lambda$-模层，其中 $M$ 是有限 $\Lambda$-模，
这就化归到引理 \ref{more-etale-lemma-constant-shriek-rel-dim-1} 中已证明的情形。
\end{proof}

\begin{lemma}
\label{more-etale-lemma-constructible-shriek-rel-dim-1}
设 $Y$ 是仿射概形，$\Lambda$ 是 Noether 环。设 $\mathcal{F}$ 是
$\mathbf{A}^1_Y$ 上的可构造挠 $\Lambda$-模层。则 $Rf_!\mathcal{F}$ 的
上同调层可构造，其中 $f : \mathbf{A}^1_Y \to Y$ 是结构态射。
\end{lemma}

\begin{proof}
设 $\mathcal{F}$ 被 $n > 0$ 消去。于是可用 $\Lambda/n\Lambda$ 替换
$\Lambda$ 而不改变 $Rf_!\mathcal{F}$。因此可且确实设 $\Lambda$ 为挠环。

\medskip\noindent
写 $Y = \Spec(R)$。若把 $R = \bigcup R_i$ 写成其有限型
$\mathbf{Z}$-子代数之并，则可找到某个 $i$，使 $\mathcal{F}$ 是
$\mathbf{A}^1_{R_i}$ 上某个可构造 $\Lambda$-模层的拉回；见“\'Etale 上同调”引理
\ref{etale-cohomology-lemma-category-is-colimit}。因此可设 $Y$ 是有限维 Noether 概形。

\medskip\noindent
设 $Y$ 是维数 $d = \dim(Y)$ 的有限维 Noether 概形，且 $\Lambda$ 为挠环。
我们对 $d$ 作归纳证明该结果。

\medskip\noindent
起始情形。若 $d = 0$，则只需证明 $R^qf_!\mathcal{F}$ 的茎是有限
$\Lambda$-模。若 $\overline{y}$ 是 $Y$ 的几何点，则由引理
\ref{more-etale-lemma-stalk-R-f-shriek}，有
$(R^qf_!\mathcal{F})_{\overline{y}} = H^q_c(X_{\overline{y}}, \mathcal{F})$。
由引理 \ref{more-etale-lemma-finiteness-curves}，这是有限 $\Lambda$-模。

\medskip\noindent
归纳步骤。只需找到稠密开集 $V \subset Y$，使
$Rf_!\mathcal{F}|_V$ 的上同调层可构造。事实上，由归纳假设以及
$Rf_!\mathcal{F}$ 的构造与所有基变换交换这一事实（引理
\ref{more-etale-lemma-base-change-shriek}），$Rf_!\mathcal{F}$ 在补集
$Y \setminus V$ 上的限制将有可构造上同调层。由可构造 $\Lambda$-模层的定义，
存在稠密开子概形 $U \subset \mathbf{A}^1_Y$，使 $\mathcal{F}|_U$ 是有限型局部常值
$\Lambda$-模层。以 $Z \subset \mathbf{A}^1_Y$ 表示补集（视为既约闭子概形）。
注意，$U$ 包含 $\mathbf{A}^1_Y \to Y$ 在 $Y$ 各不可约分支的泛点
$\xi_1, \ldots, \xi_n$ 上方纤维的所有泛点。因此 $Z \to Y$ 在
$\xi_1, \ldots, \xi_n$ 上方的纤维有限。以一个稠密开集替换 $Y$（这是允许的）后，
可设 $Z \to Y$ 有限；见“态射”引理 \ref{morphisms-lemma-generically-finite}。
由引理 \ref{more-etale-lemma-relative-triangle-associated-to-open} 的特异三角以及
$Z \to Y$ 的结果（引理 \ref{more-etale-lemma-qf-f-shriek-constructible}），化归到证明
$R(U \to Y)_!\mathcal{F}$ 的上同调层可构造。这就是引理
\ref{more-etale-lemma-loc-constant-shriek-rel-dim-1}。
\end{proof}

\begin{theorem}
\label{more-etale-theorem-constructible-shriek}
设 $f : X \to Y$ 是拟紧拟分离概形之间有限表示的分离态射，$\Lambda$ 是
Noether 环。设 $K$ 是 $D^+_{tors, c}(X_\etale, \Lambda)$ 的对象；若
$\Lambda$ 为挠环，也可设它是 $D_c(X_\etale, \Lambda)$ 的对象。则
$Rf_!K$ 的上同调层可构造；即 $Rf_!K$ 属于
$D^+_{tors, c}(Y_\etale, \Lambda)$，或者在 $\Lambda$ 为挠环时属于
$D_c(Y_\etale, \Lambda)$。
\end{theorem}

\begin{proof}
该问题在 $Y$ 上局部，故可且确实设 $Y$ 仿射。由归纳原理和引理
\ref{more-etale-lemma-relative-mayer-vietoris}，化归到 $X$ 也仿射的情形。

\medskip\noindent
设 $X$ 与 $Y$ 仿射。由于 $X$ 有限表示，可选取有限表示的闭浸入
$i : X \to \mathbf{A}^n_Y$。若 $p : \mathbf{A}^n_Y \to Y$ 表示结构态射，
则由引理 \ref{more-etale-lemma-shriek-composition}，有 $Rf_! = Rp_! \circ Ri_!$。
由引理 \ref{more-etale-lemma-qf-f-shriek-constructible}，结论对 $Ri_! = i_!$ 成立。
因此可设 $f$ 是投影态射 $\mathbf{A}^n_Y \to Y$。由于可把 $f$ 视为复合
$$
X = \mathbf{A}^n_Y \to \mathbf{A}^{n - 1}_Y \to
\mathbf{A}^{n - 2}_S \to \ldots \to
\mathbf{A}^1_Y \to Y
$$
，可设 $n = 1$。

\medskip\noindent
设 $Y$ 仿射且 $X = \mathbf{A}^1_Y$。由于 $Rf_!$ 的上同调维数有限
（引理 \ref{more-etale-lemma-derived-lower-shriek-bounded}），可设 $K$ 下有界。
使用“导出范畴”引理 \ref{derived-lemma-two-ss-complex-functor} 的第一个谱序列
（或者使用截断论证），化归到引理
\ref{more-etale-lemma-constructible-shriek-rel-dim-1} 的结果。
\end{proof}








\section{应用}
\label{more-etale-section-applications}

\noindent
本节给出定理 \ref{more-etale-theorem-constructible-shriek} 的一些应用。

\begin{lemma}
\label{more-etale-lemma-finiteness-compactly-supported}
设 $k$ 是代数闭域，$X$ 是 $k$ 上的有限型分离概形，$\Lambda$ 是 Noether 环。
设 $K$ 是 $D^+_{tors, c}(X_\etale, \Lambda)$ 的对象；若 $\Lambda$ 为挠环，
也可设它是 $D_c(X_\etale, \Lambda)$ 的对象。则对所有
$i \in \mathbf{Z}$，$H^i_c(X, K)$ 都是有限 $\Lambda$-模。
\end{lemma}

\begin{proof}
这是定理 \ref{more-etale-theorem-constructible-shriek} 与第
\ref{more-etale-section-compactly-supported-cohomology} 节中紧支撑上同调定义的直接推论。
\end{proof}

\begin{proposition}
\label{more-etale-proposition-loc-cst-torsion}
设 $f : X \to S$ 是概形的光滑固有态射，$\Lambda$ 是 Noether 环。
设 $\mathcal{F}$ 是 $X_\etale$ 上有限型局部常值的 $\Lambda$-模层，并且对
$X$ 的每个几何点 $\overline{x}$，茎 $\mathcal{F}_{\overline{x}}$ 都被某个
与 $\overline{x}$ 的剩余特征互素的整数 $n > 0$ 消去。则对所有
$i \in \mathbf{Z}$，$R^if_*\mathcal{F}$ 都是 $S_\etale$ 上有限型局部常值的
$\Lambda$-模层。
\end{proposition}

\begin{proof}
该问题在 $S$ 上局部，故可设 $S$ 仿射。对 $X$ 的点 $x$，以
$n_x \geq 1$ 表示消去 $\mathcal{F}_{\overline{x}}$ 的最小整数；这里
$\overline{x}$ 是 $X$ 中位于 $x$ 上方的某个（等价地，任意）几何点。
由于 $X$ 拟紧（它在仿射概形上固有），存在有限 \'etale 覆盖
$\{U_j \to X\}_{j = 1, \ldots, m}$，使 $\mathcal{F}|_{U_j}$ 为常值层。
由于 $U_j \to X$ 是开映射，函数 $x \mapsto n_x$ 局部常值且仅取有限多个值。
因此得到如下开闭子概形的有限分解：
$X = X_1 \amalg \ldots \amalg X_N$，使得 $n_x = n$ 当且仅当
$x \in X_n$。于是只需对所诱导的态射 $X_n \to S$ 以及 $\mathcal{F}$ 在
$X_n$ 上的限制证明该引理。因此可且确实设存在整数 $n > 0$，使
$\mathcal{F}$ 被 $n$ 消去，且 $n$ 与 $X$ 所有剩余域的剩余特征互素。

\medskip\noindent
由于 $f$ 光滑且固有，像 $f(X) \subset S$ 开且闭。因此可用 $f(X)$ 替换
$S$ 并设 $f(X) = S$。特别地，可设 $n$ 在定义仿射概形 $S$ 的环中可逆。

\medskip\noindent
本段化归到 $S$ 为 Noether 概形的情形。对某个 $\mathbf{Z}[1/n]$-代数 $A$，
写 $S = \Spec(A)$。把 $A = \bigcup A_i$ 写成其有限型
$\mathbf{Z}[1/n]$-子代数之并。可找到某个 $i$ 以及有限型态射
$f_i : X_i \to S_i = \Spec(A_i)$，使其到 $S$ 的基变换为
$f : X \to S$；见“极限”引理 \ref{limits-lemma-descend-finite-presentation}。
增大 $i$ 后，可设 $f_i : X_i \to S_i$ 光滑且固有；见“极限”引理
\ref{limits-lemma-eventually-proper}
\ref{limits-lemma-descend-smooth} 与
\ref{limits-lemma-descend-dimension-d}.
由“\'Etale 上同调”引理
\ref{etale-cohomology-lemma-category-loc-cst-is-colimit}
，可知存在某个 $i$ 以及有限型局部常值的 $\Lambda$-模层
$\mathcal{F}_i$，其到 $X$ 的拉回与 $\mathcal{F}$ 同构。由于
$\mathcal{F}$ 被 $n$ 消去，可用 $\Ker(n : \mathcal{F}_i \to \mathcal{F}_i)$
替换 $\mathcal{F}_i$，并设同一性质对 $\mathcal{F}_i$ 成立。
这就化归到下一段讨论的情形。

\medskip\noindent
设有整数 $n \geq 1$，基概形 $S$ 为 Noether 概形且位于
$\mathbf{Z}[1/n]$ 上，而 $\mathcal{F}$ 为 $n$-挠层。由定理
\ref{more-etale-theorem-constructible-shriek}，各层 $R^if_*\mathcal{F}$ 都是可构造
$\Lambda$-模层。由“\'Etale 上同调”引理
\ref{etale-cohomology-lemma-sp-isom-proper-torsion-loc-ac}，
$R^if_*\mathcal{F}$ 的特殊化映射总是同构。结论遂由“\'Etale 上同调”引理
\ref{etale-cohomology-lemma-characterize-locally-constant-module} 得出。
\end{proof}







\section{导出上叹号的进一步结果}
\label{more-etale-section-more-derived-upper-shriek}

\noindent
设 $\Lambda$ 是挠环。考虑交换图
$$
\xymatrix{
U \ar[rr]_j \ar[rd]_g & & U' \ar[ld]^{g'} \\
& Y
}
$$
，其中各概形拟紧且拟分离，$g$ 与 $g'$ 分离且有限型，而 $j$ 为 \'etale。
这诱导典范映射
$$
Rg_!\Lambda \longrightarrow Rg'_!\Lambda
$$
成立于 $D(Y_\etale, \Lambda)$ 中。具体地，由引理
\ref{more-etale-lemma-shriek-composition} 与
\ref{more-etale-lemma-Rf-shriek-for-quasi-finite}，有 $Rg_! = Rg'_! \circ j_!$。
另一方面，由于 $j_!$ 是 $j^{-1}$ 的左伴随，有余单位
$\text{Tr}_j : j_!\Lambda = j_!j^{-1}\Lambda \to \Lambda$；
也称它为 $j$ 的迹映射，见评注 \ref{more-etale-remark-trace-etale-counit}。
上述映射构造为复合
$$
Rg_!\Lambda = Rg'_!j_!\Lambda
\xrightarrow{Rg'_! \text{Tr}_j}
Rg'_!\Lambda
$$
再给定一个 \'etale 态射 $j' : U' \to U''$，其中某个
$g'' : U'' \to Y$ 分离且有限型，则映射对 $j$ 与 $j'$ 的复合
$$
Rg_!\Lambda \longrightarrow Rg'_!\Lambda \longrightarrow Rg''_!\Lambda
$$
等于为 $j' \circ j$ 构造的映射
$Rg_!\Lambda \longrightarrow Rg''_!\Lambda$。这由迹映射的相应陈述得出；
更一般的情形见引理 \ref{more-etale-lemma-trace-composition}。

\medskip\noindent
设 $f : X \to Y$ 是拟紧拟分离概形之间的有限型分离态射。于是得到函子
$$
X_{affine, \etale}
\longrightarrow
\left\{
\begin{matrix}
\text{有限型分离概形，底为 }Y\\
\text{其间的态射为 \'etale 态射}
\end{matrix}
\right\}
$$
因此上述构造确定函子
$X_{affine, \etale}^{opp} \to D(Y_\etale, \Lambda)$，
它把 $U$ 送到 $R(U \to Y)_!\Lambda$。

\begin{lemma}
\label{more-etale-lemma-describe-Rf-upper-shriek}
设 $f : X \to Y$ 是拟紧拟分离概形之间的有限型分离态射，$\Lambda$ 是挠环。
设 $K \in D(Y_\etale, \Lambda)$。对 $n \in \mathbf{Z}$，上同调层
$H^n(Rf^!K)$ 在 $X_{affine, \etale}$ 上的限制，是预层
$$
U \longmapsto \Hom_Y(R(U \to Y)_!\Lambda, K[n])
$$
所伴的层。关于 $R(U \to Y)_!\Lambda$ 的函子性质，见上面的讨论。
\end{lemma}

\begin{proof}
设 $j : U \to X$ 是 $X_{affine, \etale}$ 的对象，并令 $g = f \circ j$。
回忆，对 $D(X_\etale, \Lambda)$ 中任意 $M$，有
$\Hom_X(j_!\Lambda, M[n]) = H^n(U, M)$。于是 $H^n(Rf^!K)$ 是预层
$$
U \mapsto H^n(U, Rf^!K) = \Hom_X(j_!\Lambda, Rf^!K[n]) =
\Hom_Y(Rf_!j_!\Lambda, K[n] = \Hom_Y(Rg_!\Lambda, K[n])
$$
所伴的层。略去如下验证：转移映射由我们上面构造的各对象
$Rg_!\Lambda = R(U \to Y)_!\Lambda$ 之间的转移映射给出。
\end{proof}
